% =============================================================================
%  《钱是怎么跑起来的》—— 一本从零开始的金融世界说明书
%  编译： xelatex main -> biber main -> makeindex main -> xelatex main x2
%  或直接： latexmk -xelatex main.tex
% =============================================================================
\documentclass[
  lang=cn,
  scheme=chinese,
  color=blue,
  chinesefont=nofont,
  mode=fancy,
  math=newtx,
  device=normal,
  toc=onecol,
  11pt,
  a4paper,
]{elegantbook}

% =============================================================================
%  《钱是怎么跑起来的》—— 排版样式层
%  在 ElegantBook v4.7 之上，重建网页版的全部内容组件
% =============================================================================

\usepackage{ifthen}
\usepackage{array}
\usepackage{booktabs}
\usepackage{xltabular}
\usepackage{enumitem}
\usepackage{ragged2e}
\usepackage{needspace}
\usepackage{multicol}
\usepackage{imakeidx}

\tcbuselibrary{skins,breakable,raster}

% ---------- 颜色系统（与网页版一致） ----------
\definecolor{inkblue}{HTML}{23407A}
\definecolor{inkblued}{HTML}{1B3260}
\definecolor{inkwarm}{HTML}{8A5A24}
\definecolor{inkbody}{HTML}{16181D}
\definecolor{inkgray}{HTML}{5A5F68}
\definecolor{inkfaint}{HTML}{8B9098}
\definecolor{tintblue}{HTML}{EFF2F8}
\definecolor{tintwarm}{HTML}{F8F2E9}
\definecolor{tintgray}{HTML}{F4F3F0}
\definecolor{rulegray}{HTML}{D8D6D0}
\definecolor{ruledark}{HTML}{B9B6AE}

% ---------- 中文字体：思源宋体 / 思源黑体 / 楷体 ----------
\setCJKmainfont[
  Path        = ./fonts/,
  Extension   = .otf,
  UprightFont = *-Regular,
  BoldFont    = *-Bold,
  ItalicFont  = *-Regular,
  BoldItalicFont = *-Bold,
]{NotoSerifSC}
\setCJKsansfont[
  Path        = ./fonts/,
  Extension   = .otf,
  UprightFont = *-Regular,
  BoldFont    = *-Bold,
]{NotoSansSC}
\setCJKmonofont[
  Path        = ./fonts/,
  Extension   = .otf,
  UprightFont = *-Regular,
  BoldFont    = *-Bold,
]{NotoSansSC}
\newCJKfontfamily[kai]\kaishu{KaiTi}[AutoFakeBold=2.5]
\newCJKfontfamily[light]\cnlight[Path=./fonts/, Extension=.otf, UprightFont=*-Light, BoldFont=*-Regular]{NotoSerifSC}
\newCJKfontfamily[sansmed]\cnsansmed[Path=./fonts/, Extension=.otf, UprightFont=*-Medium, BoldFont=*-Bold]{NotoSansSC}

% ---------- 版心：A4，书籍式双面 ----------
\geometry{
  a4paper,
  inner        = 3.0cm,
  outer        = 3.6cm,
  top          = 2.9cm,
  bottom       = 3.4cm,
  headsep      = 12pt,
  footskip     = 26pt,
  marginparsep = 8pt,
  marginparwidth = 2.0cm,
}
\linespread{1.42}
\setlength{\parskip}{0.30em plus 0.10em minus 0.05em}
\setlength{\emergencystretch}{2em}
\clubpenalty=9000
\widowpenalty=9000
\displaywidowpenalty=9000
\brokenpenalty=6000
\tolerance=1800
\hbadness=2500

% ---------- 基础排字 ----------
\ctexset{
  punct        = quanjiao,
  space        = auto,
  autoindent   = true,
}
\renewcommand{\baselinestretch}{1.0}

% ---------- 小工具 ----------
\newcommand{\kicker}[1]{{\sffamily\scriptsize\color{inkfaint}#1}}
\newcommand{\thead}[1]{{\sffamily\scriptsize\bfseries\color{inkgray}#1}}
\newcommand{\chaptermeta}[1]{}
\newcommand{\hlmark}[1]{\textbf{\color{inkblue}#1}}
\newcommand{\hlwarm}[1]{\textbf{\color{inkwarm}#1}}
\newcommand{\term}[2]{%
  \textbf{\color{inkblued}#1}%
  \ifthenelse{\equal{#2}{}}{}{\,{\footnotesize\itshape\color{inkgray}#2}}%
}

% 章节之间的静默分隔
\newcommand{\bookbreak}{%
  \par\addvspace{16pt}%
  \noindent\hbox to \linewidth{\hfil{\color{ruledark}\footnotesize$\S$}\hfil}%
  \par\addvspace{16pt}}

% 盒子内的小标签
\newcommand{\boxlabel}[2]{%
  \ifthenelse{\equal{#2}{}}{}{%
    {\sffamily\bfseries\scriptsize\color{#1}%
     \raisebox{0.28ex}{\rule{1.05em}{0.55pt}}\hspace{0.55em}#2\par}%
    \vspace{5pt}}}

% ---------- 内容盒子 ----------
\tcbset{bookbox/.style={
  enhanced, breakable, boxrule=0pt, arc=1.2pt, sharp corners=northwest,
  left=11pt, right=11pt, top=10pt, bottom=10pt,
  before skip=13pt, after skip=13pt,
  colframe=white, fontupper=\small,
  parbox=false,
}}

\newtcolorbox{egbox}[1]{bookbox, colback=tintblue,
  borderline west={2.2pt}{0pt}{inkblue},
  before upper={\boxlabel{inkblue}{#1}}}

\newtcolorbox{mythbox}[1]{bookbox, colback=tintwarm,
  borderline west={2.2pt}{0pt}{inkwarm},
  before upper={\boxlabel{inkwarm}{#1}}}

\newtcolorbox{databox}[1]{bookbox, colback=tintgray,
  borderline west={2.2pt}{0pt}{inkgray},
  before upper={\boxlabel{inkgray}{#1}}}

\newtcolorbox{warnbox}[1]{bookbox, colback=tintwarm, boxrule=0.5pt, colframe=inkwarm!35,
  borderline west={2.2pt}{0pt}{inkwarm},
  before upper={\boxlabel{inkwarm}{#1}}}

\newtcolorbox{deepbox}[1]{bookbox, colback=white, boxrule=0.5pt, colframe=rulegray,
  before upper={\boxlabel{inkfaint}{#1}}}

\newtcolorbox{notebox}[1]{bookbox, colback=tintgray,
  before upper={\boxlabel{inkgray}{#1}}}

\newtcolorbox{takeaway}[1]{bookbox, colback=tintblue, boxrule=0.6pt, colframe=inkblue!28,
  arc=3pt, sharp corners=false, left=13pt, right=13pt, top=11pt, bottom=11pt,
  fontupper=\normalsize,
  before upper={\boxlabel{inkblue}{#1}}}

% 误解 → 事实
\newcommand{\mvfrow}[2]{%
  \par\addvspace{4pt}%
  {\sffamily\bfseries\scriptsize\color{inkwarm}#1}\par\nopagebreak\vspace{2pt}#2\par
  \addvspace{4pt}}

% ---------- 导语 / 引语 / 过桥 ----------
\newtcolorbox{ledetext}{bookbox, colback=white,
  left=13pt, right=0pt, top=3pt, bottom=3pt,
  borderline west={2pt}{0pt}{inkblue},
  fontupper=\kaishu\color{inkgray}}

\newtcolorbox{pullquote}[1]{bookbox, colback=white,
  left=15pt, right=6pt, top=7pt, bottom=7pt,
  borderline west={2.8pt}{0pt}{inkblue},
  fontupper=\large\bfseries\color{inkbody},
  after upper={\ifthenelse{\equal{#1}{}}{}{%
    \par\vspace{7pt}{\normalfont\sffamily\scriptsize\color{inkfaint}#1}}}}

\newtcolorbox{bridgebox}{bookbox, colback=white,
  left=0pt, right=0pt, top=9pt, bottom=0pt,
  borderline north={0.4pt}{0pt}{rulegray},
  fontupper=\small\color{inkgray}}

% ---------- 数字组 ----------
\newlength{\statwd}
\newenvironment{statrow}[1]{%
  \par\addvspace{13pt}%
  \setlength{\statwd}{\dimexpr(\linewidth - 1.6em)/#1\relax}%
  \noindent{\color{rulegray}\rule{\linewidth}{0.5pt}}\par\vspace{8pt}%
  \noindent\ignorespaces}
  {\par\vspace{8pt}\noindent{\color{rulegray}\rule{\linewidth}{0.5pt}}\par\addvspace{13pt}}
\newcommand{\statitem}[3]{%
  \begin{minipage}[t]{\statwd}%
    \raggedright
    {\fontsize{19}{21}\selectfont\bfseries\color{inkblue}#1}%
    \ifthenelse{\equal{#2}{}}{}{{\footnotesize\bfseries\color{inkgray}#2}}%
    \par\vspace{4pt}
    {\scriptsize\color{inkgray}\setlength{\baselineskip}{1.25\baselineskip}#3\par}%
  \end{minipage}\hfill\ignorespaces}

% ---------- 编号流程 ----------
\newif\ifflowfirst
\newenvironment{flowlist}{%
  \par\addvspace{12pt}\begingroup\small\setlength{\parindent}{0pt}\flowfirsttrue}
  {\endgroup\par\addvspace{12pt}}
\newcommand{\flowitem}[3]{%
  \ifflowfirst\flowfirstfalse\else
    \vspace{6pt}\noindent{\color{rulegray}\rule{\linewidth}{0.4pt}}\par\vspace{6pt}\fi
  \noindent
  \begin{minipage}[t]{2.3em}\raggedright{\sffamily\scriptsize\color{inkblue}#1}\end{minipage}%
  \begin{minipage}[t]{\dimexpr\linewidth-2.5em\relax}%
    \ifthenelse{\equal{#2}{}}{}{{\bfseries #2}\par\vspace{1pt}}#3%
  \end{minipage}\par}

% ---------- 链条 ----------
\newenvironment{chainflow}{%
  \par\addvspace{12pt}\begingroup\footnotesize\sffamily\setlength{\parindent}{0pt}%
  \raggedright\color{inkgray}}
  {\par\endgroup\addvspace{12pt}}
\newcommand{\chainnode}[1]{%
  \tcbox[on line, boxrule=0.4pt, colframe=rulegray, colback=tintgray,
         arc=1.5pt, left=4pt, right=4pt, top=1.5pt, bottom=1.5pt,
         box align=base]{#1}\allowbreak}
\newcommand{\chainarrow}{\hspace{2.5pt}\raisebox{0.15ex}{$\rightarrow$}\hspace{2.5pt}\allowbreak}

% ---------- 表格 ----------
\newenvironment{tablewrap}{%
  \par\addvspace{13pt}\begingroup
  \footnotesize\setlength{\tabcolsep}{7pt}\renewcommand{\arraystretch}{1.42}%
  \setlength{\parskip}{0pt}\setlength{\LTpre}{0pt}\setlength{\LTpost}{0pt}}
  {\endgroup\par\addvspace{13pt}}
\newenvironment{tablewrapin}{%
  \par\addvspace{7pt}\begingroup
  \setlength{\tabcolsep}{5pt}\renewcommand{\arraystretch}{1.32}%
  \setlength{\parskip}{0pt}}
  {\endgroup\par\addvspace{7pt}}

% ---------- 资产负债表 ----------
\newenvironment{bsheet}{%
  \par\addvspace{13pt}%
  \begin{tcbraster}[raster columns=2, raster equal height=rows,
    raster column skip=7pt, raster left skip=0pt, raster right skip=0pt]}
  {\end{tcbraster}\par\addvspace{13pt}}
\newtcolorbox{bssidebox}[1]{%
  enhanced, boxrule=0.5pt, colframe=rulegray, colback=tintgray, arc=1.2pt,
  left=9pt, right=9pt, top=8pt, bottom=8pt, fontupper=\footnotesize,
  before upper={{\sffamily\bfseries\scriptsize\color{inkgray}#1\par}\vspace{5pt}}}
\newcommand{\bsside}[2]{\begin{bssidebox}{#1}#2\end{bssidebox}}
\newcommand{\bsrow}[2]{%
  \par\noindent#1\hfill #2\par\vspace{2pt}%
  \noindent{\color{rulegray}\rule{\linewidth}{0.3pt}}\par\vspace{2pt}}
\newcommand{\bsrowtotal}[2]{%
  \par\vspace{1pt}\noindent{\color{ruledark}\rule{\linewidth}{0.7pt}}\par\vspace{3pt}%
  \noindent\textbf{#1}\hfill\textbf{#2}\par}

% ---------- 对话 ----------
\newtcolorbox{talkbox}{bookbox, colback=white,
  left=13pt, right=0pt, top=6pt, bottom=6pt,
  borderline west={0.8pt}{0pt}{ruledark}, fontupper=\small}
\newcommand{\talkline}[2]{%
  \par\noindent{\sffamily\bfseries\scriptsize\color{inkfaint}#1}\hspace{0.7em}#2\par\vspace{4pt}}

% ---------- 条形图 ----------
\newenvironment{barfig}{%
  \par\addvspace{12pt}\begingroup\footnotesize\setlength{\parindent}{0pt}}
  {\endgroup\par\addvspace{12pt}}
\newcommand{\barrow}[3]{%
  \par\noindent\makebox[7em][l]{#1}%
  \makebox[0.5\linewidth][l]{%
    {\color{tintgray}\rule[-0.15ex]{0.5\linewidth}{8pt}}%
    \hspace{-0.5\linewidth}%
    {\color{inkblue}\rule[-0.15ex]{\dimexpr0.005\linewidth*#2\relax}{8pt}}}%
  \hfill\makebox[3.6em][r]{\bfseries #3}\par\vspace{5pt}}

% ---------- 章标题微调 ----------
\setlist[itemize]{leftmargin=1.5em, itemsep=2pt, topsep=5pt, parsep=0pt}
\setlist[enumerate]{leftmargin=1.8em, itemsep=2pt, topsep=5pt, parsep=0pt}

% ---------- 页眉页脚 ----------
\fancypagestyle{bookmain}{%
  \fancyhf{}
  \fancyhead[LE]{\kicker{钱是怎么跑起来的}}
  \fancyhead[RO]{\kicker{\nouppercase{\leftmark}}}
  \fancyfoot[C]{\footnotesize\color{inkgray}\thepage}
  \renewcommand{\headrulewidth}{0.4pt}
  \renewcommand{\footrulewidth}{0pt}}

% ---------- 术语表 ----------
\newcommand{\glossletter}[1]{%
  \par\addvspace{10pt}%
  \noindent{\sffamily\bfseries\normalsize\color{inkblue}#1}%
  \hspace{6pt}{\color{rulegray}\leaders\hrule height 0.5pt\hfill}\par\vspace{5pt}}
\newcommand{\glossitem}[4]{%
  \par\addvspace{4pt}\noindent
  \textbf{\color{inkblued}#1}%
  \ifthenelse{\equal{#2}{}}{}{\,{\footnotesize\itshape\color{inkgray}#2}}%
  \ifthenelse{\equal{#4}{}}{}{\hfill{\sffamily\scriptsize\color{inkfaint}第 #4 章}}%
  \par\vspace{1pt}
  {\footnotesize\color{inkgray}#3\par}
  \vspace{3pt}\noindent{\color{rulegray}\rule{\linewidth}{0.3pt}}\par}
\setlength{\columnsep}{22pt}
\setlength{\columnseprule}{0pt}

\addbibresource{refs.bib}
\makeindex[title=索引, intoc=true, columns=2]

\begin{document}

% ================= 前置部分 =================
\frontmatter
\pagestyle{empty}
% ---------- 封面 ----------
\begin{titlepage}
\thispagestyle{empty}
\newgeometry{inner=3.0cm, outer=3.0cm, top=2.8cm, bottom=2.4cm}

\noindent
{\sffamily\scriptsize\bfseries\color{inkfaint} A FIELD GUIDE TO THE FINANCIAL WORLD}

\vspace{8pt}
\noindent{\color{inkblue}\rule{2.4cm}{1.4pt}}

\vspace{1.15cm}

\noindent
{\fontsize{41}{49}\selectfont\bfseries 钱是怎么\textcolor{inkblue}{跑}起来的}

\vspace{0.5cm}

\noindent
{\LARGE\kaishu\color{inkgray} 一本从零开始的金融世界说明书}

\vspace{0.22cm}

\noindent
{\sffamily\normalsize\color{inkfaint} How Money Moves}

\vspace{0.85cm}

\begin{center}
\includegraphics[width=0.94\linewidth]{cover-art.png}
\end{center}

\vfill

\noindent
\begin{minipage}[b]{0.5\linewidth}
{\LARGE 俞孜扬}\quad{\sffamily\scriptsize\color{inkfaint} 著}
\end{minipage}%
\begin{minipage}[b]{0.5\linewidth}
\raggedleft
{\fontsize{27}{27}\selectfont\bfseries\color{inkblue} 2026}\par\vspace{3pt}
{\sffamily\scriptsize\color{inkfaint} 第一版 · 五卷 29 篇 · 约 11.2 万字}\par
{\sffamily\scriptsize\color{inkfaint} 数据与政策截止于 2026 年 8 月}
\end{minipage}

\vspace{0.6cm}
\noindent{\color{inkblue}\rule{\linewidth}{2.2pt}}

\restoregeometry
\end{titlepage}
\cleardoublepage

% ---------- 版权页 ----------
\thispagestyle{empty}
\null\vfill
\noindent
\begin{minipage}{\linewidth}
\footnotesize\color{inkgray}
\setlength{\parskip}{0.7em}

{\normalsize\color{inkbody}\textbf{钱是怎么跑起来的}}\\
一本从零开始的金融世界说明书\quad·\quad How Money Moves

著者　俞孜扬

第一版　2026 年 8 月\quad·\quad 全书五卷 29 篇，约 11.2 万字

\vspace{0.4em}
{\color{rulegray}\rule{\linewidth}{0.4pt}}

\textbf{关于数据}　本书所有数据截至 2026 年 8 月，来源包括中国人民银行与全国银行间同业拆借中心、
美国联邦公开市场委员会、欧洲中央银行、日本银行、国际金融协会（IIF），以及公开财经报道。
完整来源见书末〈参考文献与数据来源〉。金融环境变化很快，请以最新官方数据为准。

\textbf{免责声明}　本书是通识科普读物，不构成任何投资建议，也不构成对任何金融产品、
机构或资产的推荐。书中所有算例均用于说明机制，不可直接套用于个人决策。
读者依据本书内容作出的任何判断与操作，风险自负。

\textbf{关于书中人物}　书中出现的所有人物、公司与交易均为说明机制而虚构，
如与现实中的个人或机构雷同，纯属巧合。引用的公开数据与机构观点除外。

\vspace{0.4em}
{\color{rulegray}\rule{\linewidth}{0.4pt}}

排版：\XeLaTeX{} + ElegantBook v4.7\quad·\quad 正文思源宋体，标题思源黑体，引语楷体\\
封面插画：手绘矢量，孔版套色
\end{minipage}
\vspace{1.5cm}
\cleardoublepage

\pagestyle{bookmain}
\tableofcontents
\cleardoublepage

% ---------- 自序 ----------
\chapter*{自序：一个没法用三句话回答的问题}
\addcontentsline{toc}{chapter}{自序：一个没法用三句话回答的问题}
\markboth{自序}{自序}

去年冬天，一个做平面设计的朋友发微信问我一件事。

她说她妈把攒了二十年的钱存了三年定期，到期一看，利息比上一回少了一大截。她妈很生气，认定是银行坑人。朋友问我：到底是不是银行坑人。

我打了三行字，删掉。又打了五行，又删掉。

因为要回答这个问题，得先说清楚存款利率是怎么定出来的；要说清楚这个，得先说银行靠什么赚钱；要说这个，绕不开净息差；要说净息差为什么被压到 1.41\%，得说贷款需求为什么起不来；而这又牵到房地产的长周期、牵到央行的政策利率、牵到整个全球的利率环境。

任何一环跳过去，答案就变成了「你就信我」。

那天晚上我意识到一件事：这不是一个能用三句话回答的问题。而普通人关于钱的疑问，几乎全都是这样。它们看上去是一个个孤立的困惑，实际上挂在同一台机器上，而那台机器从来没有人完整地拆给他们看过。

\bookbreak

市面上关于钱的书大致分两类。

一类教你怎么赚钱。这类书必须提供确定性，否则卖不动：这个方法有效，这个位置该买，这个比例是对的。

另一类是正经的教科书。它们是对的，写得也严谨，但它们默认你已经有理由读下去。第一章通常从「货币的三大职能」开始，而一个想弄明白自己家存款为什么缩水的人，读到第三页就合上了。

中间那一块是空的。

一本从零讲起、把机器拆开摆给你看、承认自己不知道很多事、并且明确不打算让你发财的书。

这本书想填的就是中间那块。

\bookbreak

写的时候给自己定了五条约束，现在写在前面，方便你随时拿它们来验收。

\begin{flowlist}
\flowitem{01}{只讲一条链，不做名词表}{名词随时能查，因果链没人讲。全书从头到尾只讲一条链：钱是记账，记账要有人信，银行放贷时创造存款，所以钱约等于债，债要付息，付息需要增长，增长需要投资，投资需要给风险定价，定价错误加杠杆等于危机。二十九篇都挂在这条链的某一环上。}
\flowitem{02}{每个术语先讲人话，再给名字}{先说清它解决什么问题、在生活里长什么样，最后才告诉你它叫什么。倒过来的顺序是教科书的顺序，不是理解的顺序。}
\flowitem{03}{所有数字都标时间}{金融世界的数字变得很快，结构变得很慢。数字请当水位刻度读，结构请当河道读。}
\flowitem{04}{不预测}{书里没有一句关于明年涨跌的判断，没有一只被推荐的标的。这不是谨慎，是这本书的立场，理由写在后记二里。}
\flowitem{05}{不知道就说不知道}{全书至少有五处明确写着「目前没人知道答案」。那五处不是偷懒，恰恰是最该被记住的地方。因为骗你的人从不说这句话。}
\end{flowlist}

\bookbreak

还有一件事得说在前面。

这本书写于 2026 年 8 月。这个时点有点特别：负利率时代刚刚结束，日本央行的政策利率回到 1.0\%，是 1995 年以来的最高；黄金在各国央行的储备里第一次超过了美国国债；一轮由人工智能驱动的资本开支正在改写全球的资金流向；而中国这边，房地产还在长周期的调整里，商业银行的净息差压在历史低位。

三年之后回头看，这本书里的每一个数字都会不对。

但那台机器还是那台机器。银行创造货币的方式不会变，利率传导的路径不会变，泡沫的五幕剧本从 1637 年演到今天也没换过。这本书真正想交给你的，是那些不会变的部分。

\bookbreak

最后。

我那位朋友的问题，后来我用了一顿饭的时间讲完。她听到一半的时候放下筷子，说了一句：

「原来不是银行坑人，是整个环境变了。」

这本书要做的就是这件事。只不过讲给更多人听，而且讲得更慢一点。

\vspace{1em}
\noindent\hfill{\kaishu\color{inkgray} 2026 年 8 月}

\cleardoublepage


% ================= 正文 =================
\mainmatter
\pagestyle{bookmain}

\part{地基篇：钱到底是什么}
\chapter{一张一百块的环球旅行}
\chaptermeta{00}

\begin{ledetext}
先不讲定义。跟着一笔钱走一圈。走完你会发现，你以为自己认识它，其实你只认识它的正面。
\end{ledetext}

\section{周五晚上，98 块}

楼下那家"老陈砂锅"，两个菜一碗汤，扫码，98。

手机叮了一声。你以为故事结束了。

真实发生的事是这样的：某个机房的硬盘上，你名下那行数字从 6432.10 变成 6334.10；老陈名下那行从 21847.60 变成 21945.60。没有一张纸移动过，没有一克金属离开过任何地方，没有人在地下金库里把哪块金砖推到另一格。

两行数字，一减一加。

这就是"付钱"的全部物理过程。

\begin{egbox}{先算一笔账}
如果你觉得"钱当然是纸做的"，那看看这两个数：中国的广义货币 M2 是三百多万亿元，而实际在流通中的现金 M0 只有十几万亿。

剩下那 95\% 以上，\textbf{从生下来就没当过纸}。它们出生是数据库里的一行记录，活着是一行记录，死了是这行记录被划掉。

\end{egbox}

这是要拆的第一层壳。第二层是：这 98 块，接下来去了哪儿。

\section{两周，四个国家}

给它装个 GPS。

\begin{flowlist}
  \flowitem{01}{周六早晨 · 老陈的账}{老陈这个月流水 19 万出头。周一他得给冻品供应商打 3 万 2。钱从他开户的城商行跑到供应商开户的股份行，这两家银行之间怎么结清？在央行的账本上。你吃一顿饭，惊动了中央银行。}
  \flowitem{02}{周一 · 供应商}{供应商收到 3 万 2，转手付掉给养殖场的 2 万 4，剩下的还了经营贷的一期。}
  \flowitem{03}{周一下午 · 一笔钱死了}{还给银行的那部分本金，从这个世界上消失了。不是转手，是消失：银行账上的贷款少一块，存款同步少一块，货币总量真的变少了。这句话现在听着像玄学，第 7 章我会把它拆开摆给你看。}
  \flowitem{04}{周三 · 老陈买了理财}{他账上留 8 万周转，拿 5 万买了一只 R2 的理财。产品说明书第 11 页写着投向：同业存单、政策性金融债、中高等级信用债。其中一小块，最后落到江苏一家做储能电芯的厂子手里，成了它三年期融资的一部分。}
  \flowitem{05}{周四 · 出国}{那家电芯厂用这笔融资订了一台德国的极片分切机。人民币在外汇市场换成欧元。这一刻，你的 98 块参与了一次汇率报价。}
  \flowitem{06}{下周 · 德国的工资单}{德国厂商收到欧元，一部分发工资，一部分进财资账户买了短期债券。全球的钱永远在比价：欧洲短债和美国短债的收益率差上 1 个百分点，就有几十亿资金换个方向。}
  \flowitem{07}{某天早上 · 回到你手机里}{你打开 App，看到基金涨了 0.7\%。}
\end{flowlist}

\begin{pullquote}{}
你花掉的 98 块，两周之内在四个国家转了一圈，顺手参与了一次银行间清算、一次债券配置、一次汇率报价和一次设备投资决策。
\end{pullquote}

这台机器是人类造过的连接最紧密的东西。全天候运转，一处出事全球感冒。而绝大多数每天泡在里面的人，对它怎么转一无所知。

这不怪谁。没人教过。

\section{为什么财经新闻你总是看不懂}

因为新闻只报道\emph{事件}，从不解释\emph{管道}。

\begin{talkbox}
  \talkline{新闻}{"美联储宣布，将联邦基金利率目标区间维持在 3.5\% 至 3.75\% 不变，这是今年连续第五次按兵不动。"}
  \talkline{你}{"哦。"}
  \talkline{翻译}{地球上最重要的那个"借钱价格"没动 → 美国国债收益率的锚没动 → 全世界所有资产在被折算成"今天值多少钱"的时候，用的那个除数没动 → 你的基金净值、你换汇的汇率、你银行理财到期给你的数字、你们公司明年敢不敢扩招，全挂在这根链子上。}
\end{talkbox}

新闻给你一个点。理解世界需要一根线。

这本书只干一件事：把线接上。

\section{不发名词表}

名词表随时能查，也随时会忘。真正稀缺的是因果链：为什么 A 会推出 B，B 又必然带来 C。所以全书从头到尾只讲一条链，反复讲，换着角度讲。

\begin{chainflow}
\chainnode{钱是记账}\chainarrow \chainnode{记账要有人信}\chainarrow \chainnode{信任靠制度和抵押品}\chainarrow \chainnode{银行放贷时创造存款}\chainarrow \chainnode{所以"钱"约等于"债"}\chainarrow \chainnode{债要付息}\chainarrow \chainnode{付息需要增长}\chainarrow \chainnode{增长需要投资}\chainarrow \chainnode{投资需要给风险定价}\chainarrow \chainnode{定价错误加杠杆等于危机}\chainarrow \chainnode{危机后重置资产负债表}\chainarrow \chainnode{循环重来}
\end{chainflow}

你在这条链的末端，是一张很小的资产负债表。你能做的不是预测下一次潮汐什么时候来，而是让这张小表在潮汐里活下来，顺便搭上增长的车。

\section{先扔掉一个东西}

\begin{mythbox}{常见误解}
\mvfrow{误解}{"钱是国家印出来的。政府想多花钱就多印一点，所以才有通货膨胀。"}
\mvfrow{事实}{现代经济里绝大部分货币不是被"印"出来的，是被商业银行"借"出来的。银行批给你 300 万房贷的那一刻，它没从任何人的账户里搬走一分钱，只是在自己账本上同时写下两行：资产端"贷款 300 万"，负债端"你的账户里有 300 万"。\textbf{新的钱就此诞生。}}
这不是什么阴谋论。英格兰银行 2014 年就在官方论文里把它写清楚了，欧洲央行和美联储的表述也一致。是很多教科书没跟上，这是教科书的问题。第 7 章会把它彻底拆开。

\end{mythbox}

\section{四层楼}

越往上越抽象，但每层都踩在下一层的地基上。你可以按顺序爬，也可以直接跳到关心的那层，每章开头会说清它需要什么前置。

\begin{tablewrap}
\begin{xltabular}{\linewidth}{>{\raggedright\arraybackslash}p{0.20\linewidth}XX}
\toprule
\thead{层} & \thead{回答的问题} & \thead{写给谁} \\
\midrule
\textbf{01 · 地基篇} & 钱到底是什么？它凭什么有价值？ & 完全零基础。读完能听懂饭桌上八成的经济话题 \\
\textbf{02 · 管道篇} & 钱怎么被造出来、怎么流动、怎么定价？ & 看得懂新闻标题，但不知道背后的机制 \\
\textbf{03 · 天气篇} & 为什么总有泡沫、危机和周期？ & 想弄明白"为什么每隔十年就来一次" \\
\textbf{04 · 落地篇} & 这些跟我的钱包有什么关系？ & 所有人。尤其是第 23 章 \\
\bottomrule
\end{xltabular}
\end{tablewrap}

\begin{databox}{三条约定}
\textbf{一，遇到英文缩写别慌。}这行黑话密度极高，但每个缩写背后都是一句大白话。本书出现的每个术语都先讲人话，再给名字。

\textbf{二，所有数字都标了时间。}本书写于 2026 年 8 月。金融世界的数字变得很快，结构变得很慢。数字当水位刻度读，结构当河道读。

\textbf{三，不预测。}任何声称能告诉你明年涨跌的人，不是在骗你就是在骗自己。这本书只负责让你看懂机器怎么转，判断留给你。

\end{databox}

\begin{takeaway}{本章一句话}
钱不是纸，是账本上的一行记录。金融不是赌场，是把"今天的钱"和"明天的钱"搬来搬去的一整套管道。

\end{takeaway}

\begin{bridgebox}
\textbf{下一章}既然钱只是一行数字，它凭什么能换来一碗砂锅？这事得从太平洋上一块沉进海里的石头说起。
\end{bridgebox}


\chapter{钱不是财富，钱是一套记账系统}
\chaptermeta{01}

\begin{ledetext}
序章那 98 块只是两行数字。可老板真的把面端上来了。这一章问的就是这件事：一行数字凭什么换来一碗面。
\end{ledetext}

\section{有一块石头沉在海底，它照样是钱}

太平洋上有个岛叫雅浦。岛上的钱是石头。

石灰岩凿的圆盘，中间开个洞，大的直径三米多，重四吨。雅浦不产石灰岩，得划独木舟去四百公里外的帕劳采，再运回来。翻过船，死过人。所以一块石头值多少，不光看大小，还看它是怎么运回来的。

石头太重，搬不动。于是付款的时候石头不动。一户人家把石头付给另一户，做法是当着一群人的面说一句"这块归你了"。石头还立在原地，青苔照长，主人换了。

1910 年前后，人类学家弗内斯在岛上碰到一户特别富的人家。全村没人见过他家的钱。

那块巨石在运回来的路上遇了风暴，沉了。

船员作证：石头确实凿出来了，尺寸成色都好，就是现在躺在海底。全村认了。这块谁也没见过、也捞不上来的石头，被这家人当财产传了几代，还照样拿出来付账。

\begin{pullquote}{}
钱不是东西。钱是一群人共同维护的一段记忆。
\end{pullquote}

1932 年，法国怕美国放弃金本位，要求把存在纽约联邦储备银行的美元换成黄金，黄金不用运回巴黎。

美联储照办了。工作人员走进金库，把相应数量的金条从一个抽屉挪到另一个抽屉，换上写着"法国"的标签。

金条一米没动。第二天报纸开始报道美国黄金外流，美元承压。

弗里德曼后来点破：纽约金库里换标签，和雅浦岛给沉在海底的石头换主人，是同一件事。区别只在于我们这边有抽屉。

\begin{egbox}{举个栗子}
假设一个村子，十七户人家，没有一张钞票，村口挂着一块黑板。

李建国帮邻居修了三天屋顶，黑板上记一笔："王有福欠李建国 3 个工。"王有福给第三家修拖拉机，再记一笔。月底村会计做一次轧差，中间那些环节全划掉，只剩下谁真欠谁多少。

这个村子有钱吗。没有一张纸，但它有价格、有债务、有清算、有支付。

黑板就是货币系统，钞票是黑板的便携版。把村子放大到 14 亿人，黑板换成机房，村会计换成商业银行和央行。

\end{egbox}

\section{教科书上那个起源故事是编的}

故事你一定听过。很久以前人们物物交换，你有一头牛想换斧头，可拿斧头的人只想要鞋，交易做不成。这叫\term{双重巧合难题}{double coincidence of wants}\index{055@双重巧合难题}。为了解决它，人类发明了货币。

出自亚当·斯密 1776 年的《国富论》，写得漂亮，被抄了两百五十年。

它是编的。

人类学家找了一百多年，没找到任何一个靠物物交换运转的社会。1985 年剑桥的汉弗莱写过一句被引到烂的话：从来没有人描述过一个纯粹的物物交换经济，更别说从里面长出货币。

田野里能看到的物物交换，基本发生在陌生人和敌人之间。彼此不信任，才必须一手交钱一手交货。

村子内部靠赊账。你今天帮我收麦子，下个月我杀猪给你留条腿。谁欠谁多少全在人情里，攒到某个时候一起算。

考古的方向也对得上。美索不达米亚的泥板早在公元前三千年就密密麻麻记着"某人欠神庙多少大麦""折合多少银子"，而最早的铸币出现在公元前七世纪的吕底亚。\hlmark{记账比铸币早了两千多年。}

\begin{mythbox}{常见误解}
\mvfrow{误解}{"人类先物物交换，后来嫌麻烦发明了货币，所以钱本质上是一种方便的商品。"}
\mvfrow{事实}{顺序反了。\textbf{先有债务和记账，后有铸币。}铸币解决的是跟陌生人、跟远方、跟军队做交易的问题，不是从村口的以物易物里自然长出来的。"双重巧合难题"是一个功能论证，它说明货币解决了什么，它不是历史记录。因为伞能挡雨就断定伞是雨发明的，逻辑是一样的。}
诚实地补一句：细节上学界还在吵。格雷伯的《债：第一个五千年》把这套说法推给了大众，其中一些史料解读经济史学界并不买账。但"信用早于铸币"这条主干，现在很少有人反对。

\end{mythbox}

\section{三个功能，一个老大}

\begin{flowlist}
  \flowitem{01}{交换媒介\enterm{medium of exchange}}{你拿它换东西。最直观，也最容易被替代。贝壳、盐、烟草、方便面都干过这活儿。}
  \flowitem{02}{记账单位\enterm{unit of account}}{价格、工资、合同、房贷、债券票息，都用它标价。它是这个世界的度量衡。}
  \flowitem{03}{价值储藏\enterm{store of value}}{今天不花，留到明天还有用。三个功能里最脆的一个，通胀会慢慢啃掉它。}
\end{flowlist}

这三个不是平级的。\textbf{记账单位才是老大。}

一国货币真正死掉的那一刻，不是大家不再拿它买东西，是商家不再拿它标价。

恶性通胀的国家里，这两件事是分开发生的。价签早改成了美元，结账还收本币，按当天黑市汇率折。这时候本币只剩一个找零工具的身份，度量衡那把椅子已经让出去了。1923 年的德国走得更远：人们干脆用香烟和实物直接换。

反过来，只要工资、合同和债务还用某种货币标价，它就有极强的粘性。你欠银行的 300 万房贷是用人民币写的，你就得想办法弄到人民币。

\begin{egbox}{举个栗子}
1945 年，拉德福德写了篇论文，讲他待了两年的德国战俘营。营里的钱是香烟。

一罐牛肉 8 支烟，洗一件衬衫 2 支。有人在营区两头跑价差，有人开"咖啡馆"赊账供应，用烟结算。

红十字包裹一到，香烟供给暴增，所有东西的烟价当天上涨。这是通胀。断供几周，大家把烟抽掉了，物价开始跌。这是通缩。

2010 年代美国不少监狱禁烟，囚犯的通货换成了方便面。

香烟和方便面成为货币，不是因为好抽好吃。是因为大家开始拿它标价。

\end{egbox}

\section{老陈那张一百块凭什么}

老陈在昆明开面馆，二十六年了。他收人民币，不收比特币，也不收你写的欠条。

问他为什么，他会说这不废话吗。

不废话。他敢收，是因为四件事同时成立。

\begin{flowlist}
  \flowitem{01}{国家拿它收税}{最硬的一根。老陈每季度要交增值税和附加，雇的三个人要交个税，全得用人民币。他就算特别喜欢黄金，也得先把黄金换成人民币再去交。几亿人被迫弄到手的东西，天然有需求。}
  \flowitem{02}{法偿性\enterm{legal tender}}{法律规定用人民币偿还境内债务，债权人不得拒收。中国人民银行公开整治过拒收现金，罚过商户。这条的真实含义是：老陈手里的钱，法院认。}
  \flowitem{03}{网络效应}{老陈收人民币，因为他房东收、员工收、菜市场收。货币跟语言一样，用的人越多越好用，也越难改。你可以自己发明一种语言，但你没法让整条街的菜贩子学。}
  \flowitem{04}{央行盯着它的购买力}{货币是有质量的产品，质量就是购买力。2026 年 8 月，美联储把联邦基金利率维持在 \textbf{3.50\%–3.75\%}，年内第五次按兵不动，理由是通胀还高。它宁可让经济难受，也要保住这张纸的信用。这份信誉攒了几十年，砸掉要不了几年。}
\end{flowlist}

\begin{mythbox}{常见误解}
\mvfrow{误解}{"钱有价值是因为国家有黄金储备。国库里的金子没了，钱就废了。"}
\mvfrow{事实}{1971 年 8 月 15 日美国关掉黄金窗口之后，\textbf{世界上没有任何一种主要货币能兑换黄金}。中国人民银行不承诺给你换金条，美联储也不承诺。黄金今天是各国央行的储备资产之一，是压舱石，不是发钞的抵押物。撑住老陈那张一百块的，是上面那四条：税收、法律、网络、信誉。}
\end{mythbox}

\section{金本位回不去了}

怀念金本位的人通常这么说：钞票能换金子，政府就没法乱来。

前半句对。后半句的代价很大。

\begin{flowlist}
  \flowitem{01}{货币供给由挖矿速度定，不由经济需要定}{1873 到 1896 年，美国物价跌了二十多年。农民卖粮的钱一年比一年少，欠银行的债还是那个数字。1896 年布莱恩喊"你们不能把人类钉在黄金的十字架上"，喊的就是这个。反过来，加州淘金和南非金矿又把物价推上去。把货币的松紧交给地质运气，这不叫纪律，叫抽签。}
  \flowitem{02}{危机沿着固定汇率传染}{一国货币被挤兑，央行只能加息紧缩来保住兑换承诺，哪怕国内已经是大萧条。经济史里有条刺眼的规律：越早脱离金本位的国家，复苏越早。英国 1931 年 9 月脱钩，比死守到最后的国家早好几年爬出来。}
\end{flowlist}

1944 年的布雷顿森林是金本位的最后一版：美元按 35 美元一盎司锚黄金，别的货币锚美元。撑了不到三十年。1971 年，美国的黄金已经远远盖不住海外的美元债权，尼克松关了窗口。1973 年，主要货币开始浮动。

从那天起全世界的钱都是\term{法定货币}{fiat money}\index{013@法定货币}。fiat 是拉丁文，意思是"就这么定了"。

\begin{databox}{2026 数据：黄金回来了，但不是以你想的方式}
欧洲央行 2026 年 6 月的《欧元的国际角色》报告：截至 2025 年底，\textbf{黄金已经超过美国国债，成为全球央行的第一大储备资产}。

\begin{barfig}
  \barrow{黄金}{27}{27\%}
  \barrow{美国国债}{22}{22\%}
  \barrow{其他美元资产}{20}{20\%}
  \barrow{欧元资产}{15}{15\%}
\end{barfig}

黄金占比一年之间从 20\% 升到 27\%，美债从 25\% 降到 22\%。但另一个数字要一起读：\textbf{美元计价资产合计仍占 42\%}，还是最大的单一货币阵营。这是渐进的多元化，不是改朝换代。

其他口径：全球央行黄金储备超 \textbf{3.6 万吨}，逼近布雷顿森林鼎盛期的约 3.8 万吨；2025 年央行净购金 \textbf{850 吨}；金价 2026 年 1 月一度接近 \textbf{5600 美元/盎司}。2025 年单一最大买家不是哪个国家，是稳定币发行商 Tether，全年买了超过 100 吨。

\end{databox}

看到这组数字，很多人的第一反应是要回金本位了。这个推论把储备和发钞基础混成了一件事。

央行买黄金，买的是一份保险：黄金不是任何国家的负债，冻结不了，也稀释不掉。但没有一家央行承诺你能拿本币换它。储备是资产端的选择，发钞靠的是负债端的信用。

\section{谁欠你的}

下面这一段是全书的伏笔，读慢一点。

\textbf{每一种钱都是某个人的负债。}你说自己"有钱"，准确的说法是"有人欠你钱"。剩下的问题只有一个：欠你的是谁。

\begin{tablewrap}
\begin{xltabular}{\linewidth}{>{\raggedright\arraybackslash}p{0.20\linewidth}XXX}
\toprule
\bthead{你手里的东西} & \bthead{它本质上是} & \bthead{谁欠你} & \bthead{出事了谁兜} \\
\midrule
\textbf{口袋里的现金} & 中央银行的负债 & 中国人民银行 & 境内最终的钱，不用别人兜 \\
\textbf{银行卡里的存款} & 商业银行的负债 & 你的开户行 & 存款保险，50 万元以内本息全额 \\
\textbf{微信 / 支付宝余额} & 支付机构的负债 & 支付机构 & 客户备付金 100\% 集中交存央行 \\
\textbf{数字人民币} & 中央银行的\textbf{直接}负债 & 中国人民银行 & 等同于现金 \\
\textbf{货币基金份额} & 基金资产的份额，不是存款 & 严格说没人"欠"你 & 没有存款保险，净值理论上会波动 \\
\textbf{美元稳定币} & 发行公司的负债 & 那家发行公司 & 看它的储备资产和监管框架 \\
\bottomrule
\end{xltabular}
\end{tablewrap}

老陈的手机里有三笔钱。银行卡里 27 万，是那家银行欠他的。微信零钱 4300 块，是支付机构欠他的，对应的备付金 100\% 交存在央行。抽屉里压着 1200 块现金，是央行欠他的。

三笔钱在他眼里长得一模一样，法律上隔着三层。

这也解释了数字人民币为什么不是"微信支付的国家队版本"。微信余额是支付机构欠你的，数字人民币是央行直接欠你的，中间少了一层机构信用。

整个体系是一座用负债搭起来的塔。塔尖是央行的负债（现金和准备金），中间是商业银行的负债（存款），再往下是各类机构的负债（余额、稳定币）。越往上越安全，越往下越好用。

老陈抽屉里那张一百块，写在央行的负债栏里。

\begin{chainflow}
\chainnode{央行的负债}\chainarrow \chainnode{你的现金}\chainarrow \chainnode{银行的负债}\chainarrow \chainnode{你的存款}\chainarrow \chainnode{机构的负债}\chainarrow \chainnode{你的余额}
\end{chainflow}

\section{比特币卡在哪一步}

钱只是记账，那能不能不要国家那本账？\term{比特币}{Bitcoin}\index{004@比特币}要回答的就是这个：用密码学和一个分布式网络，让一群互不信任的人维护同一本账。

技术上它做到了。

当日常货币它不行，卡在三个地方。

\begin{flowlist}
  \flowitem{01}{波动}{一天涨跌 5\% 是常态。老陈的一碗牛肉面标 0.00021 个比特币，明天早上他得重新贴一次价签，后天再贴一次。}
  \flowitem{02}{结算}{链上要等确认，手续费还随网络拥堵跳。买碗面付几美元手续费，这笔账算不过来。}
  \flowitem{03}{记账单位失效，最致命的一条}{就算在币圈内部，人们也是拿美元给比特币报价，不是拿比特币给面包报价。它没坐上度量衡那把椅子。另外 2100 万枚的上限意味着，它要真成了主流货币，经济在长它不长，结果是持续通缩，金本位那个老毛病的加强版。}
\end{flowlist}

这里讨论的是它能不能当日常货币，不是它值多少钱。本书不预测任何资产涨跌，也不给买卖建议。

\term{稳定币}{stablecoin}\index{065@稳定币}走的是反方向。它放弃自己当货币，改成锚美元，1 枚兑 1 美元。发行方收你的美元，拿去买短期美债和回购，在链上给你发一枚代币。

拿上面那张表看它：\hlmark{稳定币就是某个私人机构的负债}。你把美元借给一家公司，换回一张能在链上转来转去的欠条。技术是新的，法律结构老得像 19 世纪的私人银行券。

老陈的儿子做跨境电商，收过 USDT。他不关心这些，他只关心到账是三秒还是三天。

\begin{databox}{2026 数据}
美国《GENIUS Act》2025 年 7 月签署生效，给美元稳定币建了联邦层面的监管框架，核心是储备资产的质量和透明度。到 2026 年 8 月，全球稳定币市值约 \textbf{3000 亿美元}量级。

这数字听着大，放进坐标系看：中国广义货币 M2 是三百多万亿元量级。稳定币真正的战场不是取代存款，是跨境支付和 7×24 结算。传统跨境汇款要两三天，扒好几层手续费，链上几秒钟就到。

\end{databox}

\begin{warnbox}{小心}
稳定币稳的是价格锚，不是风险。2023 年 3 月硅谷银行倒掉的时候，某主流美元稳定币因为有一部分储备存在那家银行，价格一度掉到 0.87 美元。\textbf{链上的欠条，安全性取决于链下的资产。}

\end{warnbox}

\bookbreak

回到老陈那碗面。

他敢把面端上来，是因为他相信那行数字明天还有人认。他之所以相信，是因为国家拿它收税、法院认它清偿、整条街都在用它，还有一家央行盯着它的购买力。

雅浦岛现在还留着几千块 Rai。绝大多数从被凿出来那天起就没挪过窝，连沉在海底的那块也还在账上。

\begin{takeaway}{本章一句话}
钱不是财富本身，是一套记录"谁欠谁多少"的公共账本。它有价值不是因为背后有黄金，是因为税收、法律、网络和信誉在同时撑着。

\begin{itemize}
  \item 先有赊账和记账，后有铸币。物物交换起源说是功能论证，不是历史。
  \item 三个功能里记账单位最根本。商家不再拿它标价，货币才算真死。
  \item 每一种钱都是某人的负债。现金是央行的，存款是银行的，余额是支付机构的，稳定币是发行公司的。
\end{itemize}

\end{takeaway}

\begin{bridgebox}
\textbf{下一章}既然你卡里的钱只是银行欠你的，那这笔钱现在人在哪儿？打开银行金库，为什么找不到属于你的那一沓？
\end{bridgebox}


\input{chapters/ch-02}
\chapter{利息：时间是有价格的}
\chaptermeta{03}

\begin{ledetext}
银行靠利差活着。可利息本身是什么？为什么钱借出去就该收回更多？答案不在钱上，在时间上。
\end{ledetext}

\section{亚里士多德算错了}

亚里士多德对利息意见很大。他说钱是死的，不会像牛羊那样下崽，靠钱生钱违反自然。

他算错的地方在于，利息本来就不是钱在下崽。

利息是你为"提前拿到"付的钱。拆开看，任何一个利率里都塞着几样东西。

\begin{flowlist}
  \flowitem{01}{时间偏好，也就是人天生不耐}{现在给赵磊 8000 块，和一年后给他 8000 块，他选现在。不是因为他算过通胀，是因为人就这样：现在的确定，胜过将来的可能。要他愿意等，得加钱。这笔加钱是时间本身的价格。}
  \flowitem{02}{通胀预期}{一年后那 8000 块买的东西大概率少一点。这部分补的不是等待，是购买力缩水。没人愿意借出一头牛，一年后收回一只羊。}
  \flowitem{03}{违约风险}{借给财政部，和借给赵磊他姐在南京开的那家卤味店，不可能一个价。店可能明年就关了。多收的那部分，是为"可能收不回来"买单。}
  \flowitem{04}{流动性溢价}{三年定期比活期利率高，多出来的部分补的是这三年你不能动它。急用钱的时候，钱在手上和钱在合同里是两样东西。}
\end{flowlist}

赵磊他姐去年借了 41 万经营贷，年化 6.5\%。这 6.5\% 是一层一层摞上去的。

\begin{tablewrap}
\begin{xltabular}{\linewidth}{>{\raggedright\arraybackslash}p{0.20\linewidth}Xrr}
\toprule
\thead{一层一层往上加} & \thead{这一层在补偿什么} & \thead{加多少} & \thead{累计} \\
\midrule
\textbf{纯粹的不耐} & 人就是不愿意等 & +0.5\% & 0.5\% \\
\textbf{通胀补偿} & 一年后同样的钱买得更少 & +1.0\% & 1.5\% \\
\textbf{\textbf{= 名义无风险利率}} & 对照：7 天逆回购 1.4\% & — & \textbf{约 1.5\%} \\
\textbf{流动性溢价} & 这笔钱三年抽不回来 & +1.0\% & 2.5\% \\
\textbf{违约风险溢价} & 一百家小店里总有几家关门 & +4.0\% & \textbf{约 6.5\%} \\
\bottomrule
\end{xltabular}
\end{tablewrap}

这是示意性拆解，不是任何机构的真实报价。风险怎么定价，第 13 章专门讲。

利率从来不是一个数字，是一叠补偿的总和。看到 6.5\% 就骂高利贷，跟看到 1.5\% 就说钱不值钱，犯的是同一个错：只看总和，不看结构。

\section{你妈的定期}

合同上写的那个数叫\term{名义利率}{nominal rate}\index{039@名义利率}。\term{实际利率}{real rate}\index{051@实际利率}约等于名义利率减通胀。这叫费雪关系。

\begin{egbox}{举个栗子}
赵磊他妈在盐城，把 11.6 万存了一年定期，利率 2\%。到期账上躺着 118,320 元。她挺满意，说你看钱变多了。

如果这一年通胀是 3\%，去年花 11.6 万能买到的那一篮子东西，今年要 119,480 元。

钱多了 2320 块，能买的东西少了 1100 出头。实际收益率是 −0.97\%。

银行没骗她。账单上每一分利息都是真的。

偷东西的那位不开对账单。

\end{egbox}

以后看到任何一个收益率，先在心里做一次减法。名义数字是给你看的，实际数字才是你赚的。

\section{这本书最重要的一个除法}

年利率 5\%，今天的 95.24 元一年后变成 100 元。

反过来读：一年后的 100 元，今天只值 95.24 元。

这个折回今天的动作叫\term{贴现}{discounting}\index{059@贴现}，那个 5\% 叫\term{贴现率}{discount rate}\index{060@贴现率}，算出来的 95.24 叫\term{现值}{present value}\index{068@现值}。

把时间拉长，威力就出来了。

\begin{tablewrap}
\begin{xltabular}{\linewidth}{>{\raggedright\arraybackslash}p{0.20\linewidth}Xr}
\toprule
\thead{未来的 100 元} & \thead{怎么折} & \thead{今天值多少} \\
\midrule
\textbf{1 年后} & 100 ÷ 1.05 & 95.24 元 \\
\textbf{10 年后} & 100 ÷ 1.05\textsuperscript{10} & 61.39 元 \\
\textbf{30 年后} & 100 ÷ 1.05\textsuperscript{30} & 23.14 元 \\
\bottomrule
\end{xltabular}
\end{tablewrap}

世界上几乎所有金融资产（股票、债券、房子、一家公司、一条收费公路）本质上都是一串未来的现金流。给它们定价，就是把这串现金流一笔一笔折回今天，再加起来。

分母上坐着的那个贴现率，来自利率。

\begin{pullquote}{}
利率一动，所有资产的价格都要重算。\textbf{没有例外。}
\end{pullquote}

\begin{egbox}{举个栗子：分母变一下，价格砍一半}
假设有项资产每年稳定给你 3.8 万，永远给下去。它值多少钱？

价格等于每年现金流除以贴现率。

贴现率 2\%：3.8 万 ÷ 0.02 = \textbf{190 万}。\\ 贴现率 4\%：3.8 万 ÷ 0.04 = \textbf{95 万}。

现金流一分没少，经营一切照旧，\hlmark{价格砍掉一半}。只因为分母变了。

2022 年以来全球资产价格的剧烈波动，核心就是这个机制。美联储把政策利率从接近零抬到 5\% 以上，全球股债一起跌。跌得最狠的是那些现金流全在很远的将来的资产：不赚钱的成长股、长久期债券。分母抬高对遥远的现金流最致命，30 年后的 100 元本来就只值 23 块，除数一大就不剩什么了。

上一章那道题的答案也在这儿。一张 30 年期、票息 3\% 的债券，市场利率上行 1 个百分点，价格大约跌 20\%。硅谷银行资产端那 150 亿美元浮亏就是这么来的。它不是被人骗了，它是被一个除法干掉的。

\end{egbox}

\section{72 除以 3.5}

贴现是把未来折回今天。反过来把今天推到未来，就是\term{复利}{compound interest}\index{018@复利}：利息也开始生利息。

有个心算工具叫 72 法则。72 除以年化收益率，得到本金翻倍要多少年。

72 ÷ 3.5 ≈ 20.6 年。

精确算是 20.1 年。72 法则在 6\% 到 10\% 之间最准，利率低的时候会略微高估。这点误差不影响它当心算工具，但它提醒了一件事：所有关于复利的漂亮说法，对收益率假设都极其敏感。

赵磊手上有 4.6 万闲钱。按 3.5\% 滚 30 年是 129,122 元，按 2\% 滚 30 年是 83,306 元，差 45,816 元。

而 3.5\% 和 2\% 在直觉上是"差不多的小数字"。

复利的真相和关于复利的谎言，第 22 章掰扯。

\section{两个百分点，128 万}

赵磊今年 32 岁，在南京做产品经理，去年买房贷了 300 万，30 年，等额本息。

\begin{tablewrap}
\begin{xltabular}{\linewidth}{>{\raggedright\arraybackslash}p{0.20\linewidth}rrr}
\toprule
\thead{贷款条件} & \thead{月供} & \thead{30 年总还款} & \thead{其中利息} \\
\midrule
\textbf{300 万 / 30 年 / 年利率 \textbf{3.5\%}} & 约 13,469 元 & 约 484.9 万 & 约 \textbf{184.9 万} \\
\textbf{300 万 / 30 年 / 年利率 \textbf{5.5\%}} & 约 17,036 元 & 约 613.3 万 & 约 \textbf{313.3 万} \\
\textbf{差额} & 每月多 3,567 元 & 多还 128.4 万 & 多付 \textbf{128.4 万} \\
\bottomrule
\end{xltabular}
\end{tablewrap}

盯住最后一行。

纸面上只差 2 个百分点，听起来像四舍五入就没了的小数。换成现实，是赵磊每个月多掏 3567 块，三十年多付出去 128.4 万。在南京这是一辆不错的车加一整套装修，在他妈住的盐城，这是一套小房子。

还有个数更刺眼：3.5\% 那一栏，他借 300 万，最后要还 485 万。5.5\% 那一栏，他还的钱是借的两倍多。

他买的不只是房子。他还买了三十年的等待权，而这个等待权明码标价。

\begin{warnbox}{小心}
看完这张表，很多人第一反应是赶紧提前还贷。这个决定取决于你的贷款利率、你能拿到的其他资产回报、你的现金流稳定性和应急储备，是一道个人化的比较题，没有标准答案。\textbf{本书不提供投资或借贷建议}，只负责让你看懂这张表在说什么。

\end{warnbox}

\section{全世界的宗教都恨过它}

收息现在天经地义。它曾经是重罪。

《旧约》禁止向同胞取利。中世纪教会把放贷取息定为罪，1179 年第三次拉特兰会议规定放高利贷者不得领圣餐，死后不许葬进教会墓地。伊斯兰教法至今禁 riba，所以伊斯兰金融要用买卖、租赁、合伙分成的结构绕开固定利息。

反应这么激烈是有原因的。

古代社会里借钱的人几乎都是走投无路的穷人。家里歉收，孩子快饿死了，才去借粮。这时候有人在你身上加三成，你还不上就拿女儿抵债。这不是金融，是趁火打劫。禁令针对的是这种\term{消费性借贷}{consumption loan}\index{071@消费性借贷}。

转折发生在借钱的用途变了以后。

商人借钱买一船货，工厂主借钱买一台织机，借来的钱产生了新的产出。利息的性质跟着变了：它不再是从穷人碗里抢饭，是出资人分掉自己出资创造出来的那部分增量。

现代社会接受的是这套逻辑：资本有机会成本，风险该被补偿，钱流向回报更高的地方有效率。

那条古老的道德直觉并没有消失。今天各国给消费贷设利率上限、打击非法放贷、强制明示年化利率，护的还是那个走投无路才借钱的人。

\section{从倒贴，到 1.0\%}

利息为正听起来天经地义。人类历史上绝大多数时间确实如此。

2014 年，事情变得很奇怪。

那一年欧洲央行把存款便利利率降到 −0.1\%。2016 年日本银行跟进。丹麦、瑞士、瑞典也进过负值区间。

负利率的字面意思是：你借钱给别人，还要倒贴保管费。把 100 万借给德国政府，十年后拿回 99 万。2020 年，全球负收益率债券的规模一度超过 17 万亿美元。

这在人类金融史上极其罕见。它是通缩、老龄化和长期需求不足逼出来的：央行想刺激经济，发现利率已经到零，只好把它按到水面以下。

这场实验现在结束了。日本这两年走出来的路径，是最值得看的一条线。

\begin{flowlist}
  \flowitem{01}{2024 年 3 月}{结束负利率与收益率曲线控制，政策利率从 −0.1\% 抬到 0\textasciitilde{}0.1\%。十七年来第一次加息。}
  \flowitem{02}{2024 年 7 月}{加到 \textbf{0.25\%}。市场剧烈震荡，套息交易开始被迫平仓。}
  \flowitem{03}{2025 年 1 月}{加到 \textbf{0.5\%}。}
  \flowitem{04}{2025 年 12 月}{加到 \textbf{0.75\%}。}
  \flowitem{05}{2026 年 6 月 16 日}{加到 \textbf{1.0\%}，\hlmark{1995 年 9 月以来的最高水平}。底气来自工资和物价：2026 年"春斗"薪资涨幅预期约 5.02\%，5 月 PPI 同比 6.3\%。}
\end{flowlist}

欧洲央行也在 2026 年 6 月 11 日上调了存款便利利率，三年来第一次加息。

\begin{statrow}{4}
  \statitem{3.50}{–3.75\%}{美国联邦基金利率目标区间}
  \statitem{3.0}{\%}{中国 1 年期 LPR}
  \statitem{3.5}{\%}{中国 5 年期以上 LPR}
  \statitem{1.0}{\%}{日本政策利率}
\end{statrow}

\begin{databox}{2026 数据（截至 2026 年 8 月）}
美联储的联邦基金利率目标区间维持在 \textbf{3.50\%–3.75\%}，年内已连续 5 次按兵不动，理由是通胀仍然高企。美国 10 年期国债收益率在 8 月底升至 2025 年 1 月以来的最高水平。

中国方面，8 月 20 日 LPR 报价 1 年期 \textbf{3.0\%}、5 年期以上 \textbf{3.5\%}，均已连续多月持平；7 天期逆回购政策利率 \textbf{1.4\%}；货币政策基调是"适度宽松"。日本政策利率 \textbf{1.0\%}。

\end{databox}

同一天，同一个地球，钱的价格从 1.0\% 到 3.75\%，差了将近三倍。

这个差价不会安静待着。在日本借 1\% 的钱，换成美元买 3.5\% 的资产，中间的利差是白拿的，只要汇率别动。这叫\term{套息交易}{carry trade}\index{058@套息交易}，是全球市场最大的一条暗流，也是 2024 年 8 月那场全球股市急跌的导火索。它怎么转、为什么会突然反转，第 12 章算这笔账。

\begin{mythbox}{常见误解}
\mvfrow{误解}{"利率当然是越低越好，借钱便宜，大家都开心。"}
\mvfrow{事实}{这个说法是错的。利率是经济的体温，长期贴地通常说明病人虚弱，不是身体好。它意味着需求不足、投资回报低到央行必须一直踩着油门。\\ 代价也是实的：低利率鼓励人借钱去买存量资产（房子、股票）而不是投新产能，把资产价格推上去，对有房有股的人是福音，对想上车的年轻人是灾难；它让本该出清的僵尸企业靠低息续命；它压垮靠利息生活的退休者，以及必须用利息覆盖长期给付的养老金和保险公司。日本和欧洲用了整整十年才从这个坑里爬出来。\\ 至于低利率到底多大程度上推高了资产价格、多大程度上只是反映了增长本身在放缓，学界吵了十几年，目前没有共识。}
\end{mythbox}

\begin{mythbox}{常见误解}
\mvfrow{误解}{"存款利率一降再降，是银行黑心，赚我们的钱。"}
\mvfrow{事实}{因果方向反了。2026 年二季度末，中国商业银行净息差是 \textbf{1.41\%}，历史低位。银行两头被压：资产端的贷款利率跟着 LPR 和 1.4\% 的政策利率往下走，负债端的存款利率不跟着降，利差会直接变负。\\ 存款利率是被贷款利率拖着走的，而贷款利率反映的是整个经济的资金供求和政策取向。这不是说银行是慈善机构，是说这件事的因果链条，和直觉指的方向正好相反。传导的完整机制，第 9 章讲。}
\end{mythbox}

\bookbreak

学会了贴现，后面所有章节的门都开了。债券为什么涨跌、股票为什么被杀估值、房子值不值这个价、一家公司该值多少钱，全是同一个除法。

赵磊的房贷合同上写着一个年利率。

这个数字里有多少是给时间的，多少是给通胀的，多少是给"他可能失业"这件事的，合同上不写。

\begin{takeaway}{本章一句话}
利息不是钱自己生钱，是时间偏好、通胀预期、违约风险和流动性溢价叠出来的价格。它反过来用就是贴现率，而贴现率一动，世界上每一项资产都要重新标价。

\begin{itemize}
  \item 实际利率约等于名义利率减通胀。存款 2\%、通胀 3\%，一年白送出去 1100 块。
  \item 5\% 贴现率下，30 年后的 100 元今天只值 23.14 元。
  \item 300 万房贷 30 年，3.5\% 和 5.5\% 的利息总额差 128.4 万。
  \item 2026 年 8 月：美国 3.50\%–3.75\%、中国 LPR 3.0\%/3.5\%、日本 1.0\%。利差就是资金流动的燃料。
\end{itemize}

\end{takeaway}

\begin{bridgebox}
\textbf{下一章}利率里有一块叫"通胀补偿"。那通胀到底是什么？它凭什么每年从你口袋里拿走一点，还让你以为是自己不会理财？
\end{bridgebox}


\chapter{通货膨胀：每天从你口袋里拿走一点的那个人}
\chaptermeta{04}

\begin{ledetext}
上一章拆利息的时候，"通胀预期"那一块被放过去了。这一章回来算账。
\end{ledetext}

\section{赵敏芝的账本}

赵敏芝 58 岁，沈阳铁西人，退休前在一家轴承厂做了二十六年会计。退休金每月 3870，平时帮女儿带外孙，每周三上午去兴顺街的菜市场。

买菜的小票她还留着，按月夹在一个塑料夹子里。

所以她的账记得比统计局细：上个月排骨 21 块 8 一斤，这个月 24 块 5。

她的结论是，什么都在涨。

同一时期，全国 CPI 同比印出来是零点几。

两边都没撒谎。她们说的不是同一件事。

先把词抠清楚。\term{通货膨胀}{inflation}\index{061@通货膨胀}不是"某样东西涨价"，是\textbf{一般物价水平的持续、普遍上涨}。三个词都得较真：一般是整体不是某一样，持续是一直在涨不是涨完就停，上涨是水平在爬不是水平很高。

排骨那 2 块 7，动的是\term{相对价格}{relative price}\index{069@相对价格}。

相对价格是整个经济体的喊话系统。猪肉贵了，是市场扯着嗓子喊"多养点猪"；手机降价了，是市场在说"这条产线卷完了，人去干别的"。这些喊话指挥着几亿人往哪儿投钱、往哪儿投时间。

通胀不喊话。通胀干的是另一件事：把量身高的那把尺子改短。

\begin{pullquote}{}
相对价格在传递信息，一般物价水平在改变尺子。混成一件事，你会一边骂排骨，一边完全看不见真正在动的东西。
\end{pullquote}

\section{一个不存在的家庭}

CPI 的算法朴素得有点土。统计部门先假想一个"平均家庭"，把这家人一个月买的东西列成清单：粮油肉蛋、房租水电、公交地铁、理发看病、孩子的补习班，每项定个权重。然后每月跑几十万个采价点抄价，加权算总账单变了多少。

麻烦在于，这个"平均家庭"现实里一户都不存在。

\begin{egbox}{先算一笔账}
赵敏芝和老伴一个月开销 6480 块：食品烟酒 2130，居住 1180，医疗保健 940，教育文化娱乐 610，交通通信 520，其他用品及服务 460，生活用品 380，衣着 260。

这个月食品涨 5.8\%，医疗涨 2.4\%，交通通信因为油价回落降 4.1\%，其余没动。

2130×1.058 = 2253.5（多花 123.5），940×1.024 = 962.6（多花 22.6），520×0.959 = 498.7（少花 21.3）。总账单从 6480 变成 \textbf{6604.8，涨 1.93\%}。

而她走出菜市场时脑子里只有一个数：\textbf{5.8\%}。她记住的是权重最大的那次刺痛，统计算的是整张账单。

\end{egbox}

官方数字和体感对不上，有四条原因，每条都成立。

\begin{flowlist}
  \flowitem{01}{每个人的篮子不一样}{赵敏芝每月 940 块花在医药和门诊上，占她开销的 14.5\%。全国口径里医疗保健的权重要低得多。她的通胀率是她那个篮子的通胀率，跟"平均篮子"没有理由相同。她女儿家又是另一副样子：两个孩子，一个上小学一个上幼儿园，课外班加儿科能吃掉家庭开销的三成，而全国口径里教育文化娱乐的权重远没有这么重。同一个城市，同一个月，一家人喊涨，另一家人没感觉。}
  \flowitem{02}{房价不进 CPI}{CPI 量的是消费，买房在统计口径里算投资。进篮子的只有房租，外加一个叫\term{自有住房服务}{owner's equivalent rent}\index{094@自有住房服务}的虚拟项目：假设你把房子租给自己，能收多少。房价翻倍而 CPI 不动，完全可能发生。这不是造假，是定义。可对一个正在攒首付的年轻人，这定义等于把他这辈子最大的一次涨价划掉了。}
  \flowitem{03}{人对涨价敏感，对降价没记性}{她记得住排骨那 2 块 7，记不住四年前 3800 买的手机今天 2100 还更好用。早餐、打车、菜价这种高频小额支出，在记忆里的权重远高于家电和机票。行为经济学管这叫损失厌恶。人话是：疼比爽记得久。}
  \flowitem{04}{质量调整}{同样 4699 块的手机，今年性能是三年前的两倍，统计上要把这部分算成降价，术语叫\term{质量调整}{hedonic adjustment}\index{086@质量调整}。逻辑站得住。可赵敏芝并不觉得自己变富了，她只觉得手机本来就该越来越好。}
\end{flowlist}

"官方数据是编的"这个说法，中国有人讲，美国和欧洲也有人讲，讲了几十年。

我认为它是错的。CPI 的毛病不在造假，在口径，而口径是公开的、可查的、可以被批评的。值得吵的是"该不该把房价放进篮子"，不是"数字是不是编的"。

\section{三条来路}

\begin{tablewrap}
\begin{xltabular}{\linewidth}{>{\raggedright\arraybackslash}p{0.20\linewidth}XX}
\toprule
\bthead{类型} & \bthead{机制} & \bthead{典型画面} \\
\midrule
\textbf{需求拉动} & 钱多货少，总需求跑在产能前面 & 2021 年美国把支票寄进居民信箱，而工厂还没恢复 \\
\textbf{成本推动} & 产量没变，生产成本变了 & 1973 年石油危机；2026 年的关税和中东局势 \\
\textbf{预期自我实现} & 人们相信会涨，提前行动，于是真涨了 & 工人要求涨 5\%，企业提前提价 5\% \\
\bottomrule
\end{xltabular}
\end{tablewrap}

\begin{egbox}{举个栗子}
兴顺街拐角那家馄饨馆，老板姓崔。一碗三鲜馄饨 13 块，成本是：肉馅虾仁和面粉 3.9 元，一次性餐具 0.7 元，摊到每碗的房租 2.6 元，人工 3.4 元，水电燃气 1.1 元。剩 \textbf{1.3 元}归他。

猪肉和虾仁一起涨了 22\%，3.9 变成 4.76。他的利润从 1.3 掉到 \textbf{0.44 元，缩水 66\%}。

崔师傅只有两条路：提到 14 块，或者少赚三分之二。街上七八家早点铺面对同一道题，绝大多数会选前一条。

这里没有任何人多印钱，也没有任何人需求过热。物价照样涨了。

\end{egbox}

三条来路里，央行最怕第三条。前两条是外部条件变了，第三条是人的信念变了，而信念会自我实现。全社会一旦默认"每年都得涨 5\%"，这个数字就会写进工资合同、租约和供应商报价单，然后它真的成立。

央行真正在管的从来不是物价，是关于物价的信念。

信念看不见，但它有价格。

市场上同时挂着两种美国国债：一种付固定票息，一种的本金跟着 CPI 走（\term{通胀保值国债}{TIPS}\index{062@通胀保值国债}）。把两者的收益率一减，剩下那个差额，就是市场此刻押注的未来通胀，行话叫\term{盈亏平衡通胀率}{breakeven inflation}\index{077@盈亏平衡通胀率}。

央行官员每天盯的就是这条线。它一往上飘，说明"关于物价的信念"开始松动了。第 10 章会讲这类债券怎么定价。

\section{为什么不按到零}

既然通胀在偷钱，干脆按到零不行吗。

不行，零是个危险的目标。

\begin{flowlist}
  \flowitem{01}{名义工资是刚性的}{降薪极难。企业宁可裁掉 11\% 的人，也不愿给所有人普降 8\% 的工资，后者能在三天里毁掉一支队伍。而在 2\% 的温和通胀下，给某个岗位"涨 0"，他的\emph{实际}工资就自动降了 2\%，没有人需要开那个会。有点阴险，但省血。}
  \flowitem{02}{得给降息留弹药}{名义利率大致有个下限，负利率试过，副作用不小。长期通胀若是 0\%，正常时期的名义利率也会低 2 个百分点；衰退一来，能降的空间就少 2 个点。}
  \flowitem{03}{CPI 自带向上的测量偏差}{猪肉贵了赵敏芝改买鸡腿，可篮子权重更新有滞后，这叫替代偏差；新产品进篮子也慢；质量调整做不干净。统计出来的通胀通常比真实的生活成本上升高零点几个百分点。把目标定在 0，等于在瞄准负数。}
  \flowitem{04}{离通缩远一点}{目标定在 0，意味着接近一半的时间物价在跌。}
\end{flowlist}

\begin{deepbox}{进阶：可以跳过}
2\% 没有任何神圣性。

它最早被写进法律，是新西兰 1989 年的《储备银行法》。而那个数字的由来带着点偶然：财政部长先在一次电视访谈里随口说了个区间，后来才被磨成正式目标。别的国家看着好使，就抄了过去，一抄三十多年。这段掌故常被引用，细节在不同版本里略有出入。但有一点没争议：它不是从哪个方程里解出来的最优解。

学界一直在吵。一派主张提到 3\% 甚至 4\%：既然名义利率的下限卡在零附近，通胀目标高一点，正常时期的名义利率就高一点，衰退时能降的空间也大一点。另一派主张干脆换成价格水平目标制，某一年通胀低了，第二年要补回来，让物价水平本身回到轨道上，而不是只盯增速。

这件事目前没人给得出定论。

而 2026 年的现实是，主要央行都在极力避免\emph{看起来像}在放松目标。沃什那句"2\% 的通胀目标不会动摇"就是这个用途。市场一旦认为目标可以商量，前面说的盈亏平衡通胀率立刻会往上跑，届时要把它按回去的代价，比当初守住它大得多。

\end{deepbox}

第四条得展开，它最要命。核心机制是欧文·费雪 1933 年提出的\term{债务-通缩螺旋}{debt-deflation spiral}\index{083@债务-通缩螺旋}。

\begin{chainflow}
\chainnode{物价下跌}\chainarrow \chainnode{债务的实际负担上升}\chainarrow \chainnode{企业变卖资产还债}\chainarrow \chainnode{抛售压低价格}\chainarrow \chainnode{物价继续下跌}\chainarrow \chainnode{实际负担又上升}
\end{chainflow}

费雪那句总结很冷：债务人还得越多，欠得越多。他们卖资产还债的动作本身在往下砸价格，而债务的面值纹丝不动。消费者那边同时在踩刹车：反正下个月更便宜，那就等等。

\begin{egbox}{举个栗子}
赵敏芝的女婿胡建军，34 岁，在沈阳一家汽车零部件厂做设备工程师。2021 年买房贷了 94 万，现在还剩 \textbf{87.3 万}，去年到手 21.6 万。债务是他年收入的 4.04 倍。

往后十年他什么都不做，只是被动地活在两种物价环境里。

\textbf{一年 3\% 的通胀，工资同步涨。}十年后到手 21.6×1.03\textsuperscript{10} = \textbf{29.03 万}，87.3 万本金一分没涨，债务收入比降到 \textbf{3.01 倍}。

\textbf{一年 2\% 的通缩，工资也跟着降。}十年后到手 21.6×0.98\textsuperscript{10} = \textbf{17.65 万}，债务收入比升到 \textbf{4.95 倍}。

同一个人，同一笔贷款，什么都没干，负担差了 \textbf{64\%}。通缩不需要任何人犯错，就能让所有债务人同时变穷。

\end{egbox}

日本用三十年演示了这个坑有多深。泡沫 1990 年破掉之后，物价长期趴着，企业忙着还债而不是投资，居民默认明天不会更贵。日本央行 2024 年 3 月才结束负利率与 YCC，2025 年 1 月加到 0.5\%，2025 年 12 月 0.75\%，\textbf{2026 年 6 月 16 日加到 1.0\%}，是 1995 年 9 月以来的最高。

一个发达经济体花了三十年，把利率抬回 1\%。

\section{谁在偷，谁在赚}

\begin{pullquote}{}
通胀不毁灭财富，它只搬运。搬的过程没有声音，被搬的那一方通常最后才知道。
\end{pullquote}

原因在于债务合同都是用名义金额写的。胡建军欠银行 87 万 3，不管十年后一碗馄饨卖到多少钱，他还的都是 87 万 3。尺子变短，数字没变，负担自动轻了。

\begin{tablewrap}
\begin{xltabular}{\linewidth}{>{\raggedright\arraybackslash}p{0.20\linewidth}Xr}
\toprule
\bthead{谁} & \bthead{通胀来了会怎样} & \bthead{方向} \\
\midrule
\textbf{胡建军这样的房贷借款人} & 债务名义固定，收入和房价名义上涨 & 受益 \\
\textbf{赵敏芝这样把钱放定期的人} & 数字不动，购买力一年一年往下掉 & 受损 \\
\textbf{持有实物资产和股权的人} & 长期跟着名义变量走，但冲击刚发生时股债经常一起跌 & 多数受益 \\
\textbf{靠固定名义收入过活的人} & 固定票息的债券持有人、签了长约的房东、调整滞后的退休金 & 受损最重 \\
\textbf{政府} & 债务名义固定，税基跟着名义 GDP 一起膨胀 & 通常是最大赢家 \\
\bottomrule
\end{xltabular}
\end{tablewrap}

最后那一行有个名字，叫\term{没有立法的税}{taxation without legislation}\index{036@没有立法的税}。不经过任何一次投票，转移就完成了。

\begin{egbox}{举个栗子}
赵敏芝去年把铁西那套老房子卖了 96 万，全存了三年定期。她的说法是"等看清楚再说"。

按 3\% 的通胀算，十年后这笔钱能买到的东西，相当于今天的 96 ÷ 1.3439 = \textbf{71.4 万}。按 5\% 算，是 \textbf{58.9 万}。

"先放着"本身也是一个决定。它把风险从价格波动换成了购买力的慢性流失，没有把风险取消掉。

\end{egbox}

\begin{warnbox}{小心}
上面这段算术极容易被拿去当推销词。"现金会贬值，所以你必须马上买某某"是几乎所有话术的标准开场白。\hlwarm{购买力缩水是事实，"因此某项资产一定跑赢通胀"是推论，而且经常是错的推论。}本书不提供任何投资建议，第 21 章和第 23 章会专门处理。

\end{warnbox}

\section{两场相反的仗}

\begin{statrow}{3}
  \statitem{3.50 至 3.75}{\%}{美国联邦基金利率目标区间，2026 年 7 月，连续第 5 次维持不变}
  \statitem{5.02}{\%}{日本 2026 年春斗薪资涨幅预期}
  \statitem{6.3}{\%}{日本 5 月 PPI 同比}
\end{statrow}

\textbf{美国的通胀还在，而且很粘。}截至 2026 年 8 月，压力主要来自\hlmark{供给侧}：关税抬高进口成本，中东局势推高能源价格，去全球化带来重构成本，AI 资本开支拉动电力和设备需求。2025 年全球数据中心资本开支同比增长约 57\%，它们要的电和变压器都是真实且稀缺的东西。FOMC 的表述是"通胀仍然高企，部分受冲击影响"，市场上一度讨论过加息。

\begin{talkbox}
  \talkline{新闻}{美联储主席凯文·沃什 2026 年 7 月表示："通胀问题不可能在九周内解决。""2\% 的通胀目标不会动摇。"}
  \talkline{翻译}{第一句在说：这轮通胀的根子在供给侧，利率工具变不出更多的石油、芯片和电力，别指望一两次会议就有结果。第二句在说：目标短期够不着，但我绝不暗示它可以商量。一旦松口，预期那道堤坝就开始漏水，到时候要付的代价比现在大得多。}
\end{talkbox}

\textbf{日本那边，久违的螺旋居然是好消息。}2026 年春斗薪资涨幅预期约 5.02\%，5 月 PPI 同比 6.3\%。工资推物价、物价再推工资，这在别的国家是刺耳的警报，在日本几乎算一句祝福：钉了三十年的通缩心态终于松动。日本央行的任务因此换了一个，确认苗头不会变成失控。

\textbf{中国的问题是反的。}这边不是钱太多，是\textbf{需求不足、物价低迷}。货币政策基调"适度宽松"，7 天期逆回购政策利率 1.4\%，2026 年 8 月 20 日的 1 年期 LPR 3.0\%、5 年期以上 3.5\%，连续多月持平。商业银行净息差二季度末降到 1.41\% 的历史低位，这个数字本身就说明降息空间被挤得很紧。政策重心因此从"放水"挪到"让人愿意花钱"：商务部等 7 部门提出到 2030 年社零达 60 万亿元左右，8 月底又出了一套房地产组合拳。

同一个月，华盛顿在想怎么把物价按下去，东京在庆幸物价终于起来了，北京在想怎么让人愿意花钱。

教科书里那个"通胀"是单数。

\section{扔掉两个直觉}

\begin{mythbox}{常见误解}
\mvfrow{误解}{"印钱必然导致通胀。央行一放水，物价就一定涨，这是常识。"}
\mvfrow{事实}{2009 年到 2019 年，美联储、欧洲央行、日本央行做了人类历史上规模最大的\term{量化宽松}{QE}\index{033@量化宽松}，央行资产负债表扩张了好几倍，基础货币暴增。\textbf{这十年发达国家的 CPI 通胀始终低于 2\% 的目标}，欧元区和日本还长期在零附近挣扎。当年预言"恶性通胀马上来"的人排着队，等了十年也没等到。}
2021 到 2023 那一轮真正的通胀，主因是疫情造成的供应链中断和产能损失、财政把钱直接转到居民账户、需求结构从服务业剧烈切换到实物商品，以及后来的能源冲击。

关键差别在这里。

QE 不是"印钱发给老百姓"，是央行用新创造的准备金从金融机构手里买走债券。银行的资产结构变了，债券换成准备金，可你我账户里的存款没有因此多出一笔。所以 QE 推高的是资产价格，不是货架上的价格。2020 年那些寄到家里的支票是另一件事。

但别滑到另一个极端。\textbf{这不代表货币无关紧要。}

战时财政货币化、恶性通胀这类极端情形里，货币是决定性的。准确的说法是：简单的货币数量论，M 涨多少 P 就涨多少，\hlmark{已经不能用来做短期预测}。现代央行早就不拿货币总量当中介目标，改盯利率和预期。

\end{mythbox}

\begin{mythbox}{常见误解}
\mvfrow{误解}{"加了关税，进口的东西贵了，这就是通胀。"}
\mvfrow{事实}{关税制造的是\textbf{价格水平的一次性跳升}，不是通胀率的持续上升。这两件事在 2026 年格外需要分清楚。}
假设关税让某类商品一次性贵了 10\%。生效之后的 12 个月里，同比 CPI 确实会被这 10\% 抬高。可到了第 13 个月，只要关税不再加码，这个效应\textbf{自动消失}，因为做同比的基数里已经含着它了。价格停在更高的台阶上，不再往上爬。

它变成持续通胀只有一条路：\hlmark{这次涨价改变了人们的预期}，工人据此要求加薪，企业据此提前调价，第二年第三年接着走。

这就是 2026 年美联储最难受的地方：加息变不出更多的石油、芯片和电力，它能做的只有把预期按住，代价是牺牲需求和就业。沃什那句"不可能在九周内解决"不是拖延，是诚实。

\end{mythbox}

\bookbreak

回到那条主线。钱是记账，记账要有人信。通胀就是这套记账系统的计量单位在悄悄变小：它不改动任何一行数字，只改动每一行数字的意思。而债务合同按名义金额写死，单位一变，账就自己动了。

赵敏芝那 96 万还在定期里。她还在等看清楚。

\begin{takeaway}{本章一句话}
通胀不是"某样东西涨价"，是记账单位本身在缩水；它不毁灭财富，只在债权人和债务人之间做一次没人签字的再分配。

\begin{itemize}
  \item 官方通胀和体感通胀对不上，是因为篮子不同、房价不计入、人对涨价更敏感、质量调整。四条都成立，不需要阴谋论。
  \item 2\% 是安全边距不是神圣数字：给工资调整留润滑，给降息留弹药，抵消测量偏差，离通缩远一点。
  \item 通缩比通胀难治。它靠债务实际负担的上升自我强化，不需要任何人犯错。
  \item 2026 年的全球通胀主要来自供给侧，而货币政策治不了供给，只能锚住预期。
\end{itemize}

\end{takeaway}

\begin{bridgebox}
\textbf{下一章}通胀在债权人和债务人之间搬钱。可"债权"和"债务"到底是怎么记下来的？下一章给你一副眼镜，戴上它，从赵敏芝的定期到国家的财政赤字都会变得半透明。
\end{bridgebox}


\chapter{资产、负债、净值：全世界通用的三个词}
\chaptermeta{05}

\begin{ledetext}
上一章说通胀在债权人和债务人之间搬钱。想看清怎么搬，得先看清账是怎么记的。这一章给的不是知识，是一件工具。
\end{ledetext}

\section{王雪梅的三个数字}

王雪梅 31 岁，合肥人。2025 年秋天教培机构裁员，她那个课程顾问的岗位整组没了。丈夫周立在一家汽车电子公司做测试工程师。

失业第三个月的一个晚上，她把银行 App、公积金页面和两张还款计划表一起摊在餐桌上，第一次把家里的账算了一遍。

算法只用到三个词。

\term{资产}{assets}\index{093@资产}，你拥有的、以后能给你带来好处的东西：房子、存款、基金、一台还能干活的机器，还有别人欠你的钱。\term{负债}{liabilities}\index{017@负债}，你欠别人的、以后必须交出去的东西：房贷、车贷、信用卡账单、公司欠供应商的货款。\term{净值}{net worth}\index{030@净值}，两者之差。

\begin{pullquote}{}
资产 = 负债 + 净值
\end{pullquote}

它永远成立，是因为净值不是量出来的，是\textbf{被定义成残值}的：把你有的数清楚，把你欠的减掉，剩多少算多少。没有作弊的空间，也没有解释的余地。

\begin{bsheet}
  \bsside{资产 ASSETS}{%
    \bsrow{滨湖的房子（同户型挂牌价）}{232 万}
    \bsrow{存款}{18.6 万}
    \bsrow{两人的公积金账户}{9.4 万}
    \bsrow{车（二手估价）}{7.3 万}
    \bsrow{基金}{4.1 万}
    \bsrowtotal{合计}{271.4 万}
  }
  \bsside{负债 + 净值}{%
    \bsrow{房贷余额}{176.4 万}
    \bsrow{车贷余额}{2.8 万}
    \bsrow{信用卡待还}{1.9 万}
    \bsrow{净值}{90.3 万}
    \bsrowtotal{合计}{271.4 万}
  }
\end{bsheet}

这套房 2021 年买进来的时候是 268 万。

光是挂牌价从 268 万滑到 232 万这一件事，就从这张表的净值里拿走了 36 万。而房贷余额那一栏，一分钱都没跟着少。

这个不对称本章结尾还要再算一次。

左右两栏必须严丝合缝地相等。不等，就是有东西没数清楚。

公司的表长得一模一样，只是"净值"那一栏换了个名字，叫\term{所有者权益}{equity}\index{057@所有者权益}。同一个位置，同一套算法：资产数清楚，欠人的减掉，剩下的归股东。第 6 章讲"股"这张合同的时候，讲的就是这一栏。

\section{五百年前的记账法}

这套东西叫\term{复式记账}{double-entry bookkeeping}\index{019@复式记账}。规则只有一句：任何一笔交易都会同时改动至少两个科目，改完之后等式照样成立。

1494 年，方济各会修士卢卡·帕乔利在威尼斯出版了一本数学书，顺手把当地商人已经在用的这套办法写成了系统。五百多年过去，从王雪梅手机上那行流水到美联储的资产负债表，用的还是同一套语法。

它的本事在于，它逼你回答"另一边发生了什么"。

\begin{egbox}{举个栗子}
王雪梅想拿 11 万提前还房贷，理由是"这样心里踏实"。查账：

资产端存款 18.6 万变成 \textbf{7.6 万}，负债端房贷 176.4 万变成 \textbf{165.4 万}。

新合计：资产 260.4 万 = 负债 170.1 万 + 净值 \textbf{90.3 万}。

\textbf{净值一分没动。}

提前还贷让表的两边同时缩小，它改的是往后的利息支出和风险敞口，不是财富。

对照一下：周立年底要是拿到 4.2 万奖金并存下来，资产加 4.2 万，负债不动，净值从 90.3 万变 \textbf{94.5 万}。这才叫变富。

"心里踏实"是个正当理由。只是它买的不是财富，是睡眠。

\end{egbox}

\section{全书最重要的一句话}

\begin{pullquote}{}
一个人的金融资产，一定是另一个人的负债。
\end{pullquote}

王雪梅手机里那 18.6 万余额，她以为是"她的钱"。在银行的账本上，它躺在\textbf{负债}那一栏，是银行欠她的。银行拿什么撑住这笔负债？资产栏里的贷款，也就是别人欠银行的。

\begin{chainflow}
\chainnode{王雪梅的 18.6 万存款}\chainarrow \chainnode{是银行的负债}\chainarrow \chainnode{银行的资产是贷款}\chainarrow \chainnode{那是别人的负债}\chainarrow \chainnode{别人的资产是房子和厂房}\chainarrow \chainnode{实物资产，链条到头}
\end{chainflow}

链条可以拉得很长：她的存款、银行的债券、企业的应付账款、政府的国债。但它必须终止在两处之一，要么是实物资产，要么是股权。股权是对剩余收益的索取权，不是一张固定金额的欠条。

换句话说，金融本身不生产任何东西。它只把实物世界的收益重新切开，分给不同的人和不同的时间。第 6 章会把这句话变成整本书的支点。

\begin{databox}{2026 数据}
国际金融协会（IIF）的统计：截至 \textbf{2026 年 3 月底}，全球债务总额约 \textbf{353 万亿美元}，创历史新高，相当于全球 GDP 的约 \textbf{305\%}。同期，美国国债规模在 2026 年 8 月\textbf{突破 40 万亿美元}。

标题很吓人。把纸翻过来看背面：这 353 万亿同时是另一些人的 \textbf{353 万亿资产}。养老金账户里的国债，保险公司的准备金，银行资产栏里的贷款，王雪梅那 18.6 万。

所以债务的\emph{总量}本身说明不了危险。危险在结构里：谁欠的、欠谁、什么币种、什么期限、拿什么现金流还。第 14 章和第 17 章拆这个。

\end{databox}

\section{谁更有钱}

\begin{egbox}{举个栗子}
\textbf{李振宇}，34 岁，周立的同事，做项目经理，月薪 3 万 2，到手两万四出头。

资产：高新区的房子现在挂 415 万、车 22 万、存款 3.2 万，合计 \textbf{440.2 万}。负债：房贷 411 万、车贷 15.6 万、装修贷 15.4 万、信用卡分期 6.6 万，合计 \textbf{448.6 万}。净值 \textbf{负 8.4 万}。

\textbf{孙桂芳}，41 岁，在蚌埠一家医院做护士，月薪 1 万 4，到手 1 万 1。

资产：没有贷款的房子 63 万、存款 14.8 万、指数基金 6.2 万、公积金 4.7 万，合计 \textbf{88.7 万}。负债 \textbf{0}。净值 \textbf{88.7 万}。

饭桌上所有人都会说李振宇有钱。

真正分胜负的是断粮那天。李振宇每月刚性支出是房贷 2 万 1 千 8、车贷 4300、装修贷 3800、信用卡 2200、基本生活 6000，加起来 \textbf{3 万 8 千 1}。账上 3.2 万，撑 \textbf{25 天}。孙桂芳每月刚性支出 3400，可动用的金融资产 21 万，撑 \textbf{61 个月}。

\end{egbox}

收入是流量，净值是存量。流量决定你现在过得怎么样，存量决定你扛得住什么。失业、生病、行业塌方，考的全是存量。

\section{一家奶茶店的三个月}

个人有一张表就够。公司要三张。

\begin{tablewrap}
\begin{xltabular}{\linewidth}{>{\raggedright\arraybackslash}p{0.20\linewidth}XX}
\toprule
\thead{报表} & \thead{它是什么} & \thead{回答的问题} \\
\midrule
\textbf{资产负债表} & 某一瞬间的\textbf{照片} & 这一刻我有什么、欠什么、剩什么 \\
\textbf{利润表} & 一段时间的\textbf{录像} & 这三个月我到底赚了还是赔了 \\
\textbf{现金流量表} & 真金白银的\textbf{进出账} & 这三个月钱包里的钱是多了还是少了 \\
\bottomrule
\end{xltabular}
\end{tablewrap}

2026 年 3 月，王雪梅在政务区一个写字楼底商开了家奶茶店，招牌两个字是她自己在 iPad 上描的。自己拿 22.6 万，向她妈借 18 万，说好三年后还，年息 3\%。

借条是她妈手写的，压在家里的玻璃板底下。老太太很讲原则。

这 40.6 万花掉：装修 14.8 万、设备 9.6 万、押一付三的房租 5.4 万、首批原料 3.2 万，剩 7.6 万压箱底。

\begin{bsheet}
  \bsside{开业当天 · 资产}{%
    \bsrow{装修（待摊）}{14.8 万}
    \bsrow{设备}{9.6 万}
    \bsrow{预付租金}{5.4 万}
    \bsrow{存货}{3.2 万}
    \bsrow{现金}{7.6 万}
    \bsrowtotal{合计}{40.6 万}
  }
  \bsside{负债 + 净值}{%
    \bsrow{欠她妈}{18 万}
    \bsrow{净值（王雪梅的本钱）}{22.6 万}
    \bsrowtotal{合计}{40.6 万}
  }
\end{bsheet}

4 到 6 月生意不错。卖了 2.87 万杯，均价 14.5 元，营业额 41.6 万。其中 4.8 万是给两家公司做团建供应的挂账，说好 9 月结账。这在会计上叫\term{应收账款}{accounts receivable}\index{076@应收账款}，是"赚了但还没收到"的钱。实际进账现金 36.8 万。

\begin{egbox}{举个栗子}
\textbf{利润表（4 到 6 月，万元）}：营业收入 41.6，扣掉原料 12.4、人工 9.3、房租 5.4、水电杂项 2.35、折旧摊销 1.71、利息 0.135，剩下 \textbf{利润 10.31}。

\textbf{现金流量表（4 到 6 月，万元）}：收到现金 36.8，付出原料 10.8（还欠供应商 1.6）、人工 9.3、水电 2.35、利息 0.135，房租这季度实付 0（开业前就付掉了），\textbf{现金净增 14.22}。期末现金 7.6 加 14.22，等于 \textbf{21.82}。

同一家店，同样三个月：利润 10.31 万，现金多了 14.22 万。差的 3.91 万从哪来？

房租 5.4，上季度就付了现金，这季度才计成本。折旧 1.71，根本没花钱。欠供应商的 1.6，还没付。三项加起来再减去挂账没收到的 4.8，正好 \textbf{3.91}。

分毫不差。

\end{egbox}

第二幕在 7 月。要续交下季度房租 5.4 万，还供应商 1.6 万，暑期旺季还得压 5.8 万的货。账上 21.82 万，撑得住。

可要是她这时候接了个 28.5 万的挂账大单，为此提前备 11.4 万的料，利润表会更好看，而现金在往外流。再接两单这样的"好生意"，她就得去借经营贷发工资了。

利润是一种观点，现金是一个事实。

利润取决于你怎么确认收入、怎么摊销、折旧年限估几年，里面装着一堆判断。账上还剩多少钱，没得商量。

\hlmark{很多公司死的时候，利润表还是正的。}

折旧那 1.71 万里，藏着整套会计最反直觉的一个概念。

王雪梅那台 9 万 6 的设备能用 5 年。她不能在买它的那个月把 9 万 6 全记成成本，那会让当月巨亏、后面四年多虚增利润，两头都不真实。正确做法是每年"消耗"1 万 9 千 2。这笔钱没有一分从账户流出去，可它确实是成本：设备真的在变旧，五年后她真的还得再掏一次钱。

\term{折旧}{depreciation}\index{084@折旧}就是把一次性的现金流出，摊到它真正被消耗的那些年里。

关键在于，折旧年限是一个判断，不是一个事实。

\begin{egbox}{举个栗子}
一家公司买了 137 亿的服务器。按 3 年折旧，每年成本 \textbf{45.7 亿}；按 6 年折旧，每年 \textbf{22.8 亿}。

同样的机器、同样的电费、同样的收入，账面利润差 \textbf{22.9 亿}，现金流一分钱没变。

这不是学术问题。2025 年全球数据中心资本开支同比增长约 57\%，规模越大，"用几年"这个假设的杠杆越大。\hlwarm{如果服务器实际 3 年就被新一代芯片挤下去，却按 6 年摊，那些省下来的成本迟早要一次性认回来。}AI 服务器该按几年折旧，是 2026 年最热的一场会计争论，目前没有定论。第 24 章会回到这里。

\end{egbox}

\section{三个口袋}

\begin{deepbox}{进阶：可以跳过}
把复式记账的逻辑放大到整个国家，会得到一个让很多人不舒服的结论。

把一个经济体切成三个口袋：政府、私人部门（居民加企业）、对外部门。因为\textbf{每一笔支出都是另一个人的收入}，三个口袋的收支之和必然为零。整理一下就是：

\textbf{政府赤字 = 私人部门盈余 + 对外部门赤字}。写成符号是 G-T =（S-I）+（M-X）。

\hlmark{这是会计恒等式，不是政策主张。}它只陈述一个算术事实：三个口袋不可能同时净盈余。政府压缩赤字，在其他条件不变时，要么私人部门积累金融资产的速度变慢，要么贸易顺差得变大，二者必居其一。

所以"政府必须像家庭一样量入为出"这句话，在会计上是错的。一个家庭省钱会变富，可\textbf{所有人同时省钱会让所有人的收入一起下降}，因为你的支出就是我的收入。经济学管这叫节俭悖论。

\textbf{这绝不意味着政府可以无限借债。}真正的约束不是"会破产"。一个只用本币举债、又拥有自己央行的政府，名义违约是政策选择而非技术必然。真正的约束有四条：\\ ① \textbf{通胀}。政府花钱争夺的是钢筋、工人和电力这些真实资源，资源不够时，多出来的钱只会变成物价。\\ ② \textbf{外部平衡}。赤字带动进口和资本外流，汇率承压。\\ ③ \textbf{币种和期限结构}。用外币借的债，你的央行印不出来；短期债要不断滚动，市场一旦不买单，利率瞬间跳升。历史上真正的主权债务危机几乎全出在这一条。\\ ④ \textbf{利息负担}。付息会挤占其他所有支出。

2026 年 8 月，美国国债突破 40 万亿美元，10 年期国债收益率升至 2025 年 1 月以来最高。市场在给期限风险和通胀风险重新定价，不是在给破产风险定价。围绕这套逻辑（MMT 是其中最激进的版本）学界吵得很凶，本书不站队，只请你分清哪一部分是会计，哪一部分是政策判断。

\end{deepbox}

\section{两副该摘的眼镜}

\begin{mythbox}{常见误解}
\mvfrow{误解}{"我有一套两百多万的房子，所以我有两百多万。"}
\mvfrow{事实}{你有的是\textbf{房产估值减去房贷余额}。而这两个数字的性质完全不同：\textbf{估值有弹性，负债没有弹性。}}
王雪梅家挂牌 232 万，房贷余额 176.4 万，账面净值 55.6 万。真卖一次看看：谈到 214 万成交，中介费按 1.5\% 算 3.21 万，各项税费 2.4 万，实收 208.39 万，还完贷款剩 \textbf{31.99 万}。

账面 55.6 万，落袋 32 万，缩水 \textbf{42.4\%}。这还没算从挂牌到过户的几个月。急着用钱的时候，时间本身就是最贵的成本。这叫\term{变现摩擦}{liquidity friction}\index{006@变现摩擦}。

再看负债那栏：176.4 万是精确到分的、白纸黑字的、每月必须付的。它不会因为行情不好就体谅你。\textbf{估值可以商量，还款日不能商量。}几乎所有个人财务危机都长在这个不对称上。这不是说买房错了，是说别把账面净值当成能动用的钱。第 19 章会教你做一张诚实的表。

\end{mythbox}

\begin{mythbox}{常见误解}
\mvfrow{误解}{"国家欠这么多债，跟一个家庭欠了一屁股债一样，早晚要还不起。"}
\mvfrow{事实}{这个类比在四个地方都不成立。但它的反面同样错：政府债务确实有硬约束，只是约束不叫"破产"。}
\textbf{寿命不同。}家庭必须在退休前还清，国家没有生命周期，只要增长率不长期低于利率，债务可以一直滚下去。\\ \textbf{支出的性质不同。}你花的钱不是你自己的收入，政府花的钱恰恰是私人部门的收入。\\ \textbf{定价权不同。}你没法影响自己的贷款利率，拥有独立央行的政府对本国利率有相当大的影响力。\\ \textbf{债主不同。}你的债主是外人，一国国债的相当一部分是欠本国居民、本国养老金和本国银行的，那既算负债，也算国民的资产。

\textbf{以上四条都不构成"可以随便借"的理由。}天花板是通胀、汇率和币种结构。历史上出事的国家，多数不是因为债务占 GDP 的比例太高，是因为借了自己印不出来的货币。

\end{mythbox}

\bookbreak

回到主线。钱是记账，而这一章给的是读账的语法。往后无论谈银行、央行、债券、股票、杠杆还是危机，都只是在看几张资产负债表怎么互相咬合、怎么一起膨胀、又怎么在某一天同时收缩。

王雪梅那张表上，90.3 万的净值里有 55.6 万锁在一套挂了四个月没卖动的房子里。

\begin{takeaway}{本章一句话}
资产 = 负债 + 净值。这个恒等式管你、管你楼下的奶茶店、管你的银行、也管你的国家；在这套语法里，一个人的金融资产必然是另一个人的负债。

\begin{itemize}
  \item 净值才是财富，收入只是流量。提前还贷不改变净值，涨薪存下来才改变净值。
  \item 利润是一种观点，现金是一个事实。折旧、应收账款和预付租金会让两者长期背离。
  \item 部门恒等式是会计不是政策。政府债务的真正约束是通胀、外部平衡和币种结构。
  \item 估值有弹性，负债没有弹性。所有财务危机都长在这个不对称上。
\end{itemize}

\end{takeaway}

\begin{bridgebox}
\textbf{下一章}有了这副眼镜，可以问那个最根本的问题了：把所有花哨名词剥干净之后，金融这门生意到底在交易什么？
\end{bridgebox}


\input{chapters/ch-06}

\part{管道篇：这台机器怎么转}
\chapter{货币是怎么"凭空"出现的}
\chaptermeta{07}

\begin{ledetext}
第 2 章说过，你的存款不在银行的金库里。这一章要说的更难听一点：那笔钱在被借给你之前，压根不存在。
\end{ledetext}

\section{2024 年 6 月的一个下午，300 万从空气里长出来}

陈宁 34 岁，在杭州一家做工业相机的公司当结构工程师，月薪 2 万 3，到手 1 万 8 出头。2024 年 6 月他在余杭买了套 486 万的二手房，卖家是个想换大三居的中学老师，姓项。首付 186 万，剩下 300 万找银行贷。

贷款金额之所以正好是 300 万这么齐整的数，不是他挑的，是评估价乘以贷款成数的上限卡在那儿，他顶格贷满了。

那天下午三点多，客户经理发来一条微信：过了。

把这一秒放慢。批贷那一刻，银行到底干了什么？

多数人脑子里是这幅画面：银行有个大池子，装着王阿姨、你妈、隔壁老陈存进去的钱。信贷员拿个瓢舀出 300 万倒进陈宁账户，水位下降 300 万。

这幅画很直观，很符合常识，而且\textbf{是错的}。

真实发生的事情是，一位系统操作员在核心账务系统里敲下了两行记录。

\begin{bsheet}
  \bsside{放贷前 · 某支行（万元）}{%
    \bsrow{准备金}{47}
    \bsrow{持有的债券}{63}
    \bsrowtotal{资产合计}{110}
  }
  \bsside{负债 + 净值}{%
    \bsrow{客户存款}{98}
    \bsrow{资本（净值）}{12}
    \bsrowtotal{合计}{110}
  }
\end{bsheet}

\begin{bsheet}
  \bsside{放贷后 · 同一家支行（万元）}{%
    \bsrow{准备金}{47}
    \bsrow{持有的债券}{63}
    \bsrow{对陈宁的贷款}{300}
    \bsrowtotal{资产合计}{410}
  }
  \bsside{负债 + 净值}{%
    \bsrow{客户存款}{98}
    \bsrow{陈宁的活期存款}{300}
    \bsrow{资本（净值）}{12}
    \bsrowtotal{合计}{410}
  }
\end{bsheet}

资产端多了一笔"对陈宁的贷款 300 万"，负债端同时多了一笔"陈宁的存款 300 万"。两行一起出现，账自动就平了。

王阿姨的余额那天晚上和前一天一模一样。没有任何一个储户的账户被动过一分钱。

而全社会的钱，那一秒多了 300 万。

\begin{pullquote}{}
银行不是把别人的钱借给你。银行是在给你贷款的同时，顺手把你的存款也造了出来。
\end{pullquote}

回到第 2 章的结论：你的存款说白了就是银行欠你的一张欠条。所以"银行创造货币"这件听着像魔法的事，拆开看很朴素，无非是给自己记了一笔"我欠陈宁 300 万"的账。

你我也能记。我现在就可以撕张纸写"欠陈宁 300 万"。

区别只有一条，但这一条是天堑：\textbf{银行写的欠条，全社会都当钱花。}陈宁拿它付给项老师，项老师拿它交契税、付中介费、给女儿报暑期班，没有一个人问过"这真的能兑现吗"。

那个没人问的瞬间，就是整台机器的地基。

故事往下走。项老师的卡不在这家银行。

跨行付款要在央行的账本上清算，这家支行必须真金白银划出 300 万\term{准备金}{reserves}\index{089@准备金}，也就是它存在央行账户里的那笔钱。可它账上只有 47 万。

去借。

\begin{egbox}{举个栗子：钱是先放出去，还是先凑够}
这家支行的操作顺序是：先在银行间市场拆入 253 万（准备金凑到 300 万，负债端多一笔"同业拆入 253 万"），然后把 300 万准备金划给项老师那家银行，陈宁的存款同时清零。

请盯住这个\textbf{时间顺序}。贷款早就放出去了，准备金是\hlmark{事后才去凑的}。银行不是先有准备金才敢放贷，是先放贷再补窟窿。窟窿补得贵不贵，看银行间市场当天的价，而那个价是央行说了算的。

\end{egbox}

\begin{bsheet}
  \bsside{陈宁跨行付款后（万元）}{%
    \bsrow{准备金}{0}
    \bsrow{持有的债券}{63}
    \bsrow{对陈宁的贷款}{300}
    \bsrowtotal{资产合计}{363}
  }
  \bsside{负债 + 净值}{%
    \bsrow{客户存款}{98}
    \bsrow{同业拆入}{253}
    \bsrow{资本（净值）}{12}
    \bsrowtotal{合计}{363}
  }
\end{bsheet}

贷款还在资产端，负债端却从"欠陈宁"变成了"欠同业"。

货币总量一分没少。那 300 万此刻躺在项老师的银行里，只是换了个门牌号。钱在银行之间流动，改变的是它记在谁名下，不是它存不存在。

能让它消失的只有一件事。

\section{陈宁每还一笔本金，就杀死一笔钱}

把时间快进到 2029 年。陈宁的公司被并购，他手上一笔期权兑现，决定把房贷一次性结清。那时候本金还剩 269 万。

银行的操作是：把他账上的 269 万存款划掉，同时把资产端"对陈宁的贷款"抹平。

\begin{bsheet}
  \bsside{结清之后（万元）}{%
    \bsrow{准备金}{47}
    \bsrow{持有的债券}{63}
    \bsrow{对陈宁的贷款}{0}
    \bsrowtotal{资产合计}{110}
  }
  \bsside{负债 + 净值}{%
    \bsrow{客户存款}{98}
    \bsrow{陈宁的存款}{0}
    \bsrow{资本（净值）}{12}
    \bsrowtotal{合计}{110}
  }
\end{bsheet}

那 269 万去哪儿了？

哪儿也没去。它不存在了。

它没进银行口袋。银行的资本还是 12 万，一分没多。这笔货币被注销了，像一条被抹掉的数据库记录。

陈宁每月还的那笔钱也一样。利息部分变成银行收入，会重新流回经济；\textbf{本金部分是被销毁的}，贷款和存款同步减少，社会上的钱真的少了那么一点。

\begin{databox}{这解释了正在发生的一件事}
几百万个陈宁同时决定"把房贷提前还了吧"，会怎么样？

在每一张家庭账本上这都是明智事：少欠钱，少付息，晚上睡得着。几百万张账本叠起来，它换了个名字，叫\textbf{货币紧缩}。存量贷款收缩，对应的存款被同步注销，社会总货币量往下走。

于是有了那个别扭的组合：央行说政策"适度宽松"，7 天期逆回购利率稳在 1.4\%（截至 2026 年 8 月），身边的人却普遍觉得钱很紧。央行能定水价，定不了水量。水量是千千万万个陈宁在"借"和"还"之间投票投出来的。

\end{databox}

\section{教科书里那个乘数，因果讲反了}

上过经济学原理课的话，脑子里大概有这么一套流程。

\begin{chainflow}
\chainnode{央行投放准备金}\chainarrow \chainnode{银行留下 10\% 法定准备金}\chainarrow \chainnode{把 90\% 贷出去}\chainarrow \chainnode{钱又存回银行}\chainarrow \chainnode{再贷出 81\%}\chainarrow \chainnode{最终放大 10 倍}
\end{chainflow}

这套东西叫\term{货币乘数}{money multiplier}\index{024@货币乘数}。它优雅、好考、能算出漂亮的等比数列。

它是错的。

不是"不够精确"，是因果方向整个反过来了。教科书说先有准备金才有贷款，现实是先有贷款再补准备金。

这不是什么非主流学派的私货。2014 年，英格兰银行在季度公报上发了一篇官方论文《现代经济中的货币创造》，白纸黑字写着：银行放贷并不是把储户的存款转手借出，而是通过放贷创造存款；并且明确指出，货币乘数式的叙述\hlmark{不是对货币实际创造过程的准确描述}。欧洲央行后来的官方文章、美联储研究人员的表述，口径完全一致。

十二年过去了。教科书还在照原样印。

\begin{mythbox}{常见误解}
\mvfrow{误解}{"存款准备金率决定货币总量。降准 0.5 个百分点，社会上的钱就成倍增加。"}
\mvfrow{事实}{准备金率今天是\textbf{成本和流动性的旋钮}，不是货币总量的乘法器。降准给银行释放长期低成本资金，负债成本下来一点，它\emph{可能}更愿意放贷。但如果没人想借，或者银行资本不够、风险偏好塌了，释放的钱就趴着当超额准备金，什么都不会发生。}
\end{mythbox}

\begin{mythbox}{常见误解}
\mvfrow{误解}{"银行是资金的搬运工，把储户的钱搬给借款人，中间赚个差价。"}
\mvfrow{事实}{银行是\textbf{货币的制造商}。它的生意是用自己的负债（存款）买别人的负债（贷款合同），并承担期限错配和信用风险。它之所以是被管得最严的行业，就因为它手里那支笔写出来的东西，全社会都当钱用。}
\end{mythbox}

\section{那什么在拦着银行}

既然敲两行字就能造钱，银行为什么不把全世界的钱都造出来，给自己发年终奖？

约束是真实存在的，只是不在你以为的地方。

\begin{flowlist}
  \flowitem{01}{有没有值得贷的人}{贷款合同要两个签名。银行想放，得有人敢借，还得让银行信他还得起。经济下行时最缺的从来不是钱，是\hlmark{敢签字的借款人}。}
  \flowitem{02}{资本充足率}{巴塞尔 III 的核心规则：每放一笔贷款，都要按风险权重占用一块自有资本。资本是股东真金白银掏的，不能凭空写两行。\textbf{这是最硬的一道闸。}}
  \flowitem{03}{赚不赚钱}{截至 2026 年二季度末，中国商业银行\term{净息差}{net interest margin}\index{029@净息差}只有 1.41\%，历史低位。放 100 万贷款一年毛赚 1.41 万，扣掉运营成本和坏账拨备，剩下的还未必比拆借给同业强。\hlmark{放贷不赚钱，机器就不想开。}}
  \flowitem{04}{监管在盯着}{房地产贷款集中度、行业限额、压力测试、逆周期资本缓冲。监管随时可以直接把某个方向的龙头拧小，不需要经过市场同意。}
  \flowitem{05}{央行定的那个价}{事后要补的准备金是有价格的。把这个价抬高，放贷成本上升，能过审的项目变少。\textbf{这才是央行真正的抓手。}}
\end{flowlist}

\begin{egbox}{举个栗子：资本这道闸有多硬}
某城商行核心一级资本 386 亿元。个人住房抵押贷款的风险权重按抵押率分档，这里取一个便于计算的档位 50\%，监管要求的核心一级资本充足率（含储备资本）按 7.5\% 算。

它能撑起的风险加权资产上限是 386 ÷ 7.5\% ≈ \textbf{5147 亿元}，按 50\% 权重折回去，房贷规模上限大约 \textbf{1.03 万亿元}。

落到陈宁头上：每放一笔 300 万的房贷，占用风险加权资产 150 万，占用核心一级资本 \textbf{11.25 万元}。这家银行最多放 34 万个陈宁，然后就得停。哪怕它趴在一座超额准备金的山上，也一笔都不能多放。

\end{egbox}

\section{三层货币：你以为的"钱"其实分三种电压}

现在装上这一章最有用的一副眼镜。

\begin{tablewrap}
\begin{xltabular}{\linewidth}{>{\raggedright\arraybackslash}p{0.20\linewidth}XXX}
\toprule
\thead{层} & \thead{是什么} & \thead{是谁的负债} & \thead{谁能拿着} \\
\midrule
\textbf{第一层} & 央行准备金 & 央行 & \textbf{只有银行}。银行之间清算专用，你我一辈子碰不到 \\
\textbf{第二层} & 商业银行存款 & 商业银行 & 你、我、企业、政府。日常说的"钱"几乎全在这一层 \\
\textbf{第三层} & 现金 & 央行 & 所有人。唯一能真攥在手里的那种 \\
\bottomrule
\end{xltabular}
\end{tablewrap}

拿电网打比方就清楚了。准备金是高压输电线上的电，只在变电站之间跑，你家电器插上去会炸。存款是入户的 220 伏，能烧水做饭。现金是充电宝里那点电，随身带着，量很小。

而把高压变成 220 伏的那台变压器，就是银行放贷。

这个比方能一路推下去：往高压线上加多少电压，你家插座都不会变一伏，除非变压器在工作。变压器停了，高压线上电压再高也点不亮一盏灯。第 8、9 章要讲的"传导失灵"，说的就是变压器不转。

\begin{pullquote}{}
第一层的钱和第二层的钱之间没有自动通道。这就是为什么央行"放水"，你的口袋可能一滴都感觉不到。
\end{pullquote}

2009 到 2019 年那十年，最能说明这副眼镜的用处。

发达国家搞了十年\term{量化宽松}{quantitative easing}\index{033@量化宽松}，央行资产负债表膨胀了好几倍。当时一大批人斩钉截铁地预言：恶性通胀马上就到，赶紧买黄金、囤罐头。

十年过去，CPI 通胀没来。

因为 QE 创造的是第一层。央行从银行手里买走国债，付给银行的是准备金。准备金躺在央行账户里，只能在银行之间转来转去，它不会自己长腿跑进陈宁的工资卡。

那些钱去哪儿了？它把国债和其他安全资产的收益率压到极低，逼着所有拿着钱的机构去买风险更高的东西：股票、房子、信用债、私募。所以那十年物价很平静，资产价格涨疯了。

存钱的人一无所获，持有资产的人富得流油。这不是道德问题，是管道问题。水注进了第一层的池子，而资产价格恰好是第一层水位最先溢出的地方。

要让钱真正进实体经济只有一条路：绕过银行的放贷决策，让钱直接落在居民和企业的存款账户上。这件事只有\textbf{财政}干得了。政府发工资、发补贴、给工程队结算，钱是直接打进第二层的，收款人的余额当场增加。

\begin{egbox}{举个栗子：同样是"放水"，结局差了四十年}
\textbf{2010 年代：QE。}美联储买国债，创造准备金（第一层）。结果是股票和房地产大涨，CPI 常年低于 2\% 的目标，美联储那些年发愁的是通胀太低。

\textbf{2020 到 2021 年：直接支票。}财政部把钱打进几乎每一户居民的账户（第二层），叠加失业救济加码、供应链断裂、劳动力紧张。结果是 2022 年 CPI 冲上四十年未见的高位。

同一个国家，同一批人，前后不到十年。差别不在"印了多少"，在\textbf{钱进了哪一层，以及拿到钱的人有没有立刻去买真东西}。

所以本书从不说"印钞必然通胀"。通胀是支出撞上供给产生的。2026 年通胀的粘性更多来自关税、去全球化和中东局势推高的能源成本，全是供给侧的事。

\end{egbox}

\section{稳定币是一家开在链上的窄银行}

把这一章的规律拿去套 2026 年最热闹的那个新东西，严丝合缝。

\term{稳定币}{stablecoin}\index{065@稳定币}怎么运作？你交 1 美元给发行方，发行方拿这 1 美元去买短期美国国债和货币市场工具，然后发给你一个代币，承诺随时按 1 比 1 换回美元。

翻译成这本书的语言：代币是发行方的负债，撑着它的是资产端的短期国债。它非常接近教科书里的\term{窄银行}{narrow bank}\index{082@窄银行}，只吸收资金、只投最安全的短期资产、不放贷，因此不创造新货币，只是把已有的钱换了条轨道。

截至 2026 年 8 月，全球稳定币市值在 3000 亿美元量级。美国的 GENIUS Act 已于 2025 年 7 月生效，为美元稳定币建立了联邦监管框架。监管者盯的问题和这本书一样：储备是什么，能不能随时兑付，发行方破产了谁赔。

\begin{warnbox}{小心}
"稳定"两个字是营销，不是物理定律。它稳不稳取决于资产端的质量和赎回机制，和一家银行会不会被挤兑是同一道题。历史上已经有算法稳定币在几天内归零的案例。

本书不对任何具体加密资产给投资建议，只提醒你用同一把尺子量它：\textbf{谁的负债？拿什么撑着？急着要钱的那天，那个"撑着"的东西卖得掉吗？}

\end{warnbox}

银行存款靠贷款和资本撑着，现金靠国家的征税能力撑着，稳定币靠短期国债撑着。任何"钱"，都是某个主体的负债。

\section{三百万亿里，只有 4\% 当过纸}

\begin{statrow}{3}
  \statitem{300+}{万亿}{中国 M2 规模（元，量级）}
  \statitem{10+}{万亿}{流通中现金 M0（元，量级）}
  \statitem{1.41}{\%}{商业银行净息差（2026 年二季度末）}
\end{statrow}

\begin{barfig}
  \barrow{银行存款（第二层，贷款创造）}{96}{约 96\%}
  \barrow{现金（第三层，央行印的）}{4}{约 4\%}
\end{barfig}

广义货币三百多万亿元，流通中的现金只有十几万亿元。95\% 以上的"钱"从没当过一天纸，它们是商业银行给某个陈宁、某家模具厂、某个地方平台批贷款时敲键盘敲出来的。

钱和债是同一枚硬币的两面。钱要变多就得有人多欠债；债要减少，钱必然变少。整套货币体系靠的是"持续有人愿意借"。

那如果哪天真的没人愿意借了呢？

日本试过。试了三十年，至今没试出一个能让别人照抄的答案。这件事目前没人知道解法。

\bookbreak

\begin{takeaway}{本章一句话}
绝大部分货币不是印出来的，是借出来的。银行放贷时同时创造资产和负债，还款时又把它注销掉。

\begin{itemize}
  \item 贷款创造存款。准备金是事后补的，不是原料。货币乘数那套讲法把因果讲反了。
  \item 真正的约束是借款人的需求、资本充足率、1.41\% 的净息差、监管和央行定的那个价，不是"手上有多少存款"。
  \item 钱分三层。QE 只灌第一层，所以它推资产价格胜过推物价；财政直接打钱灌第二层，所以 2021 年通胀来了。
  \item 碰到任何新的"钱"，先问：谁的负债，拿什么撑着。
\end{itemize}

\end{takeaway}

\begin{bridgebox}
\textbf{下一章}既然钱主要是商业银行造的，那个天天上新闻的央行到底在管什么？它管的不是水量，是水价和水温。
\end{bridgebox}


\chapter{央行不是印钞厂，是水闸和温度计}
\chaptermeta{08}

\begin{ledetext}
上一章证明了钱主要是商业银行造的。那新闻里天天出现的央行到底在管什么？它管的不是水量，是水价。
\end{ledetext}

\section{先拆掉那台想象中的印钞机}

说起央行，很多人脑子里是一间轰鸣的车间：崭新的钞票像瀑布一样吐出来，旁边站着个戴白手套的人在按加速键。

这幅画里唯一对的部分是，纸币确实要印。

但它只占全社会货币的百分之四上下（第 7 章那张图：M0 十几万亿元，M2 三百多万亿元）。而且印钞不是"投放货币"，是换包装。陈宁去 ATM 取 800 块，是把银行存款换成现金，社会上的钱一分没多。印钞的节奏是被公众的取现需求牵着走的，不是央行想印多少印多少。

它真正的工作只有一句话。

\begin{pullquote}{}
央行决定不了这个经济体里有多少钱，它决定的是钱的价格。
\end{pullquote}

\section{三重身份，和一份长短不一的任务书}

这个机构有点分裂，因为它同时干着三份工作。

\begin{flowlist}
  \flowitem{01}{发行的银行}{它垄断法定货币的发行权。你钱包里那张纸的信用，来自这个机构和它背后国家的征税能力。这是最古老、现在也最不吃力的一份工作。}
  \flowitem{02}{银行的银行}{所有商业银行都在央行开户，账上那笔钱就是第 7 章说的准备金。银行之间转账在这本账上清算，银行缺钱来这里借，银行快倒了来这里求救。\textbf{它九成的日常工作在这一层。}}
  \flowitem{03}{政府的银行}{财政部的账户开在央行，税收进来、支出出去都走这里。但有一条红线：多数国家的法律禁止央行直接给财政透支或在一级市场买国债。这条线一旦模糊，货币的锚就松了。中国《人民银行法》也有明确规定。}
\end{flowlist}

它要往哪儿开，取决于被写进法律的那份任务书。各国的任务书长得很不一样，这直接决定了 2026 年它们为什么会走散。

\begin{tablewrap}
\begin{xltabular}{\linewidth}{>{\raggedright\arraybackslash}p{0.20\linewidth}XX}
\toprule
\bthead{央行} & \bthead{法定目标} & \bthead{后果} \\
\midrule
\textbf{美联储} & \textbf{双重使命}：物价稳定加充分就业 & 两个目标打架时（通胀高、失业也高）必须做痛苦取舍，也因此天然容易被政治拉扯 \\
\textbf{欧洲央行} & 以物价稳定为首要目标 & 目标单一、纪律性强，但面对二十个财政不统一的成员国，工具受限 \\
\textbf{中国人民银行} & 保持币值稳定并以此促进经济增长，同时兼顾金融稳定、汇率、就业和结构性任务 & 目标最多，工具箱也最杂。中国这一整套"定向""结构性"工具就是从这儿长出来的 \\
\bottomrule
\end{xltabular}
\end{tablewrap}

目标越多，可用的借口越多，独立性越难守。这不是批评，是结构。

\section{六件工具，六个能推下去的比喻}

\textbf{一、政策利率，总水阀。}这一件最要紧，其余五件都是配角。

美联储的做法叫地板系统。它给银行存在美联储的准备金付息（\term{准备金余额利率}{IORB}\index{090@准备金余额利率}），又给货币基金这类非银机构开了个隔夜逆回购窗口（\term{隔夜逆回购}{ON RRP}\index{020@隔夜逆回购}）。前者大致构成利率上沿，后者构成下沿，市场利率被夹在中间，这条通道叫\term{利率走廊}{interest rate corridor}\index{032@利率走廊}。

\begin{egbox}{举个栗子：两个报价怎么关住一整个市场}
2026 年 8 月，联邦基金利率目标区间是 3.50\% 到 3.75\%。ON RRP 利率在区间下沿附近，IORB 在区间内靠近上沿。

\textbf{下沿怎么守住。}一家货币市场基金账上趴着 84 亿美元隔夜资金。如果市场上有人只肯出 3.42\%，它宁可把钱放进美联储的 ON RRP 拿 3.50\%，零风险。所以没人能用明显低于 3.50\% 的价格借到隔夜钱。

\textbf{上沿怎么压住。}一家银行若能在市场上以 3.55\% 借到钱，转手存进美联储拿 IORB（假设 3.65\%），就是 10 个基点的无风险利差。84 亿美元的规模，年化口径 \textbf{840 万美元}。这种套利会立刻把市场利率压回来。

央行不需要投放任何钱。它把两个报价挂在墙上，整个隔夜市场就被夹住了。\hlmark{这就是"管价格"的字面意思。}

\end{egbox}

中国的路子不同，逻辑一样：以 7 天期逆回购利率为主要政策利率，截至 2026 年 8 月是 \textbf{1.4\%}；用中期借贷便利（MLF）投放中期资金；再通过 18 家报价行的 LPR 机制传导到贷款端。还有一件重型工具是存款准备金率，它调的是银行的长期资金成本和流动性，不是第 7 章批倒的那个货币乘数。

\textbf{二、公开市场操作，每天的微调。}央行在市场上买卖债券、做正逆回购，把银行体系的资金量调到让市场利率贴着政策利率跑的位置。这活儿天天干，像给恒温空调做实时补偿。

\textbf{三、QE 和 QT，把手伸到曲线的右边。}政策利率管得住隔夜，管不住 10 年期。于是央行下场大量买入长期国债，压低长端收益率和\term{期限溢价}{term premium}\index{042@期限溢价}，这叫量化宽松；反过来停止再投资、让持仓自然到期缩表，叫量化紧缩。日本央行 2026 年一边加息一边继续削减国债购买，计划到 2027 年一季度降到每月约 2 万亿日元，走的就是后一条路。

\textbf{四、前瞻指引，用嘴巴影响预期。}今天的长期利率，说穿了是市场对未来一连串短期利率的预期的平均。央行只要让市场相信"我未来两年都不加息"，长端利率现在就下来了。一分钱没花，效果已经发生。这是本世纪央行发明的最便宜的工具。

\begin{databox}{2026：沟通范式正在转向}
2026 年，美联储主席凯文·沃什明确表示要\textbf{淡化前瞻指引}。

理由不难猜。当央行反复承诺又反复被现实打脸（比如 2021 年那句"通胀是暂时的"），指引就从降低不确定性变成了制造不确定性，公信力被一次次透支。

沃什同时强调，美联储"不会对任何市场走势背书"，以及"必要和适当时会毫不犹豫地采取行动"。翻译过来：\hlmark{别指望我提前给你剧本，也别指望我托底你的持仓。}

这是过去十几年货币政策沟通方式的一次真转向。对市场的直接后果是，预期的锚变松了，波动会变大。

\end{databox}

\textbf{五、最后贷款人，消防队。}银行的生意天生借短贷长，只要储户同时来取钱，再健康的银行也会倒。所以危机时得有一个不会缺钱的人站出来放款。19 世纪的英国人白芝浩把这件事写成三条铁律，一百五十年没过时。

\begin{chainflow}
\chainnode{只救还有偿付能力的}\chainarrow \chainnode{按惩罚性高利率放款}\chainarrow \chainnode{凭优质抵押品无限量供应}
\end{chainflow}

消防队这个比喻能一路推下去。只救还救得回来的房子，是因为往一堆焦炭上浇水只是浪费水压；收水费，是为了让人别把消防队当免费喷淋用；水管够，是因为火场里最先要扑灭的不是火，是"水不够"这个念头。恐慌只有在"钱管够"这句话被相信的时候才会停。

\textbf{六、宏观审慎和结构性工具，滴灌。}总水阀只有一个，可经济的不同角落干湿不均。于是有了逆周期资本缓冲、房贷首付比例、行业贷款集中度这类\term{宏观审慎}{macroprudential}\index{022@宏观审慎}工具。中国还额外配了一整套定向再贷款：支农支小、科技创新、绿色低碳，给特定方向的贷款提供低成本资金。

这类工具精准，代价是它在干本该财政干的活。\hlmark{决定谁能拿到便宜的钱，是一个分配决策。}

\section{怎么读一张央行的资产负债表}

新闻里天天说"美联储缩表""央行扩表"，听着玄，其实那张表比你家账本还简单。

\begin{bsheet}
  \bsside{QE 之前 · 某央行（亿元）}{%
    \bsrow{持有的国债}{7840}
    \bsrow{对金融机构的贷款}{1620}
    \bsrow{外汇与黄金}{460}
    \bsrowtotal{资产合计}{9920}
  }
  \bsside{负债 + 净值}{%
    \bsrow{银行准备金}{5930}
    \bsrow{流通中现金}{3140}
    \bsrow{政府存款}{610}
    \bsrow{资本}{240}
    \bsrowtotal{合计}{9920}
  }
\end{bsheet}

资产端是它买来的和借出去的，负债端是它欠出去的。QE 就是资产端买入国债、负债端同时长出准备金，两边一起膨胀。QT 反过来，两边一起收缩。

\begin{egbox}{举个栗子：同样买 780 亿国债，结局可以完全不同}
\textbf{情形 A，卖方是一家银行。}央行资产端国债加 780 亿，负债端准备金加 780 亿；这家银行资产端国债减 780 亿、准备金加 780 亿。陈宁和你我的存款：\textbf{一分没动}。只有第一层的钱变多了，广义货币纹丝不动。

\textbf{情形 B，卖方是一家养老基金。}养老基金不能在央行开户，钱得打进它的开户行。于是央行准备金加 780 亿（记在这家银行名下），这家银行资产端准备金加 780 亿、负债端"养老基金的存款"加 780 亿。\textbf{第一层和第二层同时增加}，广义货币真的变多了。

同一句"央行买了 780 亿国债"，两种结局。所以看到"央行大放水"这四个字，先问一句：\hlmark{钱是从谁手里买的，落进了哪一层？}

\end{egbox}

\section{独立性：一个 1970 年代的老问题，2026 年又新鲜了}

为什么全世界要把货币政策交给一群不用选举、任期很长、谁也管不着的技术官僚？

答案叫\term{时间不一致}{time inconsistency}\index{049@时间不一致}。政客的时间尺度是下一次选举，而降息的好处（就业、股市、支持率）来得快，坏处（通胀）来得慢，还常常落到下一任头上。任何把利率交给政治周期的国家，长期都会得到系统性偏高的通胀。1970 年代那笔学费是用两位数通胀和一场深度衰退交的。

把方向盘交给一个不看民调的人，不是因为他更聪明，是因为他敢做那件正确但不受欢迎的事。

2026 年，这个老问题重新变得刺手。美联储主席换成了凯文·沃什，政治层面对货币政策的压力明显上升，而通胀偏偏没回到目标。沃什那两句"不会对任何市场走势背书""必要和适当时会毫不犹豫地采取行动"，一半说给市场听，一半说给白宫听。

\begin{deepbox}{进阶：可以跳过 · 这场争论和作者的倾向}
\textbf{支持独立性的一方：}历史证据很硬。独立性高的央行长期通胀更低、通胀预期更稳，而且没有以更高的失业作为代价。放弃它就是拿几十年的信用去换几个季度的舒服。

\textbf{质疑的一方：}货币政策早就不只是调利率了。QE 会推高资产价格、拉大贫富差距；救谁不救谁是重大的分配选择；结构性工具直接决定哪个行业拿到便宜的钱。这些都是政治决策，却由一群不经选举的人做出，问责在哪里？

\textbf{我的倾向是：}两边说的都对，但它们说的不是同一件事。设定利率这件事应当独立，证据足够充分，不该动。而 QE、定向再贷款、救助名单这些带明确分配后果的操作，应当由财政和立法机构明确授权、事后被审查。把这两类权力混在一起打包成"央行独立性"来辩护，才是过去十五年真正的漏洞。

\end{deepbox}

\section{2026 年 8 月：四家央行走散了}

过去二十年我们习惯了一件事：美联储一转向，全世界跟着转。2026 年这个默认设定失效了。

\begin{tablewrap}
\begin{xltabular}{\linewidth}{>{\raggedright\arraybackslash}p{0.20\linewidth}XXr}
\toprule
\bthead{央行} & \bthead{2026 年 8 月的位置} & \bthead{方向} & \bthead{关键利率} \\
\midrule
\textbf{美联储} & 2026 年已连续 5 次按兵不动（截至 7 月），FOMC 称通胀"仍然高企"，市场甚至讨论过加息 & 观望偏鹰 & 3.50\%–3.75\% \\
\textbf{欧洲央行} & 2026 年 6 月 11 日上调存款便利利率，三年来首次加息 & 转向收紧 & 已上调 \\
\textbf{日本央行} & 2026 年 6 月 16 日以 7 比 1 加息 25 个基点，为 1995 年 9 月以来最高，同时继续削减国债购买 & 正常化推进中 & 1.0\% \\
\textbf{中国人民银行} & 基调"适度宽松"，强调"精准有效""适时调整"，支持科技创新、提振消费、民营小微、绿色转型 & 宽松但克制 & 1.4\% \\
\bottomrule
\end{xltabular}
\end{tablewrap}

\begin{statrow}{3}
  \statitem{1.0}{\%}{日本央行政策利率，1995 年 9 月以来最高}
  \statitem{1.4}{\%}{中国 7 天期逆回购政策利率}
  \statitem{3.5}{\%}{中国 5 年期以上 LPR（2026 年 8 月 20 日）}
\end{statrow}

一个在收紧，一个在正常化，一个在观望，一个在宽松。

利差就是资本的坡度。日本的利率从负 0.1\% 一路爬到 1.0\%，那些年借便宜日元去买全球资产的套利盘就得重新算账；欧洲开始加息而美国不动，欧美之间的资金流向和汇率就得重新找平衡点。

\begin{pullquote}{}
央行不同步，本身就是资本流动和汇率波动的源头。世界上不再有一个统一的"钱的价格"了。
\end{pullquote}

\section{两个必须扔掉的直觉}

\begin{mythbox}{常见误解}
\mvfrow{误解}{"央行想让经济好，就能让经济好。降息、放水，经济自然起来。"}
\mvfrow{事实}{央行只能改变钱的价格和可得性，它\textbf{替不了任何人做借钱的决定}。传导链有四个地方会断：利率降到接近零仍无人肯借（流动性陷阱）；企业和家庭忙着还债，目标从赚钱变成减债（资产负债表衰退）；银行资本不足或风险偏好塌了，宁可买国债也不放贷；以及最朴素的一种，根本没有值得投的项目。\hlmark{你能把马牵到河边，逼不了它喝水。}}
\end{mythbox}

\begin{mythbox}{常见误解}
\mvfrow{误解}{"降息是利好，说明经济要转好了。"}
\mvfrow{事实}{这个说法是错的，而且通常刚好反过来。\textbf{央行降息，是因为它看见了需要被救的东西。}降息是对坏消息的反应，不是好消息本身。历史上多次深度衰退的起点，都伴着一轮急促的降息。该盯的从来不是"降了没有"，是"为什么降、降得多急"。连续大幅降息往往意味着某个地方在着火。}
\end{mythbox}

\begin{warnbox}{小心}
把央行动作当买卖信号，是散户最常见的亏钱方式之一。政策利率是\textbf{已经发生}的事，资产价格交易的是\textbf{预期}。降息公布那一刻，市场可能早把三次降息算进价格里了，靴子落地反而是利空。本书不提供任何交易建议，只提醒你：新闻标题和你的账户之间，隔着一整条下一章要拆的传导链。

\end{warnbox}

\bookbreak

\begin{takeaway}{本章一句话}
央行管的不是钱的数量，是钱的价格。它是水闸和温度计，不是印钞厂。

\begin{itemize}
  \item 三重身份：发行的银行、银行的银行、政府的银行。九成的日常工作在第二层。
  \item 工具箱六件：政策利率、公开市场操作、QE 与 QT、前瞻指引、最后贷款人、宏观审慎与结构性工具。
  \item 2026 年四家央行走散：欧洲加息、日本升到 1.0\%、美联储停在 3.50\%–3.75\%、中国维持 1.4\% 的适度宽松。不同步本身就是波动的来源。
  \item 降息通常是坏消息的反应。央行能把马牵到河边，逼不了它喝水。
\end{itemize}

\end{takeaway}

\begin{bridgebox}
\textbf{下一章}央行拧的是最上游那个阀门。这股水要经过几道弯才能流到陈宁每月 15 号被划走的那笔钱上？我们把管道一节一节拆开。
\end{bridgebox}


\chapter{利率的传导链：从政策利率到你的房贷}
\chaptermeta{09}

\begin{ledetext}
央行在最上游拧了一下阀门。这股水要过七道闸才能流到陈宁每月 15 号被划走的那笔钱上，每一道都可能放大它、拖慢它，或者把它堵死。
\end{ledetext}

\section{从 1.4\% 到 13471 元，中间隔着七道闸}

陈宁的房贷 300 万，30 年，等额本息，执行利率 3.5\%，月供 13471 元。他和妻子加起来到手 3 万 4 出头，月供占了四成。

央行公布的那个数字是 1.4\%。陈宁账单上的数字是 3.5\%。

这两个数字之间隔着一整条管道。

\begin{flowlist}
  \flowitem{01}{政策利率}{中国是 7 天期逆回购利率，截至 2026 年 8 月为 1.4\%；美国是联邦基金利率目标区间 3.50\% 到 3.75\%。全链条的锚。\emph{信号在这里最强、最干净。}}
  \flowitem{02}{银行间货币市场}{中国看 \term{DR007}{存款类机构质押式回购利率}\index{095@DR007}，也就是银行之间拿国债作抵押借 7 天钱的实际成交价；美国看 \term{SOFR}{担保隔夜融资利率}\index{096@SOFR}。\emph{季末、缴税、大额国债发行都会让它偏离政策利率。}}
  \flowitem{03}{同业存单与短期国债}{银行发同业存单筹 3 个月、6 个月的钱，价格直接挂在货币市场利率上。\emph{信用和期限的因素从这里开始掺进来。}}
  \flowitem{04}{收益率曲线}{从 1 个月到 30 年，市场报出一整条价格。\emph{整条链上信号最容易变形的一环，因为长端不归央行管。}}
  \flowitem{05}{银行的负债成本}{银行的钱从存款、同业存单、央行借款三处来。\hlmark{存款利率是刚性的。}没人会因为央行降息就接受活期变负数，银行也不敢把定期利率砍到没人肯存。\emph{信号在这里被结结实实卡掉一大截。}}
  \flowitem{06}{贷款定价}{中国走 LPR 加点，美国各家银行走自己的模型。\emph{银行的资本状况和风险偏好，决定它愿意让出多少。}}
  \flowitem{07}{陈宁的月供}{还得等到他的重定价日。\emph{信号走到这里，可能比出发时晚了整整一年。}}
\end{flowlist}

第 7 章那个电网的比喻，在这里能接着往下推。高压电送到你家插座要走上百公里，每一段线路都有损耗，每一台变压器都要吃掉一点电压。这条链最大的那处损耗在第 5 环：银行的存款成本降不下来，贷款利率就让不出去。

"央行降息 25 个基点"和"陈宁月供少 415 块"之间不是等号，是一段漫长而且可能中断的旅程。

\section{收益率曲线：一张能看出天气的图}

第 4 环那条\term{收益率曲线}{yield curve}\index{053@收益率曲线}值得单独讲，它是整个金融市场信息密度最高的一张图。

画法极简单。横轴是期限（1 个月、1 年、5 年、10 年、30 年），纵轴是对应的国债收益率，把点连起来就是一条线。

\textbf{正常形态向上倾斜。}把钱借出去 10 年，要承担的不确定性远多于借出去一天：通胀可能失控，你可能急着用钱，发行人可能变差。市场为这份不确定性索取的额外补偿，叫\term{期限溢价}{term premium}\index{042@期限溢价}。

曲线变平，说明市场认为利率不会再往上走了。

短端反过来高过长端，叫倒挂。这个形状很刺眼：\hlmark{市场在说，我宁可锁定 10 年拿一个更低的利率，也不愿意赌未来还有这么高的短期利率。}翻译过来就是，大家预期央行将来会大幅降息，而央行大幅降息通常只有一个原因。

倒挂在历史上是可靠的衰退先行指标，还有第二个更机械的理由：\textbf{银行的商业模式是借短贷长。}短端高过长端，银行的息差被直接压扁，放贷动力下降，信用收缩，然后经济真的坏了。它不只是预言，它本身就参与制造了那个结果。

\begin{deepbox}{进阶：可以跳过 · 这个指标失灵了吗}
2022 到 2023 年，美国出现了史上持续最久的倒挂之一，几乎所有人都在等衰退。衰退迟迟没来，经济反而维持了相当的韧性。

解释有三种，都讲得通，都没被证实。一是 QE 长期压低了期限溢价，曲线形状的含义已经和历史上不一样了。二是疫情后居民和企业锁定了大量长期低息负债，资产负债表异常健康，对短期利率不敏感。三是财政持续扩张，对冲掉了货币紧缩。

\textbf{这个指标是不是已经失效，目前没人知道答案。}本书的建议是把倒挂当成一个值得警惕的信号，不是一个可以照着下注的时钟。任何"指标一响就该卖出"的说法都不构成投资建议，历史样本只有寥寥几次，根本撑不起那种确定性。

\end{deepbox}

\section{LPR：一条你能亲眼看见的链子}

中国的贷款定价机制透明得罕见，每一环都能看清。

\begin{chainflow}
\chainnode{7 天逆回购 1.4\%}\chainarrow \chainnode{18 家报价行加点报价}\chainarrow \chainnode{同业拆借中心每月 20 日发布 LPR}\chainarrow \chainnode{房贷 = LPR ± 固定加点}
\end{chainflow}

请盯住"\textbf{固定}加点"这四个字。陈宁签合同那天谈定的加点，在整个 30 年里\hlmark{一个字都不会改}。会变的只有 LPR 本身，而且要等到他的重定价日才在月供上生效。

拿这套机制去解释一条真实的新闻。

\begin{talkbox}
  \talkline{新闻}{"2026 年 8 月 20 日，1 年期 LPR 报 3.0\%，5 年期以上 LPR 报 3.5\%，均连续多月持平。"}
  \talkline{陈宁}{"不是说货币政策适度宽松吗，怎么一直不降？"}
  \talkline{翻译}{拆成三层看。第一，政策锚没动，7 天期逆回购稳在 1.4\%，报价行的资金成本基准没变。第二，报价行不想动，截至 2026 年二季度末商业银行净息差只有 1.41\%，历史低位，主动下调加点等于自己割自己。第三，负债端是刚的，存款利率降不动，成本降不下来，贷款利率就腾不出空间。\textbf{LPR 按兵不动不是"政策没意愿"，是传导链第 5 环卡住了。}}
\end{talkbox}

一条数据不动，往往比它动了更有信息量，前提是你知道它被链条上哪一环卡住。

\begin{statrow}{3}
  \statitem{1.41}{\%}{商业银行净息差（2026 年二季度末，历史低位）}
  \statitem{3.5}{\%}{5 年期以上 LPR（2026 年 8 月 20 日）}
  \statitem{1.4}{\%}{7 天期逆回购政策利率}
\end{statrow}

\section{那 80 个基点，值 48 万}

"加点"这两个字听着像四舍五入的小数点。把它换算成人民币。

统一条件：300 万元，30 年，等额本息，LPR 为 3.5\%。

\begin{tablewrap}
\begin{xltabular}{\linewidth}{>{\raggedright\arraybackslash}p{0.20\linewidth}Xrrr}
\toprule
\thead{情形} & \thead{执行利率} & \thead{月供} & \thead{30 年总利息} & \thead{与基准之差} \\
\midrule
\textbf{LPR 减 30 个基点} & 3.20\% & 12974 元 & 约 167.1 万 & 每月少 497 元 \\
\textbf{LPR 加 0（陈宁）} & 3.50\% & 13471 元 & 约 185.0 万 & 基准 \\
\textbf{LPR 加 50 个基点} & 4.00\% & 14323 元 & 约 215.6 万 & 每月多 852 元 \\
\bottomrule
\end{xltabular}
\end{tablewrap}

\begin{barfig}
  \barrow{3.20\%（LPR−30bp）}{91}{12974 元/月}
  \barrow{3.50\%（LPR+0）}{94}{13471 元/月}
  \barrow{4.00\%（LPR+50bp）}{100}{14323 元/月}
\end{barfig}

最上面和最下面差 80 个基点。听着是 0.8\%，一个可以忽略的数。

换算成钱：每月 1349 元，三十年 48.5 万元。

而这 80 个基点，往往只取决于签合同那年的市场行情、你的资质，以及你有没有多问过一句"加点还能不能谈"。

\begin{egbox}{举个栗子：LPR 降 25 个基点，陈宁少还多少}
\textbf{情形一，新贷款。}300 万、30 年、等额本息，利率从 3.5\% 降到 3.25\%：月供从 13471 元降到 \textbf{13056 元}，每月少 \textbf{415 元}，全周期少还利息约 \textbf{14.9 万元}。

\textbf{情形二，陈宁这种存量借款人。}他 2024 年 6 月放款，已经还了两年，此时剩余本金约 288.3 万元，剩余期限 28 年。同样降 25 个基点，重定价日之后月供从 13471 元降到约 \textbf{13080 元}，每月少 \textbf{391 元}。

比新贷款少省 24 块。原因是他已经还掉一部分本金，能吃到折扣的基数变小了。\hlmark{降息对存量借款人的帮助，永远小于对新借款人的帮助。}

\end{egbox}

\section{"新闻说降息了，我月供怎么没变"}

这是留言区里出现频率最高的一句抱怨。答案藏在三个词里。

\textbf{重定价日。}浮动利率合同里会约定一个日期，通常是每年 1 月 1 日，或者每年放款日的对应日。只有到了那天，银行才会取用最近一次公布的 LPR 重算月供。8 月降的息，陈宁的重定价日在 1 月 1 日，他要等到明年，中间四个半月照旧按老利率还。

\textbf{固定加点。}LPR 降了，加点不跟着降。当年签的是"LPR 加 80 个基点"，那就永远比基准贵 0.8 个百分点。这是合同，不是政策。

\textbf{批量调整。}也有例外。2023 到 2024 年，中国对存量房贷利率做过集中批量下调，绕开了"加点不能改"这条规则。那是特殊时期的特殊操作，是政策选择，不是常规机制。别把它当成一种可以指望的权利。

\begin{mythbox}{常见误解}
\mvfrow{误解}{"央行一降息，我的房贷马上就便宜了。"}
\mvfrow{事实}{这个说法是错的，中间隔着七个环节。政策利率要先传到银行间市场，再传到 LPR 报价，报价行还得\textbf{愿意}让出自己那 1.41\% 的净息差。就算 LPR 真降了，你还要等自己的重定价日，而加点一分不会动。\hlmark{从政策利率变动到你的月供变动，快则几个月，慢则一年多，也完全可能中途蒸发掉。}}
\end{mythbox}

\section{传导断在哪儿}

上面讲的是"变慢"。更麻烦的是"断掉"。

\begin{flowlist}
  \flowitem{01}{资产负债表衰退}{经济学家辜朝明对日本泡沫破裂后的诊断：资产价格暴跌而债务原封不动，企业和家庭的目标会从"利润最大化"变成"\textbf{负债最小化}"。这时候利率降到 0 也没用，一个净资产为负的人，赚到的每一块钱都会拿去还债。日本走完这个过程用了将近三十年。}
  \flowitem{02}{银行惜贷}{银行资本不足，或者刚吃过一轮坏账，风险偏好整体塌了。它宁可买国债、做同业，也不放给中小企业。\hlmark{钱在银行体系里空转，就是不出门。}}
  \flowitem{03}{没有值得投的项目}{最朴素也最致命的一种。企业算完账发现需求见顶、产能过剩，借钱扩产是纯亏。利率是 3\% 还是 1\%，改不了那张现金流预测表上的答案。}
  \flowitem{04}{利率下限}{名义利率很难显著低于零，因为存款人可以取现金揣着，而现金的名义利率是 0。政策空间到这儿就见底了，剩下的只能靠 QE 和财政。}
\end{flowlist}

中国当下的提前还贷潮，同时踩中了第 1 条和第 3 条：家庭在主动缩表，把手里的存款换成"少欠银行的钱"。按第 7 章的结论，提前还贷会直接注销货币。

于是又是那个别扭的组合：政策在放松，利率在低位，社会感受到的钱依然是紧的。

\begin{pullquote}{}
利率是价格信号，不是命令。价格能改变一个人的意愿，改变不了他"我现在不想借钱"的决心。
\end{pullquote}

顺带说一句机制上的判断，不是建议：提前还贷相当于买入一笔无风险、税后年化收益等于你贷款利率的资产。3.5\% 的房贷，提前还款就是锁定 3.5\% 的无风险回报。这笔账划不划算，取决于你有没有更好的去处、以及那笔钱拿走之后你的现金流还撑不撑得住。这两件事只有你自己知道。

\section{央行控短端，市场定长端}

这条链上最反直觉的一件事留在最后。

央行牢牢控制着曲线最左端，隔夜和 7 天的利率想定多少就是多少。曲线右端的 10 年、30 年，是市场用真金白银投票投出来的，央行只能影响，不能决定。

\begin{egbox}{举个栗子：2026 年 8 月，一件"不该发生"的事}
2026 年 8 月底，美国 10 年期国债收益率升至 \textbf{2025 年 1 月以来的最高水平}。同期美联储的政策利率一动没动，2026 年已连续 5 次维持在 3.50\% 到 3.75\%。

短端纹丝不动，长端往上冲。至少四股力量在同时起作用。

\textbf{①通胀预期。}关税、能源（中东局势）、去全球化都在抬高对未来物价的预期，长期债权人要求更多补偿。\textbf{②期限溢价回归。}QE 时代央行是长债最大的买家，人为压低了这份补偿，央行退场后它自然反弹。\textbf{③财政供给。}美国国债规模 2026 年 8 月突破 \textbf{40 万亿美元}，卖方源源不断，价格自然往下、收益率往上。\textbf{④买家结构在变。}截至 2025 年底，黄金已占全球央行储备 27\%（2024 年底还是 20\%），超越美国国债成为第一大储备资产，美债占比降到 22\%。

沃什本人也承认"名义和实际国债收益率大幅上升"，并指出"从宏观角度看，市场价格表明金融环境已趋紧"。\hlmark{美联储没有收紧，市场替它收紧了。}

\end{egbox}

这件事对普通人的意义很直接。美国的房贷利率主要挂在 10 年期国债收益率上，不是挂在联邦基金利率上。所以"美联储不加息"完全不等于"美国人的房贷不涨"。中国的长期国债收益率同样不由 7 天逆回购单独决定，它反映的是市场对未来十年增长和通胀的整体判断。

\begin{mythbox}{常见误解}
\mvfrow{误解}{"长期利率是央行定的。央行不加息，长期利率就涨不上去。"}
\mvfrow{事实}{央行直接控制的只有曲线最左边那一小段。10 年、30 年的利率由\textbf{通胀预期、期限溢价和债券供求}共同定价。2026 年 8 月的美国就是活标本：政策利率五次按兵不动，10 年期收益率创下 2025 年 1 月以来新高。央行可以影响长端，这正是 QE 存在的理由，但财政供给和通胀预期一起发力的时候，它未必拧得过市场。}
\end{mythbox}

\begin{warnbox}{小心}
本章所有算术都是为了让你看懂机制，不是为了给你一个"该不该提前还贷""该不该换固定利率"的答案。那个答案取决于你的现金流稳定性、有没有更好的资金去处，以及你对未来利率的判断。最后这一项，没有人知道。\textbf{本书不构成投资或借贷建议。}

\end{warnbox}

\bookbreak

\begin{takeaway}{本章一句话}
政策利率到你的月供之间有七道闸，每一道都会削弱、拖慢甚至截断信号。央行控短端，市场定长端。

\begin{itemize}
  \item 链条：政策利率 → 银行间利率（DR007/SOFR）→ 同业存单与短债 → 收益率曲线 → 银行负债成本 → LPR 加点 → 你的月供。
  \item 2026 年 8 月 LPR 连续多月不动，卡点在 1.41\% 的净息差和降不下来的存款成本，不是政策没意愿。
  \item 300 万 30 年，加点从减 30bp 到加 50bp，月供差 1349 元，总利息差 48.5 万元。
  \item 传导会断在四个地方：资产负债表衰退、银行惜贷、没有值得投的项目、利率下限。
\end{itemize}

\end{takeaway}

\begin{bridgebox}
\textbf{下一章}说了半天利率，它到底长在哪个市场上？答案是一个规模比全球股市大得多、每天决定着陈宁的月供和一国财政命运、却几乎没人在饭桌上聊起的地方。
\end{bridgebox}


\chapter{债券：世界上最大、最无聊、最要命的市场}
\chaptermeta{10}

\begin{ledetext}
股票天天上头条。债券没人聊。可全世界所有资产的价格，起点都是从这儿抄的。
\end{ledetext}

\section{8000 块的借条}

李敏 38 岁，在南京一家医疗器械公司做出纳。

去年同事老邵装修差钱，找她借 8000。她没好意思提利息，让他写了张条。

条上四样东西，少一样这事就说不清：还多少（8000）、什么时候还（2027 年 3 月 1 日）、中间给不给利息（不给）、谁欠的（老邵）。

把这四样标准化、编上号、允许转手，它就叫债券。市场上管它们叫面值、期限、票息、发行人。

关键在"允许转手"。

李敏借出去的 8000 是死的，她得等到明年三月。可要是这张条能卖给隔壁工位的王姐，她随时能把钱拿回来，这 8000 就活了。

要能卖，条子就得长得一模一样。中国的债券面值统一定成 100 元一张，同一批发行的每一张都是同样的票息、同一个到期日。个性化的借条没法交易，标准化的借条才有市场。

全世界规模最大的市场，就是从这么一张纸长出来的。它每天成交以万亿美元计，买卖发生在机构之间的终端和电话里，没有涨跌停板，没有人在饭桌上聊它。

无聊，但要命。

\begin{statrow}{4}
  \statitem{353}{万亿美元}{全球债务总额，2026 年 3 月底（IIF）}
  \statitem{305}{\%}{相当于全球 GDP 的比例}
  \statitem{40}{万亿美元}{美国国债规模，2026 年 8 月突破}
  \statitem{30}{万亿美元+}{其中可流通的市场规模}
\end{statrow}

\section{李敏亏了 1130 块}

2021 年秋天，理财经理跟她说债基比定存强一点，稳。她挑了一只中长期纯债基金的 C 类份额，销售页面上标着"风险等级 R2"，放进去 4 万 7。

2022 年 11 月的一个早上，她在 3 号线上打开 App。

绿的。亏了 1130。

她的第一反应是：不是说好到期还本付息吗，这钱亏在哪儿了？

亏在她提前看了一眼价格。

\begin{egbox}{举个栗子}
一张面值 100 元、票息 3\%、还有 1 年到期的债券。到期那天它会付给你 103 元，本金加利息，白纸黑字，一个字都不会改。

现在市场变了，新发行的同类债券票息给到 4\%。你手里这张老债还想按 100 元出手？没人接。人家花 100 块买新债，一年后拿回 104；买你的老债，一年后只有 103。凭什么给你同样的价。

那老债该卖多少？倒着算。买家要的是 4\% 的回报，他愿意出的价格 P 得满足 P × 1.04 = 103，于是 P = 103 ÷ 1.04 = \textbf{99.04 元}。

票息没变，本金没变，发行人也没变。变的是别人那边的价码。

\end{egbox}

\textbf{价格和收益率是同一件事的两种说法，永远反向。}收益率上去，价格必然下来。这不是什么市场规律，这是一个除法。你没办法反对一个除法。

1 年期只跌了 9 毛 6，因为李敏只需要忍受一年的票息偏低。要是这张债还有 10 年到期，她得忍十年，补偿就得给够。

\begin{tablewrap}
\begin{xltabular}{\linewidth}{>{\raggedright\arraybackslash}p{0.20\linewidth}Xrr}
\toprule
\bthead{剩余期限} & \bthead{利率 3\% 时的价格} & \bthead{利率升到 4\% 后} & \bthead{跌幅} \\
\midrule
\textbf{1 年} & 100.00 元 & 99.04 元 & −0.96\% \\
\textbf{3 年} & 100.00 元 & 97.22 元 & −2.78\% \\
\textbf{5 年} & 100.00 元 & 95.55 元 & −4.45\% \\
\textbf{10 年} & 100.00 元 & 91.89 元 & −8.11\% \\
\textbf{30 年} & 100.00 元 & 82.71 元 & −17.29\% \\
\bottomrule
\end{xltabular}
\end{tablewrap}

同样一个百分点，30 年期的损失是 1 年期的十八倍。

\begin{pullquote}{}
债券本身没有变，是折算它的那把尺子变了。
\end{pullquote}

李敏买的那只基金持有的是一篮子中长期信用债。2022 年秋天资金面收紧，理财赎回一波接一波，价格一路往下走。

她没踩雷，也没碰上任何一笔违约。她只是买了一个每天都会被重新报价的东西，然后在最难看的那天早上打开了手机。

卖给她的时候没人讲过这一段。这不怪李敏，是话术的问题。

\section{硅谷银行是怎么死的}

\term{久期}{duration}\index{031@久期}可以粗糙地理解成：你平均要等多久，才能把钱收回来。

一张 10 年期的零息债，所有钱都在第 10 年到手，久期就是 10 年。一张每年付息的 10 年期债，中间陆续有钱回来，平均等待时间短一些，久期通常在 8 年出头。

它真正的用处是当尺子：\textbf{久期 8 年，市场利率每升 1 个百分点，价格大约跌 8\%。}久期 2 年就跌 2\%。粗暴，但误差不大，全世界的债券交易员每天靠这一条心算。

2023 年 3 月 10 日，硅谷银行被接管。从股价崩盘到关门，用了两天。

它买的债，本金一分钱都没损失。

\begin{egbox}{举个栗子：一家死于算术的银行}
2019 年底，硅谷银行的存款是 618 亿美元。两年后，1892 亿。它的客户是科技公司，那两年科技公司融资融到手软，钱全趴在它这儿。

钱来得太快，贷款放不出去那么多。它拿去买债：长期美国国债和抵押贷款支持证券，票息一个多点，久期普遍五年以上，光记在"持有至到期"科目下的就有 913 亿美元。

买的时候，联邦基金利率贴着零。到 2023 年初，美联储已经把它抬到 4.5\% 以上。

拿久期这把尺子量一下：4 个多百分点乘上五六年，账面浮亏就是二三十个百分点。2022 年底那本"持有至到期"的账，成本和市价之间差了 151 亿美元。

而硅谷银行全部的股东权益，163 亿。

会计规则允许这笔浮亏不进利润表，条件是你真能拿到到期。它拿不到。科技公司在烧钱，存款一天天往外走，银行必须卖债换现金。卖一笔，浮亏落地一笔。

市场看懂了这本账。48 小时里跑掉了几百亿存款。

\end{egbox}

\begin{warnbox}{小心}
银行给这件事起了个名字，叫期限错配。

把银行想成一个水池：进水管随时能被拧掉（存款当天就能取走），出水管却焊死在五年后（长期债券）。平时相安无事，因为不会所有人同一天来接水。这门生意的全部技术含量，就是赌"不会所有人同时来"。

这个比喻可以继续往下推。什么情况下大家会同时来？消息传得够快、取款够方便的时候。硅谷银行的储户是一个互相都认识的圈子，几个投资人在群里说一句"先把钱转出来"，手机银行三秒到账。两个条件同时满足了。

存款保险为什么能防挤兑？因为它把"我得赶在别人前面跑"这个念头从你脑子里拿掉。而硅谷银行的存款，超过九成在保险上限之外。

\end{warnbox}

\section{"无风险"的意思很窄}

国债收益率被叫做\term{无风险利率}{risk-free rate}\index{067@无风险利率}。

主权政府用本币借钱，几乎不会赖账。它有征税权，实在不行还能让央行下场接盘。所以它的\textbf{名义违约风险}极低。

"无风险"三个字的全部含义，到这儿就没了。

它不豁免通胀。你锁定 3\% 的票息，第二年通胀跑到 4.2\%，没有任何人赖你的账，但你每年悄悄穷 1.2 个点。

它更不豁免利率。李敏和硅谷银行手里拿的都是不会违约的债，一个亏掉 1130 块，一个亏掉了整家银行。

所以"债券保本"这句话是错的。不是不够严谨，是错的。

\begin{mythbox}{常见误解}
\mvfrow{误解}{"债券是保本的，买债基总不至于亏吧。"}
\mvfrow{事实}{2022 年是全球债券市场几十年来最惨的一年，美国长期国债那一年跌了 30\% 以上，比同期股票还难看。持有到期确实能拿回本金和票息，前提是发行人不违约。可\textbf{债券基金没有到期日}，它每天按市价估值、随时申赎。李敏买的不是一份合同，是一段价格。把"到期还本"和"随时能卖"混着讲，是这行的老毛病。}
\mvfrow{误解}{"这国国债收益率有 8\%，说明它信用好、给得起。"}
\mvfrow{事实}{恰恰可能相反。收益率是市场\emph{要求}的补偿，要得越狠，说明市场越不放心。2010 到 2012 年希腊 10 年期国债收益率一度冲到 30\% 以上，那不是希腊信用变好了。看一国的债得同时看两个数：收益率的绝对水平，和它相对于同期限最安全那条曲线的\textbf{利差}。}
\end{mythbox}

\section{打分的人，买单的人}

不是所有借条都由政府写。公司也发债，中国还有一类很特殊的城投债。

评级公司给每一张债打分。AAA 的意思接近"基本不会赖账"，BBB− 以上统称投资级，BB+ 往下叫高收益债，行里的叫法更实在：垃圾债。CCC 的含义大致是"这家公司现在就有麻烦"。

评级不是道德评价，是概率翻译。投资级债券的长期累计违约率在个位数，投机级能到两位数。你多要的那几个点利差，就是给这个概率开的价：\textbf{信用利差 = 你要的收益率 − 同期限国债收益率}。

\begin{databox}{一个十几年没改动的毛病}
给评级机构付钱的，是发债的公司，不是买债的人。你交学费，你考试，你还能挑老师。2008 年那批被打上 AAA、后来血本无归的次贷证券，就是这个模式的产物。

危机之后各国加了不少监管，付费结构基本没动。我的判断是它也动不了：让买方付费听起来很对，可你花钱买来的评级，转身就被市场上所有人免费看走了，没人愿意当那个掏钱的冤大头。这道题的名字叫搭便车，经济学解了一百年也没解开。

\end{databox}

城投债是中国债市的特产。地方政府过去通过融资平台公司举债搞基建，这些债法律上是企业债，市场却一直按"政府会不会兜底"来定价。这几年在做的事叫债务化解：把高息、短期、藏着的债，换成低息、长期、摆在明面上的债。利息负担确实降了，本金还在那儿，只是换了张脸和一个新的还款日。

这么多债，最后都堆在谁手里？

\begin{flowlist}
  \flowitem{01}{保险公司和养老金}{它们欠的是几十年后要付出去的赔款和养老金，久期长得吓人。要匹配就必须买长债。利率下行的时候它们最难受，因为负债的现值涨得比资产还快。}
  \flowitem{02}{银行}{拿存款去买债，赚一个利差，顺便满足流动性监管指标。硅谷银行干的就是这件事，只是干过了头。}
  \flowitem{03}{央行}{不为赚钱。它买卖国债是为了把政策意图传出去，价格是手段，不是目的。}
  \flowitem{04}{基金和 ETF}{替普通人打包，每天按市价估值。李敏买的是这一类，所以她能天天看见浮亏。}
  \flowitem{05}{外国投资者}{两个动机：储备要安全，利差要赚。这两个动机的相对分量正在变，下一节说。}
\end{flowlist}

这堆买家按行为分两类。保险、养老金、银行的自营账户是\textbf{配置盘}，买了就放着，只关心收益够不够覆盖负债成本。对冲基金、券商自营、一部分公募是\textbf{交易盘}，只关心三个月内价格往哪儿走。

债市剧烈波动的时候，动的一般不是配置盘，是交易盘在互相踩脚。配置盘的钱慢，交易盘的钱快，而价格永远由最快的那批钱定。

\section{第一把交椅换人了}

欧洲央行 2026 年 6 月的《欧元的国际角色》报告里，有一组数字。

\begin{barfig}
  \barrow{黄金（2024 年底）}{20}{20\%}
  \barrow{黄金（2025 年底）}{27}{27\%}
  \barrow{美国国债（2024 年底）}{25}{25\%}
  \barrow{美国国债（2025 年底）}{22}{22\%}
  \barrow{其他美元资产（2025 年底）}{20}{20\%}
  \barrow{欧元（2025 年底）}{15}{15\%}
\end{barfig}

截至 2025 年底，黄金占全球央行储备的 27\%，美国国债占 22\%。\textbf{黄金第一次超过美债，坐上了各国央行的第一把交椅。}一年之前，这个对比还是黄金 20\%、美债 25\%。

美元计价资产合计仍占 42\%，还是最大的单一币种阵营。所以这不是塌方，是配比在挪。为什么挪，第 12 章讲，那是一个关于"钱会不会被没收"的故事。

\begin{databox}{2026 数据}
IIF 在报告里警告，美国政府债务"日益走上不可持续的道路"。同一时期，外国投资者对美债的净买入总体保持稳定，但同步加配了日本和欧洲的主权债。这是分散化，不是撤离。

2025 年，欧元计价的国际债券发行量增长 30\%，接近 1 万亿欧元的历史新高。2026 年 8 月底，美国 10 年期国债收益率升到 2025 年 1 月以来的最高水平。

\end{databox}

\section{凭什么}

你手里任何一样东西的价格，本质上都是未来的现金流除以一个折现率。第 6 章推过这个式子。

那个折现率的第一块砖，就是同期限的国债收益率。它是所有定价的海拔基准。

\begin{chainflow}
\chainnode{国债收益率}\chainarrow \chainnode{信用债定价}\chainarrow \chainnode{股票估值的折现率}\chainarrow \chainnode{房贷与企业贷款定价}\chainarrow \chainnode{跨境资本流向}\chainarrow \chainnode{李敏的基金净值}
\end{chainflow}

10 年期美债收益率动 50 个基点，纳斯达克的估值模型要重算，上海写字楼的资本化率要重算，新兴市场的资金流向要重算。李敏家的房贷定价基准，也吊在同一条链子上。

\begin{talkbox}
  \talkline{新闻}{"截至 8 月底，美国 10 年期国债收益率升至 2025 年 1 月以来最高水平。"}
  \talkline{李敏}{"这跟我有什么关系。"}
  \talkline{翻译}{全世界给"未来的钱"打折时用的那个除数变大了 → 靠远期现金流撑估值的东西一起变便宜 → 已经发出去的长久期债券价格往下走 → 保险公司和养老金的账面被重估 → 借美元的新兴市场成本抬高 → 三个月后，你那只债基的净值曲线上会多出一个不太好看的弯。}
\end{talkbox}

\begin{pullquote}{}
债市打个喷嚏，全世界都得感冒。区别只在于，交易员当天就知道自己感冒了，李敏要等三个月才从基金净值上看出来。
\end{pullquote}

这本书不给任何投资建议。但有一件事可以说死：债券基金会亏钱，这是它的机制决定的，不是它出了差错。

\begin{takeaway}{本章一句话}
债券是一张标准化、可转手的借条。它的价格和收益率永远反向，而全世界所有资产的定价，都从国债收益率这块砖起步。

\begin{itemize}
  \item 价格与收益率反向是一个除法，不是市场规律。期限越长，同样的利率变动砸出来的坑越深：1 个百分点，30 年期跌 17.29\%。
  \item 久期是这个坑的深度尺。久期 8 年，利率升 1 个点，价格约跌 8\%。硅谷银行 151 亿美元的浮亏对 163 亿美元的股东权益，死因写在这把尺子上。
  \item "无风险"只指名义违约风险，不管通胀，不管利率。
  \item 截至 2025 年底，黄金以 27\% 超过美债的 22\%，成为全球央行第一大储备资产；美元计价资产合计仍占 42\%。
\end{itemize}

\end{takeaway}

\begin{bridgebox}
\textbf{下一章}债券是"先说好还你多少"。可要是我不想要那个说死的数字，我想跟着一家公司一起赚、一起赔呢？那就得先把公司切开。
\end{bridgebox}


\chapter{股票：把一家公司切成一亿片}
\chaptermeta{11}

\begin{ledetext}
买一股，到底买到了什么？不是一张纸，也不是一个代码，是排在所有人后面的那一份剩饭。
\end{ledetext}

\section{八月那张单子，钱是排队发的}

陈晓婷 32 岁，在杭州拱墅区开了家叫"三分甜"的奶茶店。八月结账，收银系统吐出来的营业额是 12 万 6400。

这笔钱不是她的。至少现在还不是。

\begin{flowlist}
  \flowitem{01}{供应商先拿}{珍珠、茶叶、鲜奶、杯子和封口膜，这个月进货 4 万 1200。不结账，下个月不发货。}
  \flowitem{02}{然后是四个店员}{工资加社保 3 万 4700。晚发三天，暑假工第二天就不来了。}
  \flowitem{03}{房东}{1 万 9600，押一付三，晚交按天算滞纳金。}
  \flowitem{04}{水电、外卖平台抽成、杂七杂八}{1 万 800。这一项每年都在往上爬，她已经不太看了。}
  \flowitem{05}{银行}{2019 年装修借的 42 万经营贷，这个月还息 1347。不还，抵押的那套小两居就不是她的了。}
  \flowitem{06}{税务局}{增值税、附加、企业所得税。有小微优惠，但不是不交。}
  \flowitem{07}{最后轮到陈晓婷}{1 万 7800 出头。去年八月对面开了家新店，这个数字是负的。}
\end{flowlist}

这个顺序不是行业习惯，是法律。公司清算的时候，职工工资和社保、税款、有担保的债权、普通债权，一层一层排下去，股东永远在最后一位。

法律给这个位置起了个名字：\term{剩余索取权}{residual claim}\index{048@剩余索取权}。买一股股票，买到的就是这个。

排最后有两个后果，是同一件事的两种长相。

坏的那面：前头任何一个人多拿一点，陈晓婷就少拿一点。营业额掉两成，房租工资利息一分不能少，那两成全从她这一份里扣。

好的那面：上不封顶。银行拿到 1347 就走了，供应商拿到货款就走了。哪个月这条街突然爆火，多出来的钱一分都不用分给他们。

\begin{pullquote}{}
股票风险高和股票长期回报高，不是两个事实，是同一个位置的两种读法。
\end{pullquote}

\section{王磊那 1800 块，三分甜一分钱没拿到}

三分甜后来开到了 14 家店，2025 年合并净利润 100 万，其中有三家还在亏钱。

这个体量当然不够上市。这本书里让它上市，是因为它的数字你能心算。

\begin{egbox}{举个栗子}
挂牌那天，三分甜发行 25 万股新股，每股 15 元，募到 375 万。这 375 万真金白银打进公司账户，用来在滨江开五家新店。这是\textbf{一级市场}，公司确实拿到了钱。

王磊在合肥一家汽车零部件厂做工艺工程师，月薪 1 万 3 出头。2021 年他拿了 3 万 8 的年终奖，抽出 1 万 2 开了证券账户。三分甜上市第二天开盘价 18 元，他觉得这生意他看得懂，买了 100 股，付出去 1800 块。

这 1800 块给了谁？给了把股票卖给他的那个人。三分甜的账户余额和昨天收盘时一模一样，一分没多。这是\textbf{二级市场}。

\end{egbox}

王磊以为自己在给公司投资。他其实是在跟另一个股民换手。

那二级市场是不是白长的？恰恰相反，只是它的作用全都隔了一层。

它给三分甜的未来标价。股价 18 元时再发 20 万股能拿到 360 万；跌到 6 元，同样 20 万股只能拿 120 万。融资成本明码写在屏幕上，谁都能看见。

它还是并购的钱。三分甜想吃掉隔壁的"茶小满"，可以不掏现金，直接拿自己的股票换。它也决定了店长手里那批行权价 10 元的期权值不值钱：股价 18 元，一股赚 8 块；股价 6 元，那就是一张彩色废纸，人留不住。

最要紧的是最后一条：它给一级市场开了个出口。

2019 年投三分甜第一笔钱的那位天使投资人，他的钱只有在能卖掉的那天才算落袋。没有二级市场，就不会有人做早期投资。风险投资、创业公司、期权激励，整条链子都吊在"将来能卖掉"这一个预期上。

把二级市场关掉，一级市场三个月内枯死。

\section{PE 是懒人版的贴现公式}

第 6 章说过，任何资产的价值，等于它未来所有现金流按某个折现率贴回今天。这个式子数学上无懈可击，实践中几乎没法用：要预测未来二十年的利润，还要挑一个折现率，手抖一下答案差三倍。

市场于是发明了一个偷懒的办法。

\begin{egbox}{举个栗子}
三分甜净利润 100 万，总股本 100 万股，每股收益 EPS = 1 元。王磊买入价 15 元，\term{市盈率}{P/E ratio}\index{052@市盈率} = 15 ÷ 1 = \textbf{15 倍}。

15 倍的字面意思：假设利润永远不变，公司要替他赚 15 年，才够把买入价还回来。

倒过来看更有味道。1 ÷ 15 = \textbf{6.7\%}，这是他的盈利收益率。同期 5 年期以上 LPR 是 3.5\%（2026 年 8 月），10 年期国债收益率更低。6.7\% 对 3.5\%，多出来那 3.2 个百分点，就是市场给"奶茶可能过气、店可能关门、陈晓婷可能不想干了"这一堆事开的价钱。

\end{egbox}

PE 把那个复杂的贴现公式压成了一个数字，代价是所有关于未来的假设全被塞进了分母里。它从来不是"高就贵、低就便宜"。

成长股 PE 高，买的是未来。一家利润每年涨 30\% 的公司，今年 PE 50 倍，三年后利润翻一倍多，股价不动的话 PE 只剩 22 倍。前提是那个未来真的来了。没来，50 倍就是 50 倍。

周期股 PE 低的时候，才是真要命的时候。

\begin{warnbox}{小心}
一家造集装箱的公司，行业景气顶点那年赚了 78 亿，股价对应 PE 只有 4.6 倍，看着比谁都便宜。

第二年运价崩了，利润掉到 8 亿多。股价一动没动，PE 立刻变成 45 倍。

\textbf{周期股 PE 最低的那一刻，通常正是利润最高的那一刻，也就是最危险的那一刻。}PE 的分母是当期利润，当期利润越不可持续，这个比值越骗人。这条规律翻过的车，能从上海排到杭州。

\end{warnbox}

\section{ROE 高不高，不如问它把钱放哪儿了}

ROE 是净利润除以股东权益，衡量股东投进去的每一块钱一年能赚回几毛。三分甜的净资产 500 万，净利 100 万，ROE 20\%。

它和增长之间有个朴素到近乎粗鲁的关系：

\textbf{可持续增长率 ≈ ROE ×（1 − 分红率）}

三分甜把 32\% 的利润分掉，留下 68\% 继续开店。它靠自己能长多快？20\% × 0.68 = 13.6\%。就这么多，除非去借钱或者增发。

所以"高 ROE"这三个字单独拿出来没有意义，得配一个问句：钱还有地方投吗？

三分甜每开一家新店都能赚回 20\%，那它就该少分红、多开店，这是一台复利机器在转。

可要是杭州的奶茶店已经饱和，新开一家只能赚 4\%，它还把利润死死留着扩张，等于把本来能赚 20\% 的钱拿去做 4\% 的生意。报表上写着"营收增长"，股东的钱在这个过程里被磨掉了。这种时候把钱还给股东，才是负责任的做法。

还钱有两种姿势。分红是直接打到王磊账上。回购是公司拿钱在市场上买回自己的股票然后注销，王磊手里的股数一股没变，但蛋糕被切成的片数少了，他那片自动变大。

\begin{egbox}{举个栗子}
三分甜净利 100 万，100 万股，股价 15 元，EPS 1 元。

它掏出 150 万现金，在市场上回购 10 万股，注销。股本变成 90 万股，利润还是 100 万，EPS 变成 100 万 ÷ 90 万 = \textbf{1.11 元}，涨了 11\%。

没多卖一杯奶茶，没多赚一分钱，EPS 涨了 11\%。

\end{egbox}

回购有两副面孔。股价低于内在价值时回购，等于替留下来的股东做了一笔划算买卖；股价高得离谱还借钱回购，那是拿股东的钱在高位接盘，顺便给 EPS 化个妆，好让管理层自己的期权顺利兑现。

税上回购确实占便宜。分红当期就要缴税，回购把收益留在股价里，不卖就不用交。A 股的股息税按持股期限差别化：持股超过 1 年暂免征收，1 个月以内全额计入应纳税所得额按 20\% 计征，中间那段减半。这条规则本身就在推着你拿久一点。

\section{称重机得有人真的去称}

格雷厄姆那句话被引用了八十年：市场短期是投票机，长期是称重机。

这句话有个从来没人明说的前提。\textbf{得有人真的在称。}

2025 到 2026 年，全球被动投资的规模第一次超过了主动管理。指数基金和 ETF 不研究公司，它们只干一件事：按指数权重买。某只股票被纳入指数，就必须买，不管贵不贵；被剔除，就必须卖，不管便宜不便宜。

王磊 2023 年也把大部分钱换成了沪深 300 指数基金，管理费从 1.5\% 降到 0.15\%，他觉得自己占了大便宜。他确实占了。

可如果所有人都这么想，谁来决定一家公司值多少钱？

\begin{deepbox}{进阶：可以跳过}
\textbf{一方说没问题。}定价是在边际上完成的。哪怕只剩 10\% 的资金还在做主动研究，只要它们能自由买卖，价格照样会被推回合理位置。指数基金按现有价格等比例买入，并不改变相对价格。何况被动化把管理费从 1.5\% 打到 0.05\%，这是实打实还给王磊们的钱。这一派背后是 Grossman 和 Stiglitz 那套老逻辑：只要主动管理还能赚到超额收益，就总有人来做，市场会自动维持一个刚好够用的研究量。

\textbf{另一方说有问题。}指数纳入和剔除本身就在制造跟基本面无关的价格波动，早被人抢跑套利。市值加权意味着涨得越多的股票权重越大，钱越往头部堆，流动性和估值一起挤在少数几家公司身上。下跌的时候被动资金的赎回同样机械，卖出不看估值，可能把踩踏放大。

\textbf{我的看法：}短期内价格发现不会失灵，主动资金的绝对规模还是太大了。但"钱按权重涌向头部"确实在改变市场结构，指数集中度处在历史高位是事实。至于被动化的比例到了多少会真正伤到价格发现，这件事目前没人知道答案，谁跟你说他知道，大概率是在卖产品。

\end{deepbox}

\section{王磊五年交了 4700 块过路费}

股市里的钱只有三个来源，性质完全不同。

第一种是企业盈利增长。三分甜今年比去年多赚了 40 万，所有股东按比例一起分，不需要任何人亏钱。这是股市长期回报唯一靠得住的来源。

第二种是估值变化。PE 从 15 涨到 20，公司一分钱没多赚，多出来那部分是后面接手的人付的。王磊赚的，正好是他多付的。

第三种是交易成本和税费。佣金、印花税、买卖价差、管理费。这部分钱谁也没赚走，它直接离开了池子。

王磊从 2021 年到 2026 年，账户从 1 万 2 加到 9 万 4，五年里买进卖出六十多次。佣金加印花税，一共交了 4700 多。

这 4700 块和大盘涨跌无关。涨也交，跌也交。

所以一群人在存量里高频倒手，整体一定是亏的。不是因为他们笨，是因为第三项永远带负号。第 20 章会把这个等式完整拆开。

\section{同一套打法，两个市场，两个结果}

\begin{tablewrap}
\begin{xltabular}{\linewidth}{>{\raggedright\arraybackslash}p{0.20\linewidth}XX}
\toprule
\bthead{维度} & \bthead{A 股} & \bthead{美股} \\
\midrule
\textbf{投资者构成} & 个人投资者贡献的成交额占比长期偏高 & 机构与被动资金主导，个人直接持股比例低得多 \\
\textbf{交易制度} & T+1；主板涨跌停 ±10\%，创业板与科创板 ±20\% & T+0，无涨跌幅限制，设有市场级熔断 \\
\textbf{退市} & 退市新规后节奏明显加快，历史存量仍在消化 & 常态化，每年大量公司因不达标或私有化退出 \\
\textbf{股息率} & 近年在政策推动下持续抬升，整体仍低于成熟市场 & 分红加回购的综合现金回报率更高 \\
\textbf{指数编制} & 上证综指覆盖全部上市股票，长期采用总股本加权 & 标普 500 由委员会筛选，自由流通市值加权，有盈利门槛 \\
\textbf{监管取向} & 融资功能与投资者保护并重，政策工具使用频繁 & 侧重信息披露与事后追责，集体诉讼是主要威慑 \\
\bottomrule
\end{xltabular}
\end{tablewrap}

王磊有个大学同学在纽约做量化，两人聊过好几次，发现同一套因子搬过去就失灵。

散户成交占比高的市场，情绪波动更大，短期反转更明显；退市不畅的市场，壳本身有价格，垃圾股会被反复炒；T+1 加涨跌停，会把今天没释放完的抛压原封不动推到明天开盘。

\begin{databox}{2026 数据}
2026 年 8 月 31 日收盘：上证指数 3986.30 点，当月涨 4.02\%；深证成指 14015 点；创业板指 3438.68 点；两市成交约 2.13 万亿元。中国公募基金规模约 38 万亿元（2026 年初）。日经 225 在 2026 年 6 月首次站上 70000 点。

多家机构判断，市场的核心逻辑正在从"估值修复"转向"盈利驱动"。这两个词对应的正好是前面那笔账：估值修复赚的是零和的那一块，盈利驱动赚的是正和的那一块。

\end{databox}

\begin{mythbox}{常见误解}
\mvfrow{误解}{"股价跌了一半，这公司肯定亏惨了。"}
\mvfrow{事实}{股价是股东之间转手的价钱，它根本不出现在公司的利润表里。三分甜股价从 18 跌到 9，它照样卖奶茶，照样赚它的 100 万。真正受影响的是间接那几条：再融资更贵、店长手里的期权变废纸、大股东质押可能爆仓、换股并购不划算。\textbf{公司亏钱是经营的事，股价下跌是定价的事。两者会互相拉扯，但不是一回事。}}
\mvfrow{误解}{"10 块钱的股票比 100 块的便宜。"}
\mvfrow{事实}{这句话是错的，错得很干净。股价是"每一片多少钱"，估值是"整个蛋糕多少钱"。A 公司切成 1 亿片，每片 10 元，市值 10 亿；B 公司切成 100 万片，每片 100 元，市值 1 亿。贵的那个反而便宜十倍。只看股价不看市值和盈利，等于问"这块蛋糕切一刀多少钱"。}
\mvfrow{误解}{"我买股票就是在支持实体经济。"}
\mvfrow{事实}{王磊那 1800 块给的是上一个股东，三分甜账上一分没多。真正把钱交到公司手里的，是参与 IPO 认购、参与定向增发、买公司新发行的债。\hlmark{但二级市场的存在，让公司能以更低成本增发、让创业者和早期投资人有地方退出、让期权有价值。}他不是在直接输血，他是在维护那套让输血成为可能的管道。这个作用是真的，只是隔了一层。}
\end{mythbox}

这本书不推荐任何标的，也不判断任何价格。上面所有的算术都是为了让你看懂这台机器怎么转。

\begin{takeaway}{本章一句话}
股票是剩余索取权：所有人拿完之后剩下的那份归你，所以它风险最大，也因此长期回报最高。

\begin{itemize}
  \item PE、PB 是贴现公式的简写，不是独立真理。周期股 PE 最低时往往最危险。
  \item 可持续增长率 ≈ ROE ×（1 − 分红率）。三分甜留下 68\% 的利润，它靠自己能长 13.6\%。
  \item 股市的钱只有三笔：盈利增长（正和）、估值变化（零和）、成本税费（负和）。王磊五年在第三笔上花掉 4700 块。
  \item 被动投资规模在 2025 到 2026 年首次超过主动管理。它会不会削弱价格发现，目前没人知道答案。
\end{itemize}

\end{takeaway}

\begin{bridgebox}
\textbf{下一章}股票和债券都用本国货币标价。可王磊要是想买德国的机器、日本的旅游、美国的芯片，他手里这串数字得先换成别人账本上的数字才算数。
\end{bridgebox}


\chapter{汇率与美元体系：为什么全世界用别人的钱做生意}
\chaptermeta{12}

\begin{ledetext}
一个宁波老板卖货给巴西人，两边都不是美国人，合同却用美元签。这背后是一套没人投票选出来的秩序。
\end{ledetext}

\section{一块钱值多少钱，得问另一块钱}

老韩 2011 年从一家台资厂出来单干，在宁波北仑做汽车密封条，厂里现在 63 个人。2024 年他接了圣保罗一家改装车厂的单子，132 万美元，账期 90 天。

报价那天美元兑人民币 7.24，钱到账那天 7.08。同样 132 万美元，少收了 21 万人民币。

他那个月的利润，就在这 16 个基点上没了。

汇率不是一个国家的考试成绩，它是两种钱互相的报价。"7.1"的意思是 1 美元换 7.1 元人民币，同一个数倒过来读，就是 1 元人民币换 0.141 美元。两个方向，谁也不比谁高级。

谁在决定这个数？

老韩的第一反应是贸易：中国出口多，外国人得先买人民币才能付货款，人民币就该涨。这个逻辑本身没错，可它解释不了他看到的行情。

全球外汇市场日均成交是万亿美元量级。而全球货物贸易一整年的总额，摊到每个工作日也就几百亿美元的量级。

\textbf{在外汇市场里，做真实生意的钱是少数派。}

绝大部分成交来自资本流动：跨境买卖股票债券、跨境放贷、对冲敞口、套利、投机。它们盘子大、掉头快、几乎没有摩擦成本。老韩那 132 万美元排队等着结汇的时候，隔壁有人几秒钟就调走了十亿。

所以看汇率，先看钱往哪儿跑，再看货往哪儿走。一个国家完全可能贸易顺差很大而货币在贬，也完全可能常年逆差而货币很硬。美国已经这么硬了几十年。

\begin{chainflow}
\chainnode{美国加息}\chainarrow \chainnode{美元资产更香}\chainarrow \chainnode{资本流入美国}\chainarrow \chainnode{美元走强}\chainarrow \chainnode{新兴市场货币承压}\chainarrow \chainnode{他们的美元债还起来更贵}
\end{chainflow}

\section{巨无霸能告诉你方向，不能告诉你时间}

《经济学人》1986 年发明了巨无霸指数，逻辑简单到可爱：同一个汉堡，在哪儿都该值一样多的钱。

假设一个巨无霸在美国卖 5.8 美元，在中国卖 25 元，"汉堡汇率"就是 25 ÷ 5.8 ≈ 4.3。要是市场汇率在 7 以上，按汉堡算，人民币被低估了四成。

这是\term{购买力平价}{purchasing power parity}\index{021@购买力平价}的直觉：长期看，汇率该让同一篮子东西在各国卖出差不多的价。

麻烦在于，汉堡运不过去。

一个巨无霸里，牛肉和面包是可以跨境运的，可门店租金、店员工资、水电、当地税全都运不走。上海店员的工资和芝加哥店员的工资差多少，汉堡价格就差多少，套利者没办法把上海的店员打包空运到芝加哥去。

这个漏洞可以一直漏下去。加上关税、增值税、饮食习惯，PPP 算出来的偏差能常年存在，一存几十年。

它是一根很长的橡皮筋，不是一张时刻表。谁拿它判断"下个季度该升还是该贬"，等于拿望远镜看牙齿。

\section{借日元，买美债：全世界最拥挤的那笔生意}

\begin{egbox}{举个栗子}
2026 年 8 月，同一个地球上，钱的价格差出好几倍。日本央行政策利率 \textbf{1.0\%}，2026 年 6 月刚从 0.75\% 加上来，是 1995 年 9 月以来最高；美联储联邦基金利率目标区间 \textbf{3.50\%–3.75\%}；中国 1 年期 LPR \textbf{3.0\%}。

于是有人做这样一笔生意：在东京借日元，成本大约 1\%；换成美元，买美国短期国债，收益大约 3.6\%。什么活都不用干，一年白赚 2.6 个百分点。

2.6\% 不解渴，那就上杠杆。5 倍，变 13\%。10 倍，变 26\%。

\end{egbox}

这就是\term{套息交易}{carry trade}\index{058@套息交易}，全球资金流动里最庞大也最不起眼的一股力量。它的规模没人能精确统计，因为它藏在无数张互不相干的资产负债表里。

它的毛病只有一句话：\textbf{利差是按天赚的，汇率是按小时还的。}

\begin{egbox}{举个栗子：2024 年 8 月那三天}
2024 年 7 月 31 日，日本央行把政策利率从 0～0.1\% 上调到 0.25\%，同时缩减购债。加息 25 个基点，听起来像挠痒痒。

但成本端动了。日元开始升值，美元兑日元从 7 月初的 161 一路走到 8 月初的 142 附近，不到一个月，日元强了十几个百分点。

借日元的人要还钱，得用更贵的日元买回来。攒了好几年的那点利差，被汇率一口吞掉。

平仓引发平仓。8 月 5 日，日经 225 单日跌 12.4\%，创下它历史上最大的单日点数跌幅；美股同步大跌，波动率指数盘中冲到 60 以上。从东京传到纽约，用了几个小时。

起因是日本一次 25 个基点的加息。

\end{egbox}

\begin{pullquote}{}
套息交易的收益曲线像爬楼梯，风险曲线像坐电梯。三年爬上去，三天掉下来。
\end{pullquote}

\section{三选二：所有国家都在这张桌子上做同一道题}

\term{三元悖论}{impossible trinity}\index{045@三元悖论}说的是：固定汇率、资本自由流动、独立货币政策，这三样你最多同时拿两样。

为什么？假设你既要汇率钉死，又要资本自由进出。美联储把利率抬到 3.75\%，你的利率是 1\%，资本立刻往美国跑，你的货币面临贬值压力。可你承诺了固定汇率，就只剩两条路：拿外汇储备去买回本币，或者把自己的利率也抬到 3.75\%。前者会把储备耗光，后者等于承认你的货币政策是美国定的。

三样都想要的，历史上没有一个撑住。1992 年英镑被迫退出欧洲汇率机制，1997 年泰铢崩盘，做错的都是同一道题。

\begin{tablewrap}
\begin{xltabular}{\linewidth}{>{\raggedright\arraybackslash}p{0.20\linewidth}XXX}
\toprule
\bthead{经济体} & \bthead{放弃了什么} & \bthead{换来了什么} & \bthead{代价长什么样} \\
\midrule
\textbf{香港} & 独立货币政策 & 港元钉住美元在 7.75–7.85 区间，资本完全自由进出 & 美联储加息，香港必须跟着加，哪怕本地楼市和零售正在挨打 \\
\textbf{欧元区成员国} & 独立货币政策 & 成员国之间等于彻底固定汇率（用同一种钱），资本自由流动 & 2010–2012 年主权债务危机时，希腊没法靠贬值恢复竞争力，只能靠内部通缩，也就是直接砍工资 \\
\textbf{美国、日本} & 固定汇率 & 完全独立的货币政策，资本自由流动 & 汇率大幅波动。2024 年 8 月日元三周走强十几个百分点，全球市场跟着地震 \\
\textbf{中国} & 资本完全自由流动 & 独立的货币政策，有管理的浮动汇率 & 跨境资金进出有额度和审批，老韩换汇要提交合同和报关单 \\
\bottomrule
\end{xltabular}
\end{tablewrap}

香港金管局跟着美联储走不是没主见，是合同条款；欧洲央行一套利率管二十个国家，必然有人被夹疼，这是设计使然；中国能在美联储加息的时候降息，靠的就是老韩填的那些表格。

\section{老韩为什么不用人民币报价}

他试过。巴西客户第一反应是"那我的银行怎么给我报价"。

\begin{flowlist}
  \flowitem{01}{网络效应}{老韩的橡胶原料按美元报价，海运费按美元结算，货运险和信用证也是美元。他一个人换币种，等于把自己从整张网里摘出来。}
  \flowitem{02}{有个足够深的池子停钱}{收到美元不能压箱底，得有地方放。美国国债可流通市场规模超过 30 万亿美元，今天卖 10 亿也砸不动价格。这种深度全世界只有这一家。}
  \flowitem{03}{随时能换出来}{这一条比"收益高"重要得多。储备货币的第一要求是想走的时候走得掉。}
  \flowitem{04}{法治与合约的可预期性}{钱放在纽约，出了纠纷有一套大致能预判的司法程序。这条听着虚，定价的时候非常实。}
  \flowitem{05}{清算管道}{SWIFT 报文加上美元清算体系，是跨境支付的默认路径。用了半个世纪，改造成本高到没人愿意先动。}
\end{flowlist}

\begin{deepbox}{进阶：可以跳过}
1960 年，比利时经济学家罗伯特·特里芬指出了一个绕不过去的矛盾。

全世界都想用美元做储备和结算，那美元就必须源源不断流到美国之外去。怎么送出去？靠贸易逆差。美国进口比出口多，把美元付给全世界。

可长期逆差意味着美国持续对外欠债。欠得越多，别人越怀疑这些美元将来还值不值这个价。

\textbf{世界需要美元的方式，恰恰会削弱美元。}这就是特里芬难题，它没有解，只有拖。布雷顿森林体系 1971 年就是被这道题拖垮的，当时的表现形式是美国的黄金储备不够兑付境外的美元。今天它换了张脸：美国国债在 2026 年 8 月突破 40 万亿美元，IIF 警告这条路"日益走上不可持续的道路"。

\end{deepbox}

\section{2026：储备资产的第一属性变了}

欧洲央行 2026 年 6 月的《欧元的国际角色》报告，给了一组数字。

\begin{barfig}
  \barrow{黄金}{27}{27\%}
  \barrow{美国国债}{22}{22\%}
  \barrow{其他美元资产}{20}{20\%}
  \barrow{欧元}{15}{15\%}
\end{barfig}

截至 2025 年底，黄金占全球央行储备的 27\%，美国国债占 22\%。\textbf{黄金第一次超过美债，成为各国央行的第一大储备资产。}2024 年底这个对比还是黄金 20\%、美债 25\%，其他美元资产从 22\% 降到 20\%，欧元稳在 15\%。

美元计价资产合计仍占 42\%，依然是最大的单一币种阵营。这不是崩塌，是配比在挪。

\begin{statrow}{3}
  \statitem{3.6}{万吨+}{全球央行黄金储备，接近布雷顿森林体系鼎盛期的约 3.8 万吨}
  \statitem{850}{吨}{2025 年全球央行净购金}
  \statitem{5600}{美元/盎司}{2026 年 1 月金价一度接近的水平}
\end{statrow}

为什么突然开始抢金子？

分水岭是 2022 年。俄乌冲突之后，西方国家冻结了俄罗斯央行数千亿美元规模的外汇储备。

在那之前，各国央行管钱的考核表上写着三项：安全、流动性、收益。"会不会被没收"这一项根本不在表上。

那天之后，所有非西方阵营的储备管理者都做了同一道算术：我持有的美债，本质上是别人账本上的一行数字。那行数字能不能被划掉，不取决于我。

黄金没有这个问题。它是唯一一种不构成任何人负债的储备资产。搬进自己的金库，它就在那儿，不需要谁替你记账。代价是它不付利息，还得花钱看守，一吨金子一年的保管费不是小数目。

\begin{pullquote}{}
储备资产的第一属性，从"收益"变回了"不会被没收"。
\end{pullquote}

\begin{databox}{2026 数据：一个有点荒诞的细节}
2025 年全球最大的单一黄金买家不是哪国央行，是稳定币发行商 Tether，一年买进超过 100 吨。

一家专门给美元稳定币做储备的公司，拿赚来的钱去买金条。连最"美元原生"的那批玩家，也在给自己的储备垫一层不依赖任何账本的东西。

\end{databox}

\section{人民币走到哪一步了}

跨境贸易结算里人民币的占比这些年在稳步上升，CIPS 的参与机构和覆盖国家持续扩大，数字人民币的跨境支付网络 CBETS 在 2026 年新增了 26 家银行。老韩现在有两个南美客户愿意用人民币结算，三年前一个都没有。

这些是实打实的进展。但真正的约束不在管道上。

\textbf{第一，资本项目还没有完全可兑换。}别人愿意拿一种货币当储备的前提，是想走的时候能走。这条和三元悖论直接顶上：中国选了独立货币政策加有管理的汇率，就必须留着资本管制这道闸。要放开，就得在另外两样里让出一样。这是个真实的取舍，不是意愿问题。

\textbf{第二，国债市场的深度和开放度。}储备货币需要一个巨大的、外国人能自由进出的、买卖不怎么影响价格的安全资产池。中国国债市场的规模已经不小，但外资持有比例和交易便利程度，跟美债还不在一个量级。

什么时候能实现资本项目基本可兑换？这件事目前没人知道答案，包括正在做这个决定的人。它取决于国内金融体系的稳健程度、外部环境、以及那一年的风险偏好。

\begin{mythbox}{常见误解}
\mvfrow{误解}{"本币贬值就是坏事，升值说明国力强。"}
\mvfrow{事实}{得分三条线看。\textbf{出口企业：}人民币贬值，老韩那 132 万美元换回更多人民币，报价也更有竞争力，他是受益方。\textbf{进口方和你我：}贬值让原油、芯片、留学、出国旅游全部变贵，输入型通胀跟着来。\textbf{有外币负债的：}这条最凶险。老韩 2023 年买德国挤出机借了 260 万美元设备贷，但他七成收入是美元，这叫自然对冲，扛得住。扛不住的是那种只赚人民币却借了美元债的公司：人民币贬 10\%，收入一分没变，债务折成人民币凭空多出一大块。1997 年亚洲金融危机里倒下的那批企业，死因高度一致，行话叫\term{币种错配}{currency mismatch}\index{005@币种错配}。}
\mvfrow{误解}{"美元马上要崩溃了" / "人民币很快就会取代美元"。}
\mvfrow{事实}{这两句话都是错的，而且错在同一个地方：把趋势的方向当成了趋势的速度。截至 2025 年底，美元计价资产仍占全球央行储备的 42\%，是最大的单一阵营；同时黄金一年之内从 20\% 升到 27\%，2025 年欧元计价国际债券发行增长 30\%、接近 1 万亿欧元的历史新高，国际投资者净流入欧元区资产 8500 亿欧元。\hlmark{这叫渐进多元化，不叫替代。}储备货币的更替以几十年计：英镑让位给美元花了大半个世纪，中间还隔着两次世界大战。}
\end{mythbox}

本章不预测任何货币的走势，也不给任何换汇建议。老韩的锁汇决定是他自己做的，他做对过，也做错过。

\begin{takeaway}{本章一句话}
汇率是两种钱的相对价格，今天主要由资本流动而不是贸易决定；美元体系靠的是网络效应、市场深度和清算管道，而 2026 年正在发生的是渐进多元化。

\begin{itemize}
  \item 购买力平价是长期引力，不是短期指南。它告诉你方向，不告诉你时间。
  \item 套息交易赚的是利差，还的是汇率。2024 年 8 月日经单日跌 12.4\%，起因是日本加息 25 个基点。
  \item 三元悖论三选二：香港和欧元区成员国放弃独立货币政策，美日放弃固定汇率，中国放弃资本完全自由流动。
  \item 截至 2025 年底黄金占全球央行储备 27\%，首次超过美债的 22\%；美元计价资产合计仍占 42\%。
\end{itemize}

\end{takeaway}

\begin{bridgebox}
\textbf{下一章}到这里，管道全部铺好了：钱怎么造出来、怎么定价、怎么跨境流动。可管道里流的东西，为什么隔几年就要炸一次？下一层我们抬头看天气。
\end{bridgebox}



\part{天气篇：潮汐、风暴与泡沫}
\chapter{风险的价格：为什么有人借钱 3\%，有人 12\%}
\chaptermeta{13}

\begin{ledetext}
管道已经铺好，钱能流到任何地方。但流过去要收多少过路费，是另一门手艺。这门手艺一半是算术，一半是人心。
\end{ledetext}

\section{郑姐那天下午问了四个地方}

郑姐在常州开一家做汽车密封条的小厂，二十三个工人，2025 年出货一千一百多万。2026 年 8 月的一个下午，她要凑 80 万买一台新的注塑机。旧的那台漏油漏了大半年，修一次两万三，一年修了四回。

她问了四个地方，拿回四个数字。

\begin{tablewrap}
\begin{xltabular}{\linewidth}{>{\raggedright\arraybackslash}p{0.20\linewidth}XXr}
\toprule
\bthead{路子} & \bthead{拿什么做保} & \bthead{条件} & \bthead{年化} \\
\midrule
\textbf{房贷（她家那套 2019 年买的）} & 房子本身 & 30 年，等额本息 & 3.5\% \\
\textbf{抵押经营贷} & 住宅二次抵押 + 担保公司兜底 & 3 年，先息后本 & 3.9\% \\
\textbf{信用卡分期} & 什么都不用押 & 12 期，月费率 0.6\% & 13.03\% \\
\textbf{某持牌消费金融的小额经营贷} & 只看征信和流水 & 12 期，随借随还 & 23.9\% \\
\textbf{（对照）美国的私募信贷} & 企业资产 + 一堆财务契约 & 5 到 7 年，浮动 & 8\%–12\% \\
\bottomrule
\end{xltabular}
\end{tablewrap}

房贷那一行的 3.5\%，是 2026 年 8 月 20 日公布的 5 年期以上 LPR，已经连续多月没动。最后一行是给你看看，同一个世界的另一头，钱是什么价。

同一个人，同一天，同一笔 80 万。最贵的比最便宜的贵了将近七倍。

这不是谁比谁善良。

\begin{egbox}{先把那个 13.03\% 算出来}
客户经理跟郑姐说的是：月费率只有 0.6\%，一年下来 7.2\%，比经营贷贵一点点，好处是当天到账。

7.2\% 是假的。

她先试了 12 万，分 12 期。每月还本金 1 万，手续费按\textbf{最初的 12 万}算，120000 × 0.6\% = 720 元，所以每月还 10720。全年手续费 8640 元，8640 ÷ 120000 = 7.2\%，账面上没错。

错在分母。第一个月她手上确实压着 12 万，第二个月只剩 11 万，到第十二个月只剩 1 万。整段时间里她平均占用的钱是 6.5 万左右，不是 12 万。

拿 8640 除以 65000，13.3\%。严格解出内部收益率，是 \hlmark{13.03\%}。

银行说的每个字都对，它只是没提分母在缩水。中国从 2021 年起要求消费信贷明示年化利率，所以现在很多 App 上会有一行小字写着"年化利率（单利）"。看那一行。

\end{egbox}

\section{一笔利率是四块钱叠起来的}

任何一个利率数字，撕开都是同样几层。

\begin{flowlist}
  \flowitem{01}{无风险利率}{把钱借给几乎不可能赖账的对手，同样的期限，能拿多少。中国这边的锚是国债和政策利率，2026 年 8 月 7 天期逆回购利率是 1.4\%。这一块跟你是谁没关系，只跟时间有关。}
  \flowitem{02}{\term{期限溢价}{term premium}\index{042@期限溢价}}{借 3 年比借 7 天贵。放贷的人把钱锁死三年，这三年里通胀怎么走、政策怎么变、他自己会不会急用，全都没得反悔。这份不自由要收钱。}
  \flowitem{03}{\term{信用利差}{credit spread}\index{072@信用利差}}{借款人可能还不上。这是四块里最大、也最难算的一块，下一节整节都在算它。}
  \flowitem{04}{\term{流动性溢价}{liquidity premium}\index{035@流动性溢价}}{如果放贷方半路急需用钱，这笔债权能不能出手、要打几折。上市公司的债几分钟就能卖掉，郑姐这笔贷款卖给谁？}
  \flowitem{05}{税和条款}{抵押、担保、提前还款罚金、财务契约、要不要开一个受监管的资金账户。条款每收紧一格，利率就能松一格。这是可以谈的部分，也是小企业主最少去谈的部分。}
\end{flowlist}

教科书就分到这里为止。银行内部的定价表长得不太一样，它更实用，也更冷酷。

下面这个拆法是示意，不是哪家银行的真实报价单。假设这笔小微经营贷最后定在 8.1\%：

\begin{barfig}
  \barrow{资金成本}{25}{2.0\%}
  \barrow{预期损失}{30}{2.4\%}
  \barrow{运营成本}{27}{2.2\%}
  \barrow{资本回报}{15}{1.2\%}
  \barrow{税与杂项}{3}{0.3\%}
\end{barfig}

最长的那根不是银行的利润，是预期损失。

\section{2.4 个百分点，是九十六个人替四个人付的}

信贷定价的核心只有一个公式，写出来三个字母：

\begin{pullquote}{}
预期损失 = 违约概率 × 违约损失率 × 敞口
\end{pullquote}

这个式子朴素到有点无聊，但它决定了这个星球上每一笔贷款的价格。

\begin{egbox}{给一百家小微企业放贷}
一家城商行放 100 笔，每笔 80 万，总敞口 8000 万。

一年下来，4 家还不上。这是\term{违约概率}{probability of default}\index{063@违约概率}，4\%。

违约不等于全赔光。抵押品处置、诉讼、执行，平均两年半之后，每家收回 40\%。损失掉的是 60\%，这是\term{违约损失率}{loss given default}\index{064@违约损失率}。

算总账：4 笔 × 80 万 × 60\% = \textbf{192 万}。摊到 8000 万的盘子上，2.4\%。

也就是说，这家银行光要把本金亏掉的部分挣回来，利率就得比无风险利率高出 2.4 个点，一分钱不赚。

再加上资金成本 2.0\%、运营成本 2.2\%、股东要的资本回报 1.2\%，8\% 出头。\hlmark{一分利润没多拿。}

\end{egbox}

运营成本那 2.2\% 是最容易被忽略的一块。一笔 80 万的小微贷，客户经理要上门、要看厂房、要核订单、要查上下游，放款之后每季度还得跟踪，出问题要催收。这一整套动作，和一笔 8000 万的公司贷款做的事差不多。

人力成本几乎一样，分母差一百倍。

所以"小微企业融资难融资贵"这句话，前半段的第一层解释不是道德，是四家企业还不上，剩下九十六家在替他们付钱，外加一百套尽调的工时。

\section{但数学解释不了全部}

如果只有数学，郑姐的价格就该是 8\%，不该有 3.9\% 那一行。

3.9\% 是怎么来的？她把住宅做了二次抵押，还找了本地一家政府背景的融资担保公司。这两件事同时干了两件活：抵押品把违约损失率从 60\% 压到 25\% 左右，担保把银行实际承担的违约概率又切掉一大半。预期损失从 2.4\% 掉到不足 0.6\%，利率自然就下来了。

征信数据干的是另一半活。银行能合法看到她的增值税发票流水、社保缴纳、用电量、上下游的应收账款，那个 4\% 的违约概率就不再是一个行业平均数，而是她这一家的数。\hlwarm{风险定价贵，很多时候不是因为风险真的大，是因为看不清。}看不清的部分要按最坏的算。

我的判断是这样的：中国小微融资贵这件事，数学占大头，摩擦占的那一小半是可以往下压的，而且这些年确实压下去了一截。压它的从来不是文件里的道德呼吁，是税务数据、电力数据、供应链单据能不能被银行合法看见、看见之后能不能被采信。数据确权往前走一步，利率就往下挪一小格。这条路很慢，走了十几年才挪了这么点，但它是真的路。

\begin{mythbox}{常见误解}
\mvfrow{误解}{"利率收这么高，这家机构就是黑心。"}
\mvfrow{事实}{先把 8\% 和 24\% 分开看。8\% 里面绝大部分是预期损失和运营成本，是算出来的。24\% 那一档，预期损失可能真的有十几个点，因为来借的人本身就是被前面几道门槛筛剩下的。\\ 真正该骂的不是数字大，是\textbf{把 13.03\% 说成 7.2\%}，是把砍头息、服务费、保险费拆成七八个名目让你算不清总价。价格高不是罪，价格不透明才是。}
\end{mythbox}

\section{风险、不确定性，和一个消不掉的东西}

1921 年，芝加哥大学的弗兰克·奈特写了一本书，把两样长得很像的东西分开了。

骰子是\textbf{风险}：六个面，每一面的概率你都知道，你算不准下一把，但你算得准一千把。明天开门的那家新店能不能活过三年是\textbf{不确定性}：你不知道有几个面，甚至不知道它是不是个骰子。

金融业干的事，是把不确定性硬塞进风险的模子里。办法是找历史数据来充当那六个面。

数据有多长，你的骰子就有几个面。2007 年以前，美国全国房价从来没有同时下跌过，模型里根本不存在这一面。不是模型算错了，是那一面在骰子上没刻。

\begin{warnbox}{小心：VaR 那句话的后半句}
风控报表上最常见的数字叫\term{在险价值}{Value at Risk}\index{081@在险价值}。一句标准表述是："99\% 的置信度下，本组合单日亏损不超过 3200 万。"

这句话听着让人安心，因为它只说了 99\% 的日子。剩下那 1\% 的日子亏多少？\textbf{VaR 不回答。}它的定义里就不包含那部分。

更麻烦的是，那 1\% 的日子有个坏习惯：它们喜欢连着来。一年 250 个交易日里的两三天，往往不是散落在三月、七月和十一月，而是挤在同一个星期里。

\end{warnbox}

说完 1\% 的日子，说说另一件被误会得更深的事。

把钱摊到很多篮子里，能消掉的叫非系统性风险：这家厂子着火、那家公司的 CFO 跑路、某个专利官司输了。这些事互相不认识，一多就互相抵消。

消不掉的叫系统性风险：利率、汇率、战争、政策、全社会同时不敢花钱。它作用在所有篮子上。

决定你到底分散了多少的，只有一个变量：相关性。

\begin{egbox}{同样 20 只股票，三种命运}
假设你持有 20 只股票，每只 5\%，每只的年化波动都是 30\%。

如果它们两两之间完全不相关（相关系数 0），组合的波动会降到单只的 22.4\%，也就是 6.7\%。这就是分散化被称为"免费午餐"的原因，你没牺牲任何期望收益，风险少了四分之三。

把两两相关系数调到 0.6，同样这 20 只，组合波动只能降到单只的 \textbf{79\%}。

调到 0.9，降到 \textbf{95\%}。持有 20 只和持有 1 只，几乎没有区别。

2008 年 10 月和 2020 年 3 月，很多平时相关系数只有 0.2 的东西，在同一周里一起冲到了 0.9 以上。股票、公司债、新兴市场货币、黄金、REITs，全跌。

\end{egbox}

因为在恐慌里，所有人卖的不是某个资产，是"任何能卖掉的东西"。这时候资产之间的基本面关系失效了，剩下的只有一个共同因子：谁需要现金。

\begin{pullquote}{}
危机中所有相关性都趋近于 1。你精心搭的那个组合，恰好在最需要它的那一天不再是组合。
\end{pullquote}

\section{为什么好人会先离场}

1970 年，阿克洛夫写了一篇被三家期刊退稿的论文，讲二手车市场。买家分不清好车和坏车，只肯出一个平均价；好车车主觉得亏，不卖了；市场上剩下的车更差；买家再压价。最后好车全部退场，只剩柠檬。

把二手车换成借款人，这个循环一模一样。

\begin{chainflow}
\chainnode{分不清好坏}\chainarrow \chainnode{只肯按平均定价}\chainarrow \chainnode{好借款人嫌贵走人}\chainarrow \chainnode{池子里的人更差}\chainarrow \chainnode{只好再涨价}\chainarrow \chainnode{更好的人再走}
\end{chainflow}

这叫\term{逆向选择}{adverse selection}\index{040@逆向选择}。它发生在借钱之前。

借到钱之后还有第二个问题。郑姐拿这 80 万去买注塑机，银行放心。可如果她拿去做了别的，比如投给她表弟那个据说很赚钱的项目，银行是不知道的。钱一旦到手，用钱的人和出钱的人，利益就分岔了。这叫\term{道德风险}{moral hazard}\index{010@道德风险}。

金融业几百年来发明的那些看着很烦的东西，全是在治这两个病。抵押品治的是逆向选择（敢押房子的人对自己有信心），也治道德风险（乱来会失去房子）。分期放款治道德风险：钱按工程进度给，你没法一次性拿走。财务契约治的也是它：负债率超过某个数，银行有权提前收贷。征信记录治的是逆向选择，它把"你上次是什么样"变成了这次的价格。

评级机构本该是治逆向选择的最大一件工具。它的工作就是替所有人分辨好车和柠檬。

它的收费结构是：\textbf{被评的那家公司付钱。}

这个安排的问题不需要解释。2008 年，一大批贴着 AAA 标签的结构化产品变成了废纸，那是史上最贵的一次演示。危机之后确实改了不少：内部防火墙、方法论公开、监管备案、发行人不能挑评级师。但付钱的还是被评的人。

所以我对评级的态度是：当一份有用的、带着已知偏向的参考读，不当结论用。一家公司被降级往往说明市场早就知道了，评级只是最后一个承认的。

\section{2026 年：利差是最早的那支体温计}

\begin{statrow}{3}
  \statitem{2.3}{万亿美元}{全球私募信贷管理规模（2025 年）}
  \statitem{8–12}{\%}{私募信贷典型借款利率}
  \statitem{200–400}{bp}{比银行贷款高出的部分}
\end{statrow}

\begin{databox}{2026 数据}
私募信贷 2025 年全球管理规模约 2.3 万亿美元，2020 年以来翻了一倍，年化增速超过 17\%。借款利率通常在 8\% 到 12\%，比银行贷款高 200 到 400 个基点；二顺位和夹层这类更靠后的层级，12\% 到 16\%。

它的行业敞口里软件业约占 30\%，其中投向 AI 领域的约 2000 亿美元，占总资产 9\% 左右。

2025 年 10 月，美国两家中小银行 Zions Bancorporation 和 Western Alliance 先后出现风险事件。中金把私募信贷称作"2 万亿美元的灰犀牛"，认为风险可能在未来 12 到 24 个月逐步传导。这是一家机构的判断，不是已发生的事实。

\end{databox}

为什么这些利差数字值得盯着？因为信用市场的报价比股市诚实。

\begin{talkbox}
  \talkline{新闻}{"某中型企业的债券收益率较同期限国债的利差，从 320 个基点走阔至 610 个基点。"}
  \talkline{翻译}{三个月前，把钱借给这家公司比借给政府多赚 3.2 个点；现在要多赚 6.1 个点才有人肯借。市场对它违约概率的估计，翻了将近一倍。它的股价可能还在原地，因为股东在做梦，债权人不做梦。}
\end{talkbox}

股东买的是这家公司变成十倍的可能性，所以坏消息可以被"故事"抵消一阵子。债权人最多拿回本金加利息，上不封顶的东西跟他没关系，他只关心一件事：这家公司会不会死。

所以坏消息在信用市场里先有价格。2007 年 2 月，次级抵押贷款相关指数已经在往下掉，标普 500 还在往上爬，一直爬到 10 月才见顶。八个月的时间差。

\begin{mythbox}{常见误解}
\mvfrow{误解}{"高风险高收益。我年轻，能承受风险，所以我该去拿高收益。"}
\mvfrow{事实}{风险溢价是一个\textbf{期望值}，不是一份合同。它的完整意思是"平均而言、长期来看、在足够多的独立样本上"，多出来那几个点会兑现。你这一次不是样本，你这一次是一次抽样。\\ 更要命的是另一半：市场只补偿\emph{没法被分散掉的}那部分风险。你把全部身家压在一只股票上，承担了比市场高得多的波动，但其中大部分是可以被分散掉的非系统性风险，理论上一分钱溢价都不该给你。这叫承担了不被补偿的风险。\\ 还有第三种：某些资产的高利率不是给你的风险补偿，是给渠道的销售费用。这不构成投资建议，只是提醒你去看费率表。}
\end{mythbox}

\bookbreak

郑姐最后选了 3.9\% 那条路，把房子押了。她跟我说了一句话，我一直记着："押房子我心里踏实，因为我知道我肯定能还上。要是我自己都不敢押，那我凭什么让银行信我。"

这句话把这一章的道理讲完了。

\begin{takeaway}{本章一句话}
利率不是一个数字，是一份拆得开的账单：无风险利率 + 期限溢价 + 信用利差 + 流动性溢价，加上一堆条款和摩擦。

\begin{itemize}
  \item 预期损失 = 违约概率 × 违约损失率 × 敞口。小微贷款利率高，一大半是这道算术的结果。
  \item 另一小半是摩擦，可以靠抵押品、征信数据和担保压下去，但压不到零。
  \item 风险是可测的，不确定性不是。VaR 只告诉你 99\% 的日子，麻烦全在剩下那 1\%。
  \item 分散化的效果由相关性决定，而危机里所有相关性都往 1 跑。
  \item 信用利差是这台机器上最早报警的那支温度计，比股市早。
\end{itemize}

\end{takeaway}

\begin{bridgebox}
\textbf{下一章}定价错了，未必致命。真正把小麻烦变成灾难的，是另一样东西。它就藏在你家的房贷合同里，而且大多数人从没意识到自己正用着它。
\end{bridgebox}


\chapter{杠杆：金融世界的放大镜与绞肉机}
\chaptermeta{14}

\begin{ledetext}
同样的判断，同样的方向，同样的市场。有人赚了 40\%，有人归零。中间隔着的那个东西，你家客厅的天花板下面就有一份。
\end{ledetext}

\section{96 万，两种活法}

许工在合肥一家做汽车电子的公司当硬件工程师，太太在医院药房。2026 年，两个人凑出 96 万，包括双方父母出的一部分。

他们面前有两条路。

一条是全款买一套 96 万的小两居，从此没有月供，晚上睡得着。另一条是拿 96 万当首付，贷 224 万，买一套 320 万的房子，位置好，学区还行。

第二条路的杠杆倍数是 320 ÷ 96 = \textbf{3.33 倍}。这个数字接下来会把所有事情乘以 3.33。

\begin{tablewrap}
\begin{xltabular}{\linewidth}{>{\raggedright\arraybackslash}p{0.20\linewidth}rr}
\toprule
\thead{~} & \thead{全款买 96 万} & \thead{首付 96 万 · 贷 224 万 · 买 320 万} \\
\midrule
\textbf{自己的钱} & 96 万 & 96 万 \\
\textbf{手上握着的资产} & 96 万 & 320 万 \\
\textbf{房价涨 12\% 后的净值} & 107.5 万（+12\%） & 134.4 万（+40\%） \\
\textbf{房价跌 12\% 后的净值} & 84.5 万（-12\%） & 57.6 万（-40\%） \\
\textbf{房价跌 30\% 后的净值} & 67.2 万（-30\%） & 0 \\
\bottomrule
\end{xltabular}
\end{tablewrap}

看最后一行。

房子跌到 224 万，正好等于贷款余额。卖掉，还清银行，一分不剩。中介费和税还得自己再掏。

整个过程里，许工的判断没有错过一次，市场也没有针对他。只是他把一个 30\% 的价格波动，通过 3.33 倍的杠杆，翻译成了 100\% 的净值消失。

\begin{egbox}{还没算利息}
那 224 万不是白借的。按 2026 年 8 月 5 年期以上 LPR 3.5\% 算，30 年等额本息：

月供 \textbf{10059 元}。三十年总共还回去 362 万，其中利息 \hlmark{138 万}。

换句话说，这套 320 万的房子，许工实际要付出 96 万 + 362 万 = 458 万。房价三十年里不涨 43\%，他就是亏的。

这不是说买房不该贷款。这是说，杠杆的成本不在首付那一刻，在后面那 360 个月里，每个月 10059 块，风雨无阻。

\end{egbox}

\begin{pullquote}{}
你的房贷就是杠杆。而且对绝大多数中国家庭来说，它是这辈子唯一一笔、也是最大的一笔。
\end{pullquote}

\section{银行的 ROA 只有 0.8\%，凭什么给股东 10\%}

杠杆的正式定义很朴素：用别人的钱，控制比自己的钱更多的资产。

它对股东回报的影响可以拆成一个乘法，这个拆法一百年前是杜邦公司的会计发明的：

\begin{chainflow}
\chainnode{ROE 股东回报}\chainarrow \chainnode{ROA 资产回报}\chainarrow \chainnode{权益乘数}
\end{chainflow}

\term{权益乘数}{equity multiplier}\index{043@权益乘数}就是总资产除以自有资本。你手上有 1 块钱，撑起 12 块钱的资产，乘数就是 12。

看一家中等规模的银行长什么样：

\begin{bsheet}
  \bsside{资产 ASSETS}{%
    \bsrow{各类贷款}{682 亿}
    \bsrow{债券投资}{214 亿}
    \bsrow{现金与存放央行}{104 亿}
    \bsrowtotal{合计}{1000 亿}
  }
  \bsside{负债 + 净值}{%
    \bsrow{存款}{771 亿}
    \bsrow{同业拆借与发债}{147 亿}
    \bsrow{股东的钱}{82 亿}
    \bsrowtotal{合计}{1000 亿}
  }
\end{bsheet}

权益乘数 1000 ÷ 82 = 12.2 倍。

这家银行一年在 1000 亿资产上赚 8 亿，资产回报率 0.8\%，放在任何一个制造业老板眼里都是笑话。但对股东来说，8 亿 ÷ 82 亿 = \textbf{9.8\%}，是个体面的数字。

同一个乘数，反过来读一遍：\hlwarm{资产端只要出 82 亿的窟窿，也就是资产总额的 8.2\%，股东的钱就一分不剩。}

银行这门生意的结构就是这样。它不是因为管理不善才脆弱，它生下来就脆弱。全世界的银行监管、存款保险、最后贷款人，整套东西都是围着这 8.2\% 建的。

\begin{databox}{2026 数据}
中国商业银行净息差 2026 年二季度末为 \textbf{1.41\%}，处在历史低位。净息差被压薄，直接压的是 ROA 那一项。

ROA 掉下去，ROE 要维持住，数学上只有两条路：把权益乘数加上去，或者把成本降下去。而权益乘数是被巴塞尔协议卡死的。

我的判断是，2026 年中国银行业最真实的压力不在坏账，在息差。坏账是周期问题，息差是结构问题。

\end{databox}

\section{杀死你的不是下跌，是那通电话}

杠杆本身不杀人。

许工那套房如果跌了 30\% 他不卖，只要月供还得上，银行不会来敲门。中国的房贷极少设置盯市追加条款。他可以趴在那里等十年。

真正的绞肉机是另一套机制：\term{保证金追缴}{margin call}\index{002@保证金追缴}。

\begin{flowlist}
  \flowitem{01}{价格跌了一点}{你的仓位市值缩水，但欠的钱一分没少。}
  \flowitem{02}{担保比例跌破线}{券商或交易对手每天盯市，比例低于合同约定的数就触发条款。这是自动的，不需要任何人做判断。}
  \flowitem{03}{电话来了}{要么当天补钱，要么减仓。留给你的时间通常是一个交易日，有时候是几小时。}
  \flowitem{04}{你补不上}{能补的人早补了。被追保的时候，你身上其他资产大概率也在跌。}
  \flowitem{05}{强制平仓}{系统按市价卖，不看基本面，不等反弹，不听解释。}
  \flowitem{06}{卖压把价格再压低}{而这个更低的价格，是所有人共用的那个价格。}
  \flowitem{07}{回到 01}{另一批人的担保比例跌破线。}
\end{flowlist}

这个圈叫\term{流动性螺旋}{liquidity spiral}\index{034@流动性螺旋}。它一旦转起来，跟谁对谁错没有关系。

\begin{pullquote}{}
在这个循环里，你即使判断完全正确也会死。因为你活不到被证明正确的那一天。
\end{pullquote}

\begin{mythbox}{常见误解}
\mvfrow{误解}{"加杠杆没问题，我提前想好止损位就行了。"}
\mvfrow{事实}{止损是一份需要别人配合才能兑现的合同。你挂一个卖单，得有人接。\\ 2015 年 7 月上旬，A 股连着几天上千只股票跌停，很多人的止损单老老实实挂在那里，一天，两天，三天，成交不了。跌停板上是没有对手方的。\\ 2020 年 3 月，美股十天内四次熔断，很多流动性极好的 ETF 出现了远低于净值的成交。你以为你有一个出口，其实那扇门在最需要它的时候是从外面锁着的。\\ \textbf{止损管理的是你的纪律，管不了市场的流动性。}把仓位控制在"就算完全卖不掉也死不了"的水平，比设止损位有用得多。}
\end{mythbox}

\section{你没觉得自己借了钱的那些杠杆}

杠杆很少写着"杠杆"两个字。

\textbf{\term{期限错配}{maturity mismatch}\index{041@期限错配}。}用 7 天的钱去买 5 年的债，赚中间那段期限利差。只要每 7 天都能续上，这生意能一直做下去。整套模式的全部重量，压在"每 7 天都能续上"这一件事上。2013 年 6 月中国银行间市场的钱荒、2019 年 9 月美国回购市场那一夜隔夜利率冲到 10\%，都演过续不上是什么样子。

\textbf{\term{回购}{repo}\index{023@回购}。}拿债券做抵押借隔夜钱，借来的钱再买债，新债再拿去抵押，再借。抵押品打个 2\% 的折扣，理论上能滚出 50 倍。实践中滚不到那么高，但把自有资金放大八到十倍，在债券市场是常规操作。

\textbf{结构化产品的分层。}一个池子切成优先级和劣后级，劣后级用自己一份钱替优先级的四份钱兜底，收益也归它。它一分钱没借，但它的风险敞口是自有资金的五倍。

\textbf{期权。}一张合约对应 100 股或者一手标的，权利金只有标的价值的百分之几。买期权本身就是一种付了保险费的杠杆。

\textbf{经营杠杆。}这一种连金融都不算。一家影院，租金、折旧、最低人力是固定的，上座率从 45\% 掉到 35\%，收入掉两成，利润可能直接翻到负数。它没有任何借款，但它的成本结构就是杠杆。2020 年倒掉的那一批线下生意，很多死在这上面，不是死在负债表上。

监管机构花了很多年才认清一件事：这些形态各异的杠杆没法一个一个去堵，堵完一个还会长出下一个。只能从背面管。管不住你用多少种方式放大敞口，就管你自己垫了多少钱。

2008 年之后的巴塞尔 III，核心就三个数：资本充足率、杠杆率、流动性覆盖率。

第一个说的是你自己得垫多少钱。第二个是个不看风险权重的兜底约束，防止你把资产都塞进"看起来风险很低"的篮子里。第三个说的是：真出事的头 30 天，你手上得有足够立刻能变现的东西。

整套逻辑一句话说得完：

\begin{pullquote}{}
监管不是不让你冒险。监管是要求你冒险之前，自己先垫钱。
\end{pullquote}

为什么是资本不是存款？因为存款是别人的钱，你亏了要还；资本是股东的钱，亏了就亏了。资本这一层的厚度，决定了在损失砸到储户和纳税人身上之前，有多少缓冲。

2008 年之前，一些欧美投行的权益乘数做到 30 倍以上。3.3\% 的资产损失就能抹平全部资本。今天的大银行普遍在 10 到 15 倍。这个变化，是过去十七年金融监管做成的最实在的一件事。

\section{三倍做多，为什么长期一定亏}

这一节几乎没有科普讲清楚过，但它是纯算术，五分钟能懂。

杠杆 ETF 承诺的是\textbf{每天}跟踪标的涨跌的三倍，不是一个月三倍，不是一年三倍。为了做到每天三倍，它必须每天收盘后调仓：涨了就加仓，跌了就减仓。

这个"每天重置"的动作，在震荡市里会持续吃掉你的钱。

\begin{egbox}{两天，标的回到原点，你亏 5.45\%}

\begin{tablewrapin}
\begin{tabularx}{\linewidth}{>{\raggedright\arraybackslash}p{0.20\linewidth}rr}
\toprule
\thead{~} & \thead{标的指数} & \thead{三倍做多产品} \\
\midrule
\textbf{起点} & 100.00 & 1.0000 \\
\textbf{第一天 +10\%} & 110.00 & 1.3000 \\
\textbf{第二天 -9.09\%} & 100.00 & 0.9455 \\
\textbf{两天合计} & 一分没动 & -5.45\% \\
\bottomrule
\end{tabularx}
\end{tablewrapin}

第二天的 -9.09\% 是把 110 打回 100 需要的跌幅。三倍就是 -27.27\%，1.3 × 0.7273 = \textbf{0.9455}。

标的回到了出发点，你少了 5.45\%。

再来一遍，0.9455 × 0.9455 = 0.894。来十个这样的来回：\hlmark{0.571}。

标的指数在原地上下晃了十趟，一分钱没涨没跌。你的三倍产品少了 \textbf{42.9\%}。

\end{egbox}

这个损耗有个名字，叫\term{波动损耗}{volatility decay}\index{007@波动损耗}。它跟管理费无关，跟基金经理水平无关，它是每日重置这个机制的数学后果，躲不掉。

标的走得越颠簸，损耗越大。单边行情里三倍产品能跑出超过三倍的收益，震荡行情里它稳定失血。

所以杠杆 ETF 的说明书上那句"适合有短期交易需求的投资者"，不是免责套话，是产品的物理属性。拿三倍产品做长期持有，你赌的不是方向，是波动率。这不构成投资建议，这是算术。

\section{三个坟头}

\textbf{1998 年，LTCM。}合伙人名单里有默顿和斯科尔斯，两个人 1997 年刚拿了诺贝尔经济学奖。他们做的是相对价值套利，赚的是两个几乎一样的债券之间那点微小价差，每一笔的利润薄得像纸，所以必须上极高的杠杆才有意义。表内杠杆二十几倍，算上衍生品的名义敞口更高。

模型说，出现 1998 年 8 月那种规模的价格偏离，概率小到几百上千年才有一次。

基金 1994 年 2 月开张。1998 年 8 月，俄罗斯违约。

四年。

\textbf{2015 年，A 股场外配资。}1 比 5 是常见的，1 比 10 也有。上证从 5178 点掉下来的时候，配资盘的平仓线写在系统里，到点就砍，不看基本面，不接受申辩。卖出压低股价，压低股价触发下一批平仓线。这就是上一节那个循环，只是跑得特别快。

\textbf{2021 年，Archegos。}Bill Hwang 用总收益互换在六七家投行分别开仓，每家投行只看得见自己那一部分。等到 ViacomCBS 增发砸了股价，第一家开始平仓，剩下几家才发现原来还有别人。两天，大约 100 亿美元的自有资金归零；瑞信在这一单上亏了 55 亿美元，是它后来一系列麻烦的开头之一。

这三件事隔了二十三年，行业不同，工具不同，剧本一模一样：\textbf{高杠杆，集中持仓，对手方看不见全貌。}

\section{2026 年的杠杆在哪儿}

今天最被盯着的那块是私募信贷。它的杠杆到底高不高，答案有点反直觉。

\begin{tablewrap}
\begin{xltabular}{\linewidth}{Xr}
\toprule
\thead{时期与主体} & \thead{典型杠杆倍数} \\
\midrule
\textbf{2007 年前后的 CDO 结构} & 10–20 倍 \\
\textbf{2007 年前后的欧美投行} & 20–30 倍 \\
\textbf{2025–2026 年的私募信贷基金} & 0–2 倍 \\
\bottomrule
\end{xltabular}
\end{tablewrap}

低得多。它的负债端也更稳定，锁定期长，不像回购那样每天都要续；层层嵌套的结构也少。这几点是中金在把它叫"灰犀牛"的同时明确指出的，跟次贷不是一回事。

真正需要看的数字在另一边。

\begin{statrow}{3}
  \statitem{5000–6000}{亿美元}{美国银行体系流向私募信贷的资金}
  \statitem{3}{\%}{占银行总资产}
  \statitem{30}{\%}{占银行资本}
\end{statrow}

3\% 听起来像个零头。30\% 就不是了。

差别在分母。资产是银行手里的东西，资本是垫在最下面、亏光就完蛋的那一层。前面那张资产负债表已经算过：资产端 8.2\% 的损失就能吃掉全部资本。

套到这里：如果这 5500 亿美元的敞口最终损失三分之一，那是 1833 亿美元，等于把整个美国银行体系的资本削掉 \textbf{10\%}。不是削掉 1\%，是 10\%。

这是杠杆这一章最该记住的读数方法。\hlwarm{看一笔风险有多大，永远除以资本，不要除以资产。}除以资产是给自己看的，除以资本才是给风险看的。

\begin{mythbox}{常见误解}
\mvfrow{误解}{"杠杆是坏东西，我这个人比较保守，我从不用杠杆。"}
\mvfrow{事实}{许工的杠杆是 3.33 倍，比大部分对冲基金的股票策略还高。\\ 房贷之所以是一种"安全"的杠杆，靠的不是倍数低，而是三件事：不盯市（跌了没人打电话）、期限极长（30 年，你有时间等）、现金流来自工资而不是资产本身（房子不涨也能还上）。这三条同时成立，3.33 倍才睡得着。\\ 只要抽掉其中任何一条，同样的倍数就会变成另一种东西。比如你还款靠的是这套房的租金，而租金正好跟房价一起跌。}
\end{mythbox}

\begin{deepbox}{进阶：可以跳过}
企业和银行也有杠杆，但读法不同。企业看净负债除以 EBITDA，4 倍以上算高；银行看权益乘数和资本充足率；基金看总敞口除以净资产。

这些指标都有一个共同的盲点：它们是快照。真实的风险是"在最坏的那一天，你的资产能卖出什么价，以及有多少债在那一天到期"。这两件事不在任何一张报表上。

\end{deepbox}

\begin{takeaway}{本章一句话}
杠杆不改变你判断的对错，它只改变你能撑多久。放大收益的那个倍数，和让你归零的那个跌幅，是同一个数的正反面。

\begin{itemize}
  \item 3.33 倍杠杆下，价格跌 30\% 你的净值就是 0。倒数关系，算得清清楚楚。
  \item ROE = ROA × 权益乘数。银行 0.8\% 的资产回报能变成 9.8\% 的股东回报，代价是 8.2\% 的资产损失就归零。
  \item 致命的从来不是杠杆，是被迫平仓。止损单在流动性枯竭时执行不了。
  \item 期限错配、回购、结构化分层、期权、经营杠杆，都是杠杆，都不叫杠杆。
  \item 三倍做多产品在震荡市里稳定失血，十个来回亏 42.9\%，标的一分没动。
  \item 衡量任何一笔风险，除以资本，不要除以资产。
\end{itemize}

\end{takeaway}

\begin{bridgebox}
\textbf{下一章}杠杆能把风险放大，另一类工具能把风险切下来卖给别人。它被骂了十七年，骂得最凶的人大多没搞清它是什么，而它最早的用途，是替一个伊利诺伊的农民保住秋天的麦子价。
\end{bridgebox}


\input{chapters/ch-15}
\chapter{影子银行：管道之外的管道}
\chaptermeta{16}

\begin{ledetext}
有一整套系统，干着和银行一模一样的活，却不穿银行的制服。它不违法。它是这台机器的正常器官，而且拆不掉。
\end{ledetext}

\section{常州那笔 4200 万}

常州武进有家做精密减速器的厂子，老板姓夏。2025 年秋天他接到一个人形机器人关节的订单，要扩一条产线，缺 4200 万。

他先去了银行。信贷员很客气，材料也收得认真：厂房是租的，专用设备做抵押折价折得狠，近三年收入复合增长 40\% 多，但净利润薄，2024 年才刚刚转正。

信贷员没说"不行"。信贷员说的是"我们再研究研究"。

这不是他人品有问题。银行放这笔款要占用资本，权重一算，对银行的资本回报不划算。他手上还有八个客户等着。

夏总最后是在一家美元基金旗下的信贷平台拿到钱的。四年期，前两年只付息，年化 9.8\%，附一堆条款：净负债对 EBITDA 不得超过 3.5 倍，每季度交一次表，超了就重新定价。

这笔钱不是任何人的存款变的。它来自一家北欧的养老金、一家新加坡的保险公司，以及几个连夏总名字都没听说过的家族办公室。

\begin{egbox}{算算他多付了多少}
如果银行肯放，以夏总的资质，报价大概在 5.9\%。他实际拿到的是 9.8\%，中间差 390 个基点。

4200 万 × 9.8\% = 411.6 万，这是他一年的利息。\\ 4200 万 × 5.9\% = 247.8 万，这是他本来可以付的。\\ 一年多掏 163.8 万，四年多掏 655 万。

夏总认了。银行给他的选项不是"更便宜的钱"，是"没有钱"。多付的 390 个基点买的是这条产线今年就能开，而不是等他把利润做厚两年。到那时候订单已经给别人了。

\end{egbox}

\section{银行的三件事，谁都能干}

把银行这门生意拆开，核心只有三件：

\textbf{期限转换。}借进来的是短的，活期存款随时能提；放出去的是长的，房贷三十年。\\ \textbf{流动性转换。}负债端随时能变现，资产端一时半会变不成现金。\\ \textbf{信用转换。}吸进来的是被当作零风险的存款，放出去的是有可能收不回的贷款。

任何机构，只要它同时干这三件事，它做的就是银行的生意。它自称什么不重要。

\term{影子银行}{shadow banking}\index{078@影子银行}的定义就是这么来的：在银行监管框架之外完成这三种转换的一整套机构和市场。

它缺的是两样东西。一是不受资本充足率、拨备覆盖率、流动性覆盖率那一整套约束；二是出事的时候没有存款保险，背后也没有央行的贴现窗口接着。

把持牌银行想成有牌照的出租车：能上专用道，出事有保险，但要装计价器、年检、交份子钱。影子银行是网约车刚出现那两年，同样把人从 A 送到 B，价格更灵活、接单更快，撞了车谁赔当时没人说得清。

这个比喻可以往下推一步：监管的思路从来不是"禁止用私家车拉人"，而是"什么时候发牌、装什么表、买哪种险"。而每一轮发牌落地之后，总会冒出一种还没被写进条例的新送人方式。这不是监管无能，这是定义问题：规则只能写已经存在的东西。

\begin{mythbox}{常见误解}
\mvfrow{误解}{"影子银行就是非法金融、地下钱庄那一套。"}
\mvfrow{事实}{它绝大部分是持牌机构在做合规业务。货币市场基金、资产证券化、信托、私募信贷基金，每一样都有监管，都有信息披露，只是不归"银行"那套监管管。名字里的"影子"说的是它站在银行监管这盏探照灯照不到的角度，不是说它在地下室。真正的非法集资是另一回事，第 23 章会单独拆。}
\end{mythbox}

\section{两条腿，砍掉一条它还站着}

影子银行不是被谁发明出来的，是被两个力量长出来的。

\textbf{第一条腿是监管套利。}巴塞尔框架下，一笔对中型企业的普通贷款风险权重按 100\% 算，按 8\% 到 10\% 的资本要求，放 1 亿要拿出八百万到一千万自有资本压在下面。这笔资本压住了就不能再用。挪到表外，通过一只基金、一个资管计划、一个特殊目的载体去做，资本就腾出来了。同一笔生意，同一个借款人，换个位置放，占用的资本能差一个数量级。

\textbf{第二条腿是真实需求。}银行的风控模型是为标准化客户造的：有抵押、有历史、现金流平稳。夏总这种公司不符合模型，可他的订单是真的，设备是真的，招进来的 40 个人也是真的。这笔钱总得有人给。

只要资本监管还在，只要银行的模型还会筛掉一部分真实的融资需求，影子银行就不可能被消灭。这一句是判断，我把理由摆在这儿了：你可以关掉 P2P，可以让理财净值化，可以把信托的通道砍断，但被赶走的资金和被拒绝的需求会在别的地方接上头，换个名字继续。过去二十年，它至少换过五个名字。

\begin{pullquote}{}
监管能改变影子银行的形状，改变不了它的存在。你堵掉一条通道，水会自己找下一条。
\end{pullquote}

\begin{mythbox}{常见误解}
\mvfrow{误解}{"监管收紧就能把影子银行消灭掉。"}
\mvfrow{事实}{2018 年资管新规之后，中国表外理财的规模确实压下来了，非标转标，通道大幅收缩。但企业的融资需求没有消失，它转到了标准化债权、私募基金、供应链金融和产业链的应收账款里。统计口径变了，链条的形状变了，业务本身还在。监管的现实目标从来不是消灭，是让它可见、可计量、有一层资本垫着。}
\end{mythbox}

\section{过去二十年，它换过几个名字}

\begin{flowlist}
  \flowitem{01}{货币市场基金}{对客户的承诺是"和现金一样，1 美元进 1 美元出"，实际干的是买商业票据。2008 年 9 月 16 日，Reserve Primary Fund 因持有雷曼商票，净值跌到 0.97 美元。承诺碎了一角，之后几天全行业被赎回上千亿美元。它从不自称银行，但它是一家没有存款保险的银行。}
  \flowitem{02}{资产证券化 ABS / MBS}{把一堆房贷打包，按偿付顺序切成不同优先级卖给不同风险偏好的人。逻辑本身没错，今天还在用。错在两处：假设各地房价不会同时跌，以及切完之后拿切剩的部分再切一遍。}
  \flowitem{03}{回购市场}{2008 年真正的挤兑不发生在银行网点。它发生在回购市场：对手方突然要求更高的折扣率，或者干脆不再续做隔夜融资。贝尔斯登和雷曼死于隔夜的钱接不上，不是死于储户排队。这一点在中文报道里几乎从没被讲清楚过。}
  \flowitem{04}{银行理财、信托、委外}{2013 到 2017 年的中国版本。理财资金经信托、券商资管、基金子公司层层通道，投向地方融资平台和房地产。刚兑预期让它在老百姓眼里像存款，底层却是非标债权，期限还错配着。}
  \flowitem{05}{P2P}{高峰期平台数超过五千家。本质上就是没有牌照、没有资本金、没有拨备的银行，负债端全是散户，期限短到一个月。它爆得最惨，因为连"机构投资者"这层缓冲都没有。}
\end{flowlist}

五个名字，五种形态，一套机制。每一次的爆点都落在同一个位置：负债端跑得比资产端快。

\bookbreak

\section{2026 年的主角叫私募信贷}

\begin{statrow}{3}
  \statitem{2.3}{万亿美元}{全球私募信贷管理规模（2025 年）}
  \statitem{17}{\%+}{2020 年以来的年化增速，五年翻倍}
  \statitem{30}{\%}{美国银行流向该领域的敞口占银行资本比重}
\end{statrow}

\term{私募信贷}{private credit}\index{056@私募信贷}就是夏总拿到那 4200 万的地方。不公开发行，不在交易所交易，一家基金直接把钱借给一家公司，条款一对一谈。

钱从哪儿来？养老金、保险公司、大学捐赠基金、家族办公室。这几类机构的共同点是负债期限极长。一家养老金要三十年后才开始大规模付钱，它完全承受得起资产五年不能变现，它想要的是可预期的票息，不是每天跳动的报价。

钱到哪儿去？中型、成长快、抵押品少、现金流还不够漂亮的公司，也就是传统银行覆盖不足的那一段。行业敞口里软件业占三成左右，其中投向 AI 相关领域的约 2000 亿美元，占全部资产 9\%。

价格 8\% 到 12\%，比银行贷款高 200 到 400 个基点；二顺位和夹层 12\% 到 16\%。这个价差不是暴利，它对应的是更差的抵押、更靠后的清偿顺序、更长的锁定期。

\section{银行没有退出，银行绕了个弯}

很多人以为私募信贷的兴起意味着银行在这块业务上撤了。银行确实不直接借给夏总了，但银行借给了借给夏总的那个人。

\begin{chainflow}
\chainnode{你的存款}\chainarrow \chainnode{银行}\chainarrow \chainnode{给私募信贷基金的授信}\chainarrow \chainnode{私募信贷基金}\chainarrow \chainnode{夏总的产线}
\end{chainflow}

这个绕弯叫\term{再中介化}{re-intermediation}\index{080@再中介化}。银行给基金提供循环授信、认购优先级份额，间接参与了企业信贷，账面上记的却是"对金融机构的贷款"，风险权重更低，资本占用更少。

规模有多大：美国银行体系流向私募信贷的资金约 5000 到 6000 亿美元，占银行总资产 3\%，占银行资本 30\%。

\begin{egbox}{把这两个比例做一次除法}
3\% 的资产，对应 30\% 的资本。这两个数放在一起，等于说这批银行的资本对资产比大约是 10\%。它也等于说：银行自有资本的三成，压在这一条通道上。

我见过不少人把"占资本 30\%"念成"银行三成资本要没了"。\textbf{这是错的。}敞口不等于损失，要变成损失得先算回收率。这条通道整体损失 10\%，损失额 550 亿美元，占银行资本 3\%；损失 30\%，占资本 9\%。难受，不是灭顶。

该担心的不是这个数字，是它在四年里从几乎为零涨到了这里，以及没有任何外部机构能确认这些底层贷款是不是用同一套假设估的值。

\end{egbox}

\section{保险公司这一环，问题不在关联交易}

另一件报道不多但很要紧的事：另类资管机构这几年大量收购保险公司，尤其是年金业务。动机很直白，保险公司手里是长期、稳定、成本可控的负债，正好是私募信贷最想要的资金。买下来之后，这些钱可以配置到集团内部发起的资产上。

部分私募股权支持的保险公司，关联方资产配置比例达到 20\% 到 40\%。行业平均水平是 7\%。

\begin{barfig}
  \barrow{私募股权支持的保险公司（区间上沿）}{80}{40\%}
  \barrow{私募股权支持的保险公司（区间下沿）}{40}{20\%}
  \barrow{行业平均}{14}{7\%}
\end{barfig}

问题不在关联交易违法。它合规，有披露，有独立董事签字。问题在估值：你自己发起的贷款，卖给你自己控股的保险公司，价格由你自己的模型给。这条链上没有一个愿意和你对着下注的第三方。

市场价格的全部意义，就是有人愿意在另一边掏钱。没有那个人，剩下的只是一个数字。

\section{把它卖给个人的那一刻}

机构资金和这类资产本来是匹配的：长钱配长资产，锁三年就锁三年。

然后有人想到，这么好的收益率，为什么不卖给个人。

于是有了半流动基金、上市的 BDC（\term{商业发展公司}{business development company}\index{046@商业发展公司}），以及包着这些资产的 ETF。到 2025 年，零售端产品规模约 4000 到 5000 亿美元，占私募信贷总规模两成。

底层贷款要三五年才能收回本金。产品对个人承诺的是季度赎回。

这就是全部的风险所在。

\begin{warnbox}{小心：慢资产配了快负债}
半流动基金一般设有赎回上限（Gate），比如每季度最多赎回净资产的 5\%。平时它是一条摆设，因为正常季度申请赎回的人只有一两个百分点。

它在市场平静时永远不会被触发，在市场紧张时一定会被触发。而它被触发的那一天，就是所有人知道"排队开始了"的那一天。

\end{warnbox}

\begin{egbox}{Gate 是怎么把恐慌放大的}
一只规模 86 亿的半流动基金，季度赎回上限是净资产的 5\%，也就是 4.3 亿。

某个季度它收到 10.4 亿的赎回申请，占净资产 12.1\%。基金按比例配售：每个人只能拿回申请金额的 41\%（4.3 ÷ 10.4）。剩下的 59\% 排队等下一个季度。

下一个季度会发生什么？看到"某基金触发赎回上限"这条新闻的人，比上个季度多得多。原本不打算赎的人也会去排队，因为排在后面的人拿不到钱。

这不是有人在恶意做空。这是每个人做出对自己最合理的选择之后，加总起来的结果。挤兑从来都是这么来的。

\end{egbox}

\begin{talkbox}
  \talkline{新闻}{"某上市 BDC 股价较其最近一期资产净值折价 18\%，创上市以来最大折价。"}
  \talkline{翻译}{这只基金说自己手上的贷款值 100 块，愿意在公开市场上买它的人只肯出 82 块。}
  \talkline{那 18 块是什么}{是市场对两件事同时报价：一，它的内部估值模型可能偏乐观；二，你想立刻拿到现金，就得付这个代价。这个折价是眼下最诚实的温度计，因为它是真金白银交易出来的，不是模型算出来的。}
\end{talkbox}

\section{它是不是下一个次贷}

这个问题被问了两年。两边的论据都不是空穴来风，得两边都摆。

\begin{tablewrap}
\begin{xltabular}{\linewidth}{rXX}
\toprule
\bthead{维度} & \bthead{2007 年的次贷} & \bthead{2026 年的私募信贷} \\
\midrule
\textbf{信用扩张环境} & 长期低利率、全球追逐收益 & 低利率十年之后的高利率，同样在追逐收益 \\
\textbf{风险下沉} & 从优质房贷下沉到次级、无收入证明 & 从大型企业下沉到中型、二顺位、夹层 \\
\textbf{底层透明度} & 打包后无法穿透到单笔贷款 & 不公开交易，按季度按模型估值 \\
\textbf{与银行的耦合} & 银行自持大量超优先级份额 & 银行向基金放贷，敞口占银行资本 30\% \\
\textbf{利率上行的冲击} & 浮动利率重置，还款额跳升 & 多为浮动利率，借款人利息负担同步上升 \\
\textbf{嵌套层数} & CDO、CDO 平方，多层再证券化 & 基本没有再证券化 \\
\textbf{杠杆倍数} & CDO 10–20 倍，投行 20–30 倍 & 基金层面 0–2 倍 \\
\textbf{负债端稳定性} & 隔夜回购、ABCP，一夜之间能抽干 & 养老金保险的多年期承诺，锁定期长 \\
\textbf{风险释放速度} & 几周 & 按季度重估，几个季度到几年 \\
\bottomrule
\end{xltabular}
\end{tablewrap}

上半张表说明它像。下半张表说明它不像。两边都是真的。

中金把它称为"2 万亿美元的灰犀牛"，判断是风险可能在未来 12 到 24 个月逐步传导，但目前不构成黑天鹅式的系统性危机。

这句话要按它本来的分量读：\textbf{它是一个判断，不是一个事实。}它站在三个假设上：负债端不会突然变短，零售化比例不会失控，投向 AI 的那 2000 亿不会集体出问题。三个里第一个最结实，第三个最悬。

我的看法：如果这轮真出事，触发点大概率不在企业违约率本身。企业是慢慢坏的，一个季度坏一点，机构消化得了。快的是产品端的赎回。快的一头撞上慢的一头，中间那条缝就是风险。

这是对机制的判断。什么时候发生、会不会发生，我不知道。

\section{你能盯的四个数}

这行有个好处：它的体温计是公开的。

\begin{flowlist}
  \flowitem{01}{BDC 的价格对净值比}{日频，公开可查。折价持续扩大并且伴随成交量放大，说明市场在给模型估值打问号。单日折价没意义，连续两个季度走扩才是信号。}
  \flowitem{02}{高收益债利差和杠杆贷款价格}{私募信贷自己不交易，它的表兄弟天天交易。杠杆贷款平均报价跌破 95 通常意味着压力开始了。表兄弟先咳嗽。}
  \flowitem{03}{大型银行和另类资管机构的 CDS 报价}{市场为"这家机构会不会出事"直接标的价，比任何评级都灵敏，因为报价的人要用自己的钱兑现。}
  \flowitem{04}{赎回上限有没有被触发}{这类消息通常不上头条，藏在产品公告里。第一家触发的时候，市场往往还很平静。}
\end{flowlist}

说清楚：这四个数不是买卖信号，本书不提供投资建议。它们唯一的用处，是让你在新闻标题变得吓人之前，先知道水位到了哪一格。

\section{中国这边是反过来走的}

2018 年资管新规是个分水岭。打破刚兑，产品净值化，非标转标，多层嵌套的通道被砍断。信托被迫转型，从"影子银行的通道"往资产服务信托、资产管理信托、公益慈善信托这三类走，很多公司到今天还没转完。

地方融资平台的债务在化解，手段是展期、置换、降息：短的换成长的，贵的换成便宜的，隐性的挪到显性的账上。这不消灭债务，它买时间。

方向和美国正好相反。美国是银行受资本约束往后退，把资产让给非银机构；中国是先长出一个庞大的影子银行体系，再用行政力量把它往表内和标准化市场里赶。

共同点只有一个：无论赶到哪边，那笔钱最后还是要落到某个像夏总一样的人手上，去买一台设备，雇几十个人。中间过几道手，决定的是谁承担损失，不是钱去干什么。

\begin{takeaway}{本章一句话}
影子银行不是银行的敌人，是银行监管的影子。它干的是同一件事：拿短钱做长事。它省下的是资本占用，代价是出事时没有存款保险，也没有最后贷款人。

\begin{itemize}
  \item 判断标准只有一条：谁在做期限、流动性、信用这三种转换，谁就在做银行的生意。
  \item 它有两条腿：监管套利和真实需求。这两条腿在，它就永远在，只是换名字。
  \item 2026 年它最大的一块叫私募信贷，2.3 万亿美元。风险不在违约率，在慢资产配了快负债。
  \item "占银行资本 30\%"是敞口不是损失。看数据要先弄清楚它是哪一种。
\end{itemize}

\end{takeaway}

\begin{bridgebox}
\textbf{下一章}影子银行自己不会引爆什么。它得先和另一样东西凑到一起：一个所有人都相信、而且确实是真的新故事。这两样一凑，就会重演一部上演了四百年的老剧本。
\end{bridgebox}


\chapter{泡沫与崩溃的通用剧本}
\chaptermeta{17}

\begin{ledetext}
四百年来它只演过一出戏，换的只是布景和主角的行头。剧本你看完就会背，难的是知道自己坐在第几幕。
\end{ledetext}

\section{牛顿的那笔账}

1720 年，艾萨克·牛顿在南海公司的股票上亏掉了大约两万英镑。当时他是英国皇家铸币厂厂长，年收入两千英镑上下。

二十年薪水。

流传最广的那句"我能计算天体的运行，却算不出人类的疯狂"，出处一直存疑，最早的记录出现在他死后很多年。但那笔亏损是有账可查的。

更要命的是细节：他不是从头套到尾的。1720 年初他买入，四月前后卖出，赚了大概七千英镑，本金翻了一倍。然后他站在场外，看着股价继续往上冲。六月他又进去了，这次买在接近顶部的位置。

让他破产的不是贪婪。是眼看着别人赚，而自己已经下车。

\section{五幕}

海曼·明斯基和查尔斯·金德尔伯格把这出戏拆成了五幕。这个结构从 1637 年用到 2026 年，一次没换过。

\begin{flowlist}
  \flowitem{01}{位移}{一件真事发生了。蒸汽机、铁路、电力、互联网、页岩油、大模型。它真的改变了利润的分布，也真的创造出一个以前不存在的赚钱机会。\textbf{泡沫的种子从来不是谎言，是真话。}}
  \flowitem{02}{繁荣}{信贷跟上来。银行、券商、基金开始往这个方向配钱。价格涨，早进去的人赚到了，而"有人真的赚到钱"这个事实，比任何论证都有说服力。}
  \flowitem{03}{狂热}{新人涌进来。估值和现金流脱钩，开始改用新指标：点击量、日活、GMV、算力储备。出现"这次不一样"的论述，而且这套论述通常有一半是对的，这才是它难对付的地方。}
  \flowitem{04}{获利了结}{最早进场的人开始减持。他们不发声明，他们只是卖。这个阶段价格往往还在涨，因为接的人比卖的人多。事后回看，这一幕最清楚；身在其中，这一幕最像"回调"。}
  \flowitem{05}{恐慌}{一件具体的、往往很小的事情触发它。流动性消失，杠杆被强平，强平压低价格，价格触发下一批强平。}
\end{flowlist}

\begin{chainflow}
\chainnode{价格下跌}\chainarrow \chainnode{追加保证金}\chainarrow \chainnode{被迫卖出}\chainarrow \chainnode{价格再跌}\chainarrow \chainnode{下一批被迫卖出}
\end{chainflow}

五幕之间没有钟点表。第三幕可以持续三个月，也可以持续四年。"识别泡沫"和"知道它什么时候破"是两个难度差着一个数量级的问题，绝大部分把这两件事混为一谈的人，都在第三幕里亏掉了钱。

\section{四个零件，缺一个就不是泡沫}

一个真实的新故事，便宜的信贷，杠杆，大量新参与者。四样凑齐才叫泡沫，少一样就只是一次普通的上涨。

第一样最反直觉，所以多说两句。

铁路是真的。1840 年代英国铁路狂热崩掉之后，英国确实得到了世界上第一张全国铁路网，那些铁轨后来跑了一百多年。互联网是真的，2000 年之后光纤、机房和一整代工程师留在了地上，被下一批公司用得很爽。AI 也是真的。

泡沫化的从来不是那个故事。

\begin{pullquote}{}
市场很少猜错方向，它经常猜错速度：把十五年才能兑现的现金流，按三年兑现折进今天的价格里。
\end{pullquote}

\begin{egbox}{把"时间表"算成数字}
一家公司今年收入 12 亿，市场给 40 倍市销率，市值 480 亿。

这个定价隐含的意思是：它得把收入做到 80 亿以上，并且做出 20\% 的净利率（16 亿利润），才能落到 30 倍市盈率这种"不算离谱"的位置。收入要涨到 6.7 倍。

每年增长 45\%，需要 5 年零 2 个月。\\ 每年增长 25\%，需要 8 年零 6 个月。

同一家公司，同一个故事，增速差 20 个百分点，兑现时间差三年多。而市场愿意等的时间，通常比它嘴上说的短得多。第三幕的崩塌几乎从来不是因为"故事被证伪"，是因为\textbf{某个季度的增速从 45\% 掉到了 25\%}。

\end{egbox}

\section{点名，每个只说最要命的那一句}

\textbf{1637 年，荷兰郁金香。}这个故事你一定听过：一颗球茎换一栋运河边的房子。

近二十年的档案研究给出的图景要冷淡得多。卷进去的主要是一个不大的、彼此认识的商人和手艺人圈子；成交大量是期货合约而不是现货交割；价格崩溃之后，荷兰经济没有出现可观察的衰退；法院基本以"这类合约不可强制执行"了事。它的经济影响被后世严重夸大了。

它被引用了四百年，是因为它简短、荒诞、好用。真正的教训不在球茎里，在你为什么会毫不怀疑地接受一个自己从没核对过来源的故事。这一段我特意放在最前面：怀疑现成叙事，本身就是这一章要教的东西。

\textbf{1720 年，南海公司。}牛顿那笔账。补一句：南海公司的商业模式是接管英国国债，它和政府的关系比它和南海贸易的关系密切得多。

\textbf{1929 年，美股。}要害是保证金。

\begin{egbox}{10\% 首付买股票}
1928 年，你拿出 1000 美元，按 10\% 的保证金买入 10000 美元的股票，剩下 9000 向券商借。

股票涨 20\%，市值 12000，还掉 9000，你手里剩 3000。本金翻两倍。\\ 股票跌 10\%，市值 9000，你的 1000 归零，追加保证金的电话打过来。\\ 股票跌 12\%，你不但归零，还倒欠券商 120。

1929 年 10 月底那几天，指数单日跌幅超过 12\%。电话打不通的结果是券商直接平仓，平仓压低价格，价格触发下一批平仓。第 14 章那台绞肉机的历史现场，就在这儿。

\end{egbox}

\textbf{1990 年，日本。}顶点那几年流传过一个说法：东京都的土地按市价算，卖掉可以买下整个美国。这个算法本身很粗糙，但它当时被广泛相信，这件事比算得准不准更重要。日经 225 在 1989 年 12 月 29 日收于 38957 点，此后花了 34 年才回到这个位置。

\textbf{2000 年，纳斯达克。}5048 点跌到 1114 点，跌掉 78\%，回到 5048 用了 15 年。

同一时期，亚马逊从最高点跌了约 94\%，从一百多美元跌到五块多。后来它涨了几百倍。

同一个泡沫，同一个方向判断（"互联网会改变世界"这句话完全正确），标的选择决定了两种完全相反的命运。而这两种命运对普通人的意义又完全不同：拿着一篮子纳斯达克，考验的是耐心；拿着亚马逊，考验的是你在跌掉 94\% 的那段路上会不会被迫卖出。

被迫卖出。这四个字才是重点。你有没有杠杆，会不会失业，家里有没有别的钱要用，决定了你能不能活到兑现那一天。\hlmark{在泡沫里，"看对趋势"的价值远低于"活到兑现"。}而后者取决于杠杆和现金流，不取决于眼光。

\textbf{2008 年，次贷。}一句话讲完：把一个统计观察（美国房价从未在全国范围内同时下跌）当成了物理定律，用它给证券定价，请评级机构盖章，然后在上面叠了二三十倍杠杆。假设错一次，杠杆把这个错误放大二三十倍还给你。

\textbf{2023 年，硅谷银行。}它没有坏账。它买的是美国国债和机构 MBS，信用风险接近于零。它错在久期：用随时可以提走的存款，买了平均久期五六年的债。利率从零升到 5\% 以上，债价跌，浮亏挂在"持有至到期"科目里不用计提，直到它必须卖出变现的那天。

真正的新东西是速度。2023 年 3 月 9 日一天被提走 420 亿美元，第二天预定还要走 1000 亿。1930 年代的挤兑要排队，2023 年的挤兑是几个微信群和 Slack 频道里的截图，加上手机银行的转出按钮。这是四百年里唯一一条被技术真正改写的规则。

\begin{tablewrap}
\begin{xltabular}{\linewidth}{>{\raggedright\arraybackslash}p{0.20\linewidth}XrX}
\toprule
\bthead{泡沫} & \bthead{顶到底} & \bthead{最大跌幅} & \bthead{回到原点用了} \\
\midrule
\textbf{1929 道琼斯} & 1929/09 – 1932/07 & -89\% & 25 年 \\
\textbf{1989 日经 225} & 1989/12 – 2009/03 & -82\% & 34 年（2024 年 2 月） \\
\textbf{2000 纳斯达克} & 2000/03 – 2002/10 & -78\% & 15 年 \\
\textbf{2007 标普 500} & 2007/10 – 2009/03 & -57\% & 5 年半 \\
\bottomrule
\end{xltabular}
\end{tablewrap}

看最后一列。这一列决定的不是收益率，是你的人生安排：25 年和 5 年半，对一个 55 岁准备退休的人来说是两个物种。

\bookbreak

\section{破了之后，只有三种结局}

\textbf{温和出清。}价格跌下来，杠杆低的机构接手便宜资产，两三年翻篇。1987 年的股灾是这一类，一天跌 22\%，第二年就修复了。前提是杠杆没有渗进银行体系。

\textbf{资产负债表衰退。}日本那一种。企业的资产跌破了负债，但现金流还在，于是它们把经营目标从"利润最大化"改成"负债最小化"：赚来的钱全部拿去还债，不再借钱扩张，哪怕利率是零。

这时候货币政策失效，不是因为央行不努力，是因为问题不在钱贵，在没人想借。辜朝明在 1990 年代提出这个描述时不被主流接受，现在它被广泛引用。日本用了差不多三十年才走出来。

\textbf{系统性危机加政府兜底。}2008 年那一种。政府接管、注资、担保，把私人部门的债务转成公共部门的债务。危机被中止了，代价是政府债务率永久性抬高一个台阶，而且回不去。美国联邦债务占 GDP 的比重在 2007 年是 60\% 出头，此后再没回到那个水平；2026 年 8 月，美国国债规模突破 40 万亿美元。

三种结局的分野只有一条：\textbf{杠杆有没有渗进银行体系。}没渗进去，是财富重新分配；渗进去了，是支付系统出问题。这条线在第 14 章画过一次，这里再画一遍。

\section{稳定本身孕育不稳定}

明斯基最反直觉的一句话，也是他最重要的一句。

\begin{deepbox}{进阶：可以跳过}
明斯基把借款人分成三类。\textbf{对冲型}：现金流能覆盖利息和本金。\textbf{投机型}：现金流只够付利息，本金到期要靠续借。\textbf{庞氏型}：现金流连利息都不够，全靠资产升值和再融资撑着。

他的论点是：一个经济体在长期平稳中，会自动从第一类往第三类漂移。不是因为人变坏了，是因为前面几年里谁激进谁赚钱，激进的机构份额变大，保守的基金经理被赎回或者被换掉。市场在用最有效率的方式，筛选出最脆弱的结构。

\end{deepbox}

日子越安稳，人越敢借；越敢借，系统越脆弱；越脆弱的系统在没出事之前，看起来越稳。

所以"这次很稳"这句话本身就是一个风险指标。听到它的频率越高，位置越靠后。

\section{AI 是不是泡沫}

这是 2026 年被问得最多的问题。我不会给答案，但我会把两边的证据都摆出来，最后给一个你自己能用的框架。

\begin{statrow}{3}
  \statitem{57}{\%}{2025 年数据中心资本开支同比增速}
  \statitem{2000}{亿美元}{私募信贷投向 AI 领域的规模}
  \statitem{78}{\%}{纳斯达克 2000–2002 最大跌幅（参照）}
\end{statrow}

\textbf{说它不是泡沫的证据。}2025 年数据中心资本开支同比增长约 57\%，出钱的主力是现金流极强的巨头，不是靠融资活着的初创公司。收入是真的，需求是真的，算力目前仍然供不应求。美联储主席凯文·沃什在 2026 年 7 月公开表示，人工智能相关投资正在为未来经济增长奠定基础。

\textbf{值得警惕的信号。}不止一条：

\begin{warnbox}{小心这四件事}
\textbf{折旧年限的会计假设。}把 GPU 的折旧年限从三年拉长到五年六年，当期利润立刻好看。如果这批芯片的实际有效寿命只有三年，多出来的那部分利润是从未来借的。

\textbf{供应链上的循环交易。}A 投资 B，B 拿这笔钱买 A 的产品，A 的收入增长了。分开看每一笔都合规，合并看是左手倒右手。这种结构在 2000 年的电信设备行业出现过。

\textbf{债务融资的比例在上升。}早期是现金流出钱，现在开始发债、做结构化融资、找私募信贷。上一章那 2000 亿美元就在这里。零件二和零件三正在到位。

\textbf{市场集中度极高。}指数涨幅高度依赖少数几家公司，这意味着"分散投资"提供的保护比你以为的薄。

\end{warnbox}

\begin{talkbox}
  \talkline{沃什，2026 年 7 月}{"人工智能相关投资正在为未来经济增长奠定基础。"}
  \talkline{同一场发布会}{"资本支出供应副作用的发生时间难以预测。"}
  \talkline{翻译}{方向他认，节奏他不认。这两句话不矛盾，它们合起来正好是本章的主题：故事可能是真的，时间表可能是错的。}
\end{talkbox}

给你三个可以自己查、自己算的指标，不需要任何内幕消息：

\textbf{一，这一轮扩张的钱从哪来？}经营现金流的占比在上升还是下降？债务融资的占比在往哪个方向走？这个比例的拐点比总量重要得多。

\textbf{二，收入跟得上资本开支吗？}把两条曲线的同比增速画在一张图上。差距是在收窄还是在扩大？扩大了几个季度？

\textbf{三，新进场的人有没有加杠杆？}散户是用闲钱买，还是用融资余额买？相关主题的杠杆 ETF 规模在涨吗？零件四到位没有，看这个。

这三个问题我现在都没有确定答案。\textbf{判断框架给你，答案时间会给。}

\begin{mythbox}{常见误解}
\mvfrow{误解}{"读完这一章我就能提前认出泡沫，然后在破之前跑掉。"}
\mvfrow{事实}{事后看曲线，一切清清楚楚；事前你手里只有一个概率分布。格林斯潘 1996 年 12 月说出"非理性繁荣"的时候，纳斯达克在 1300 点附近，此后它又涨了将近三倍，2000 年 3 月才见顶。他说对了性质，说早了三年多，而早三年的代价可能比说错更大。你真正能做的只有两件：控制杠杆；保证自己在下跌 50\% 的时候不会被迫卖出。这不是投资建议，这是关于"什么做得到、什么做不到"的说明。}
\end{mythbox}

\begin{takeaway}{本章一句话}
泡沫需要四个零件：一个真实的新故事、便宜的信贷、杠杆、大量新参与者。故事通常是真的，错的是价格和时间表。

\begin{itemize}
  \item 五幕：位移、繁荣、狂热、获利了结、恐慌。幕与幕之间没有时刻表。
  \item 破裂后只有三种结局，分野在于杠杆有没有渗进银行体系。
  \item 稳定孕育不稳定。"这次很稳"是风险指标，不是安慰。
  \item 2000 年纳斯达克回本用了 15 年，同期亚马逊跌 94\% 后涨了几百倍。看对方向不等于活到兑现。
\end{itemize}

\end{takeaway}

\begin{bridgebox}
\textbf{下一章}泡沫破了会痛，但它还会再来。因为它不是孤立事件，它挂在一个更大的钟摆上，而那个钟摆有好几个，长度还不一样。
\end{bridgebox}


\chapter{周期：为什么好日子和坏日子会轮流来}
\chaptermeta{18}

\begin{ledetext}
你算不出下一个拐点，这件事没有例外。但你能判断自己现在站在哪一段，而这才是有用的那个问题。
\end{ledetext}

\section{先把钟表扔了}

"下一个拐点什么时候到"是这本书里我可以最干脆回答的问题：不知道，而且没有人知道。

"周期"这个词本身就在骗人。它让人想起钟摆和正弦曲线，想起"上一轮走了 8 年，所以这一轮也是 8 年"。真实的周期长度从三年到三十年不等，方差大得离谱，而且只能事后量出来。

能做的只有一件事：判断自己大概站在什么位置。不是问"什么时候转向"，是问"现在算早、算中、还是算晚"。

这两个问题的难度差着一个数量级。

像在山里判断天气。你说不出下午三点几分下雨，但你能看云、看风向、看气压计，判断出"今天大概率要下"，于是带伞、不上山脊、下午四点前撤。做不到预测，做得到不被淋死。

这个比喻还能再推一步：山里真正的失误从来不是"没带伞"，是明明看见云压下来了，因为已经爬了三个小时不甘心，硬要往上走。周期里也一样。绝大部分亏损不来自看错，来自不肯认。

\section{猪}

河南南阳有个养猪场，场主姓宋，2019 年的时候存栏三千头。

那年非洲猪瘟叠加产能出清，猪价一路冲到 40 块一公斤。2020 年出栏一头猪，毛利两千多。宋老板做了一个所有人都会做的决定：扩栏，从 3000 头加到 8000 头。

问题在于从决定到出肉要多久。母猪配种到产仔 114 天，仔猪育肥到出栏 6 个月上下。如果是从买后备母猪开始算，完整一轮接近 14 到 18 个月。

他 2020 年做的决定，2021 年下半年才变成猪肉。

而全国所有看到同一个价格的养殖户，在同一段时间做了同一个决定。

2021 年 10 月，生猪价格跌到 10 块出头一公斤，全行业亏损。宋老板开始淘汰母猪。全国都在淘汰母猪。淘汰的效果同样要等一年多才显现，于是 2022 年下半年供给缺口出来了，价格又涨回 28 块。

然后新一轮补栏开始。

\begin{chainflow}
\chainnode{猪价高}\chainarrow \chainnode{补栏}\chainarrow \chainnode{14 个月后集中出栏}\chainarrow \chainnode{价格暴跌}\chainarrow \chainnode{淘汰母猪}\chainarrow \chainnode{一年后供给不足}\chainarrow \chainnode{价格暴涨}
\end{chainflow}

\begin{egbox}{宋老板的两笔账}
2020 年：养殖完全成本约 15.6 元/公斤，售价 33 元/公斤。一头 120 公斤的猪，赚 \textbf{2088 元}。

2021 年 10 月：饲料涨价，成本变成 17.2 元/公斤，售价 10.8 元/公斤。同样一头 120 公斤的猪，亏 \textbf{768 元}。

同一个猪场，同一个人，同样的管理水平，一头猪的盈亏差了 2856 元。

中间他什么都没做错。他只是在 2020 年做了一个当时看完全正确、而且所有同行都在做的决定。

\end{egbox}

这套机制有个名字，叫\term{蛛网模型}{cobweb model}\index{088@蛛网模型}：生产决策依据的是上一期的价格，而产品真正出来的时候，价格已经被所有人的同一个决策改掉了。把价格和产量画在一张图上，轨迹会绕成一圈一圈的蛛网。

关键在于它收敛还是发散。如果供给对价格的反应比需求更敏感，蛛网就越绕越大，波动自我放大。生猪是典型的发散型。

往下推一步就明白了：\hlmark{时滞越长的品种，波动越猛。}鸡蛋的周期比猪短，大蒜和生姜比猪更短更疯。多晶硅料的产能建设周期是 18 到 24 个月，它的价格从 2022 年高点到 2024 年低点跌掉八成以上，用的是同一个模型。锂也是。

把这个模型放大到整个经济，就是库存周期。

\section{四个钟摆挂在一起}

\begin{tablewrap}
\begin{xltabular}{\linewidth}{>{\raggedright\arraybackslash}p{0.20\linewidth}rXX}
\toprule
\bthead{周期} & \bthead{长度} & \bthead{驱动力} & \bthead{盯什么} \\
\midrule
\textbf{库存（基钦）} & 3–4 年 & 订单波动沿产业链被层层放大 & PMI 新订单减产成品库存、工业企业产成品存货同比 \\
\textbf{设备资本开支（朱格拉）} & 7–11 年 & 设备寿命、技术换代、产能利用率 & 固定资产投资中的设备工器具购置、产能利用率 \\
\textbf{房地产建筑（库兹涅茨）} & 15–25 年 & 人口迁移、城镇化、家庭形成 & 新开工面积、土地成交、二手房挂牌量 \\
\textbf{债务长周期} & 50–75 年 & 债务对收入比的累积与清算 & 宏观杠杆率、利息支出占 GDP 比重 \\
\textbf{政治周期} & 看换届 & 选举与换届改变政策的松紧节奏 & 财政赤字率、专项债发行节奏、规划年份 \\
\bottomrule
\end{xltabular}
\end{tablewrap}

它们同时在走，互相干扰。所以你看到的从来不是一条干净的正弦曲线，是四五条不同频率的波叠加之后的那团毛线。

2008 年之所以那么惨，是因为库存周期、设备周期、地产周期和债务周期在同一个时间点全部转向下行。这种事情几十年才凑一次。

\section{周期为什么非存在不可}

\textbf{产能建设有时滞。}一座 12 英寸晶圆厂从动土到量产要 24 到 30 个月，一个铜矿从勘探到出矿七到十年。所有人在同一个价格信号下同时决定扩产，两年后同时投产。宋老板的猪场只是这件事最好懂的版本。

\textbf{信贷会加速一切。}抵押品升值，你能借到更多；借到更多，你买更多资产；资产更贵，抵押品又升值。这个正反馈叫\term{金融加速器}{financial accelerator}\index{028@金融加速器}。它在上行时把繁荣推得更高，在下行时同样对称地起作用，而且下行更快，因为催收比放贷容易。

\textbf{人会从众，而且是被制度逼的。}一个基金经理跑输两个季度要写报告解释，跑输四个季度可能要换人。"和大家一起错"的职业成本，远低于"一个人对"。凯恩斯那句话讲得最准：宁可循规蹈矩地失败，也不要标新立异地成功。

第三条最容易被当成道德问题。它不是。你把任何一个理性的人放进那个考核机制里，他都会做出同样的选择。

\bookbreak

\section{那根最长的钟摆}

短周期决定你今年好不好过，债务长周期决定你这一代人赶上的是什么剧本。

机制其实很朴素：加杠杆，资产涨，抵押品升值，于是借得更多。只要债务增速持续高于收入增速，这个循环就能自我强化很多年，看起来像永动机。

直到某个时点，利息负担吃掉了现金流。

\begin{egbox}{利息是怎么悄悄吃掉一家公司的}
一家公司年收入 3.2 亿，息税前利润 4100 万，负债 2.6 亿，平均融资成本 4.2\%。年利息 1092 万，利息保障倍数 3.8 倍。健康。

它每年借新钱扩张。三年后：负债 4.9 亿，收入 4.4 亿，息税前利润 5300 万。同期利率从 4.2\% 升到 6.5\%。

年利息变成 3185 万。利息保障倍数掉到 \textbf{1.66 倍}。

看它的报表：收入涨了 37\%，利润涨了 29\%，还在增长。但它离"利息吃掉全部经营利润"只剩不到一倍的空间。再来一次加息，或者一个季度的订单下滑，它就得开始卖资产。

这不是虚构的极端案例。过去三年很多中型企业走的就是这条路。

\end{egbox}

然后进入去杠杆。可用的手段只有四种：

\begin{tablewrap}
\begin{xltabular}{\linewidth}{>{\raggedright\arraybackslash}p{0.20\linewidth}XX}
\toprule
\bthead{手段} & \bthead{谁承担损失} & \bthead{后果} \\
\midrule
\textbf{紧缩（少花多还）} & 债务人 & 所有人一起省钱，收入同步下滑，债务对收入比可能不降反升 \\
\textbf{债务重组与违约} & 债权人 & 债权人的资产负债表受损，可能引发连锁 \\
\textbf{财富转移（税收与转移支付）} & 高财富群体 & 政治阻力最大，推进最慢 \\
\textbf{印钞与债务货币化} & 所有持币者 & 名义收入上升摊薄债务，代价是货币贬值和通胀 \\
\bottomrule
\end{xltabular}
\end{tablewrap}

四种手段的配比不同，结果完全不同。

全靠第一种，是 1930 年代的通缩螺旋：每个人都在还债，每个人的收入都是别人的支出，于是收入下降得比债务还快，杠杆率越去越高。全靠第四种，是货币贬值和恶性通胀。四者平衡、让名义增长率略高于名义利率、债务比重慢慢化掉而通胀不失控，达利欧管这叫\term{美丽的去杠杆}{beautiful deleveraging}\index{037@美丽的去杠杆}。

\begin{pullquote}{}
这四个旋钮怎么拧，从来不是经济学问题，是政治问题。它决定的是谁来承担损失，而这件事不由计算器决定。
\end{pullquote}

\section{中国站在哪一段}

不给预测，只摆位置。

房地产处在库兹涅茨长周期的调整段。这是十五到二十五年尺度的东西，不是一个季度的政策能翻转的，2026 年 8 月底出台的政策"组合拳"改变的是斜率和节奏，不是周期本身。

地方债务在化解，用的是展期、置换、降息，把短的换成长的、贵的换成便宜的。对照上面那张表，这是第一种和第二种手段的温和混合，避开了第四种。

人口结构过了拐点。这一条会同时压住地产周期和债务周期的分母。

产业升级和部分行业的产能过剩同时发生。光伏、锂电、一部分新能源车产业链，是朱格拉周期在同一时点被政策和资本一起推起来之后的自然结果。被叫做"内卷"的那个东西，用蛛网模型解释比用道德解释准确得多。

\begin{statrow}{3}
  \statitem{4.7}{\%}{中国 2026 年上半年 GDP 同比}
  \statitem{60}{万亿元}{2030 年社会消费品零售总额目标}
  \statitem{3.5}{\%}{5 年期以上 LPR（2026 年 8 月 20 日）}
\end{statrow}

商务部等 7 部门提出，到 2030 年社会消费品零售总额达到 60 万亿元左右。

这句话的信息量比数字本身大。\textbf{政策重心从投资转向消费，这本身就是周期位置的一个信号。}一个经济体处在投资周期的上行段，政策语言是"扩大有效投资"；处在投资周期的下行段和债务周期的去杠杆段，语言会变成"提振消费""扩大内需"。

官方文件的措辞变化，通常比统计数据更早。

\section{全球的账，和一个不能不提的尾部风险}

\begin{databox}{2026 数据}
国际金融协会（IIF）数据：截至 2026 年 3 月底，全球债务规模约 \textbf{353 万亿美元}，创历史新高，占全球 GDP 约 \textbf{305\%}。

IIF 点名的结构性推升力量有五个：人口老龄化、国防开支增加、能源安全与来源多元化、网络安全、AI 相关资本支出。

\end{databox}

这五样有一个共同点：都是不得不花的钱。

这一点让本轮债务扩张和过去几轮很不一样。过去的债务扩张主要由消费和地产驱动，那种债务在利率上行时是可以被压下去的：少买一套房，晚换一辆车，需求就萎缩了。

这一轮压不下去。养老金的支付义务不会因为国债收益率上行而减少；国防预算不会因为融资成本高就砍掉；算力如果不投，输掉的是一个国家未来二十年在产业链上的位置。

所以本轮债务对利率的敏感度比过去低。这句话听起来像好消息，其实是坏消息：\hlwarm{不敏感意味着即使利率很高，债务也停不下来。}

\begin{warnbox}{小心：2026 年的这个组合}
供给冲击（中东局势推高能源价格）＋ 高利率（美联储 3.50\%–3.75\%，日本央行 6 月加息至 1.0\%，欧洲央行 6 月三年来首次加息）＋ 高债务（353 万亿美元，占 GDP 305\%）。

债务可持续性说到底取决于一个不等式：名义利率 r 和名义增长率 g 谁大。r 小于 g，债务占 GDP 的比重会自己往下走；r 大于 g，就算一分钱新债都不借，比重也会自己往上爬。

滞胀是唯一一种能同时把这两边推向坏处的组合：通胀逼着央行抬高 r，停滞压低 g。

这不是预测。没有人知道下半年油价往哪走。这是在指出这个组合为什么危险。

\end{warnbox}

\begin{egbox}{r 和 g 差两个百分点是什么概念}
一个国家债务占 GDP 200\%，名义 GDP 增速 5\%，平均融资成本 3\%。就算财政收支完全平衡，债务比重每年也会自己下降约 3.8 个百分点：200\% × (3\% − 5\%) ÷ 1.05。

换个环境：增速掉到 2\%，融资成本升到 4.5\%。同样一分钱新债不借，比重每年自己上升约 4.9 个百分点。

两种情形里没有任何人多借一分钱，也没有任何政策失误。差别只来自两个数字谁大谁小。

这就是为什么各国政府对通胀的态度那么纠结：温和的通胀能抬高 g，帮着化债；失控的通胀会先把 r 顶上去，把好处全吐回来。

\end{egbox}

\begin{mythbox}{常见误解}
\mvfrow{误解}{"周期有固定长度，把历史数据一平均就能算出下一个拐点。"}
\mvfrow{事实}{那些长度是事后统计出来的平均值，方差大到没法用来做决策。美国战后最短的一轮扩张只有 12 个月，最长的一轮 128 个月（2009 到 2020）。拿平均值去推算下一次，等于拿"人均寿命 79 岁"去推算某一个具体的人哪天走。周期是存在的，节拍器是不存在的。}
\end{mythbox}

\begin{mythbox}{常见误解}
\mvfrow{误解}{"政府有那么多工具，早晚能把周期熨平。"}
\mvfrow{事实}{政策能改变周期的形状，削平最高点，垫高最低点，但改变不了它存在。更要紧的是，政策本身就是下一个周期的原因之一。2008 年之后的超低利率托住了资产价格，也养出了后来的高杠杆；2020 年的财政刺激避免了衰退，也贡献了 2021 到 2023 年的通胀。救助的代价通常有两种形式：一种是把问题挪到未来，一种是把问题挪给别人（储蓄者、后来买房的人、拿固定收入的人）。这两种代价都不会出现在救助当天的新闻标题里。}
\end{mythbox}

说清楚一件事：这一章没有任何内容可以用来择时，本书也不提供投资建议。判断位置和预测拐点是两回事，前者影响你该背多重的包，后者不影响任何事，因为它做不到。

\begin{takeaway}{本章一句话}
周期不是钟表，是四五个不同长度的钟摆挂在一起。你算不出下一个拐点，但你能判断自己站在哪一段：债务是在长还是在还，产能是在建还是在关，新进场的人是多还是少。

\begin{itemize}
  \item 库存 3–4 年，设备 7–11 年，地产 15–25 年，债务 50–75 年。同时在走，互相干扰。
  \item 时滞、信贷加速器、从众。三个机制里只有第三个看着像人性问题，其实是考核机制问题。
  \item 债务增速长期高于收入增速，终点是确定的，时间是不确定的。
  \item 去杠杆的四个旋钮怎么拧，决定谁来承担损失。这是政治问题。
\end{itemize}

\end{takeaway}

\begin{bridgebox}
\textbf{下一章}天气看完了。接下来只剩一个问题：这些东西跟你银行卡里那点钱，到底有什么关系？
\end{bridgebox}



\part{落地篇：你在其中怎么活}
\chapter{给自己做一张资产负债表}
\chaptermeta{19}

\begin{ledetext}
第 5 章那三个词，汇丰在用，你们楼下便利店也在用。这一章读完要动手，不动手等于没读。
\end{ledetext}

\section{林悦不知道自己有多少钱}

林悦 32 岁，在杭州一家做企业 SaaS 的公司当产品经理。税前 2 万 6，五险一金扣完到手 1 万 9 出头。2021 年秋天她买了房，总价 386 万，首付 116 万，贷了 270 万，三十年。

过年饭桌上有亲戚问她现在有多少身家。她说：一套房，还欠着银行。

这个回答等于没说。

第 5 章那个式子是\textbf{资产 = 负债 + 净值}，全世界的会计都在用它，它不挑规模。套在林悦身上，需要的只是两列数字和一次减法。

唯一的技术难点在怎么填房子那一栏。买入价 386 万是历史，跟今天没有关系。\hlmark{房子按现在能卖掉的价钱记，房贷按银行系统里的余额记。}她翻了小区最近半年的成交记录，同户型同朝向，328 万落槌。就用 328 万。

\begin{bsheet}
  \bsside{资产 ASSETS}{%
    \bsrow{活期存款}{8,000}
    \bsrow{货币基金}{4.2 万}
    \bsrow{股票型基金}{11 万}
    \bsrow{公积金账户}{9.6 万}
    \bsrow{自住房（2026 年 8 月市价）}{328 万}
    \bsrowtotal{合计}{353.6 万}
  }
  \bsside{负债 + 净值}{%
    \bsrow{住房贷款余额}{247 万}
    \bsrow{信用卡未还}{1.3 万}
    \bsrow{车贷}{0}
    \bsrow{净值}{105.3 万}
    \bsrowtotal{合计}{353.6 万}
  }
\end{bsheet}

105.3 万。这是林悦真正拥有的东西，是把所有杠杆抽掉之后剩下的那一块。

接着是难受的一段算术。

\begin{egbox}{房子这一栏，五年下来是赚是赔}
林悦在这套房上投进去的本金是：首付 116 万，加上五年里还掉的本金 23 万（贷款从 270 万降到 247 万），合计 \textbf{139 万}。同期还掉的七十多万利息不算在内，那是买了五年的居住权。

今天这套房属于她的部分是：328 万市价减去 247 万余额，\textbf{81 万}。

139 万进去，81 万出来。这一栏五年少了 58 万。

她的净值总共 105.3 万，其中房子贡献 81 万。\textbf{她一直觉得房子是自己最成功的一笔投资，实际上是亏得最多的那一笔。}

\end{egbox}

\begin{warnbox}{小心：这个差额可以是负数}
林悦算是运气好的。换一个 2021 年在环沪、环郑、环沈那些地方接盘的人，跌四成不稀奇：386 万的房子如果今天只能卖 230 万，而贷款余额还有 247 万，资产负债表上这一栏就是 \textbf{−17 万}。

整张表的净值可能因此变成负的。人还是那个人，工作还是那份工作，只是表翻过去了。这个状态有个名字，叫\term{资不抵债}{negative equity}\index{092@资不抵债}。2008 年之后的美国有上千万套住房处在这个状态里，2026 年的中国有多少，没有公开统计。

不做这张表，你不会知道自己在不在里面。

\end{warnbox}

\begin{mythbox}{常见误解}
\mvfrow{误解}{"我有一套房，所以我有钱。"}
\mvfrow{事实}{房子记在资产端，房贷记在负债端，属于你的只有中间那个差。而且这个差有三个毛病：不能吃，不能应急，不能切一块下来卖。林悦手里能在 24 小时内变成钱的东西一共 5 万，她的表上却写着 353.6 万的资产。\textbf{净值是身家，流动性才是生活。}}
\end{mythbox}

\section{那张表是照片，这张才是心电图}

资产负债表拍的是某一天的静态。真正决定你三年后长什么样的，是每个月流过你账户的那道水。

结构简单到不需要软件：\textbf{收入 − 刚性支出 − 弹性支出 = 储蓄}。刚性支出是砍不掉的（月供、房租、保费、给父母的钱），弹性支出是失业第二个月就会自动消失的（旅游、外卖、健身卡、想买的那件外套）。

林悦一直觉得自己存不下钱。她把数字列出来之后愣了一下。

\begin{tablewrap}
\begin{xltabular}{\linewidth}{>{\raggedright\arraybackslash}p{0.20\linewidth}Xr}
\toprule
\bthead{项目} & \bthead{说明} & \bthead{每月} \\
\midrule
\textbf{工资到手} & 税后打到卡上的 & 19,000 \\
\textbf{公积金进账} & 单位 12\% + 个人 12\%，基数 2.6 万 & 6,240 \\
\textbf{房贷月供} & 刚性。剩 25 年，执行利率 3.35\% & −12,170 \\
\textbf{给父母} & 刚性 & −2,000 \\
\textbf{保险} & 刚性。一份 316 元/年的意外险，摊到月 & −26 \\
\textbf{弹性支出} & 吃饭通勤、健身卡、衣服、旅游摊销 & −5,800 \\
\textbf{结余} & 储蓄率 20.8\% & 5,244 \\
\bottomrule
\end{xltabular}
\end{tablewrap}

顺便说个真实细节：她 2021 年签的合同利率是 5.65\%，月供 15,585。经过 2023 年 9 月和 2024 年 10 月两轮存量房贷利率批量下调，现在执行 3.35\%，月供 12,170。这三千多块钱是政策给的，她自己什么都没做，甚至一开始没注意到。

更关键的是那个 6,240。公积金是她的收入，但她感觉不到，因为那笔钱直接进了另一个账户，还没经过她的手就被逐月提取去还了贷。\hlmark{感觉不到的收入，往往也是唯一存下来的收入。}

\begin{pullquote}{}
收益率不归你管，储蓄率归你管。它是这本书里唯一一个你能百分之百控制的变量。
\end{pullquote}

\begin{mythbox}{常见误解}
\mvfrow{误解}{"先想办法投资，光靠攒钱太慢了。"}
\mvfrow{事实}{算一遍就知道谁快。假设年化 5\%，每月存 2,500 块，存十年，终值 38.8 万。\\ 现在给你两个选项。\textbf{选项一：把收益率从 5\% 提到 8\%。}这件事极难，需要你连续十年跑赢绝大多数专业机构。做到了，终值 45.7 万，多出 \textbf{6.9 万}。\\ \textbf{选项二：把月存从 2,500 提到 5,200。}这件事只需要你重新排一遍支出，收益率还是那个 5\%。终值 80.7 万，多出 \textbf{41.9 万}。\\ 6.9 万对 41.9 万。前十年，储蓄率的权重压倒性地高于收益率。等你的本金攒到七位数，两者的地位才开始互换。}
\end{mythbox}

\section{三个数，自己量一下}

下面这几个门槛是经验值，不是定律，没有哪条监管规定写着 40\%。它们的作用像体检报告上的参考区间：超出去不代表你有病，代表该做个进一步检查。

\textbf{应急储备，按刚性支出算，不按总支出算。}林悦的刚性是月供 12,170 加赡养 2,000 加保费 26，再加上吃住通勤里怎么也省不掉的 2,800，大约 17,000。三个月 5.1 万，六个月 10.2 万。她手上活期 8,000 加货基 4.2 万，正好 5 万，卡在下限。

为什么用刚性支出？因为失业的第二个月你不会去旅游，不会续健身卡，弹性支出的定义就是它能归零。用总支出算会把这个数虚高一截，看着安心而已。

放的地方也有讲究：得是当天能到账的。货基快赎单日有一万的额度上限，七天通知存款、国债逆回购到期都行。不能放在 R2 以上的理财里，更不能放在需要 T+7 才回到卡上的产品里。急用钱的时候差的往往不是钱，是七天。

林悦另外那 11 万股票基金呢？那不是安全垫。那是一笔在你最需要用钱的时候，大概率也正好在跌的东西。经济不好的时候公司裁员和市场下跌是同一件事的两个面。

\textbf{月供收入比，看你怎么定分母。}林悦 12,170 除以到手的 19,000，是 64\%，吓人；把公积金算进收入，12,170 除以 25,240，是 48\%。同一个人，同一笔月供，两个结论差一倍。

这个指标最容易被自己骗过去。经验上超过 40\% 会开始难受，超过 50\% 基本没有抗风险能力，一次降薪或者三个月的空窗期就能击穿。林悦在 48\%，她自己觉得"还行"，因为公积金那部分她感觉不到。感觉不到不等于不存在。

\textbf{净值增长率，比收入增长率诚实。}今年涨薪 12\% 是件好事，但它不告诉你任何关于财富的信息，你完全可能同时把支出提高了 15\%。净值从年初的 92 万变成年末的 105 万，涨了 14\%，这个数才说明问题。它还有个好处：它会诚实地把房价跌的那部分写进去，不给你留面子。

\section{你最大的一笔资产，任何对账单上都没有}

林悦到手一年 23 万上下。她 32 岁，就算干到 60 岁，还有 28 年可以卖。这笔未来的收入有个名字，叫\term{人力资本}{human capital}\index{044@人力资本}。

\begin{egbox}{给自己估个价}
用第 6 章那套折现的做法：未来 28 年，每年 23 万，按 4\% 折回今天。

年金现值系数 =（1 − 1.04 的 −28 次方）÷ 0.04 = \textbf{16.66}。

23 万 × 16.66 = \textbf{383 万}。

折现率换成 3\%，这个数变成 432 万；换成 5\%，变成 343 万。区间不小，但结论不受影响：\textbf{383 万比她那套 328 万的房子还大，是她 105 万净值的三倍多。}

她全部资产里最大的一项，是她自己。而这一项没有对账单，没有净值曲线，也没人每天提醒她它今天跌了多少。

\end{egbox}

\begin{pullquote}{}
对绝大多数三十来岁的人来说，最大的一笔资产是自己，而且是唯一一笔没上保险的资产。
\end{pullquote}

把人力资本当成一只债券来看待，很多事情会立刻清楚。这只债券每月付息 1 万 9，到 60 岁那天本金归零，中间存在违约的可能。比喻可以继续推下去：面对一只有违约风险的债券，你会做两件事，一是给它买保险，二是不要在组合里再堆一大堆和它高度相关的东西。

第一件事推出的结论是：\textbf{年轻人真正的最大风险从来不是股票跌了。}林悦的基金全亏光是 11 万；她要是因为一场病三年不能工作，损失是三年 69 万，加上回来之后议价能力的下降。量级差六倍以上。可她花在挑基金上的时间，是花在体检和保单上的几十倍。这个精力分配是反的。

第二件事推出的结论更有意思：\textbf{你的人力资本长得像债券，还是像股票？}

\begin{tablewrap}
\begin{xltabular}{\linewidth}{>{\raggedright\arraybackslash}p{0.20\linewidth}XX}
\toprule
\bthead{职业} & \bthead{现金流的样子} & \bthead{和股市的相关性} \\
\midrule
\textbf{公务员、电网、三甲医院医生} & 稳定、可预测、涨得慢 & 接近零，像 AAA 级长期债券 \\
\textbf{企业 SaaS 产品经理（林悦）} & 看客户预算，客户预算看资本开支 & 中高，偏科技板块 \\
\textbf{券商营业部、投行、量化} & 底薪低，奖金和市场同呼吸 & 极高，像一只高贝塔股票 \\
\bottomrule
\end{xltabular}
\end{tablewrap}

在券商工作的人重仓券商股，等于把工资、奖金、股票账户、期权全押在同一个变量上。市场好的时候四个数一起涨，人会真心觉得自己有天赋；市场转向的时候，部门缩编、奖金归零、持仓腰斩会在同一个季度里发生。2015 年下半年和 2018 年，这种事在陆家嘴和金融街成批出现过。

林悦的客户预算和 AI 资本开支高度绑定。她的人力资本偏股票，还偏科技股。那 11 万基金里如果有 7 万是 AI 主题，她是在同一件事上加倍下注。

所以我给一个判断，而且认为它几乎没有例外：\textbf{大多数人的组合里最该被减配的，是自己公司的股票和自己行业的主题基金。}这跟看不看好那个行业无关，跟你已经在里面押了多少无关，跟你已经押了 383 万有关。

\section{四个盒子，按日期分，不按胆量分}

钱不该混在一个池子里。分开不是为了仪式感，是因为不同的钱有不同的到期日。

\begin{tablewrap}
\begin{xltabular}{\linewidth}{>{\raggedright\arraybackslash}p{0.20\linewidth}XXX}
\toprule
\bthead{账户} & \bthead{装什么} & \bthead{期限} & \bthead{放什么} \\
\midrule
\textbf{日常} & 一个月的开销 & 30 天内 & 活期、余额宝 \\
\textbf{安全垫} & 3–6 个月刚性支出 & 随时可能用 & 货基、七天通知存款、逆回购 \\
\textbf{目标} & 三年内确定要花的：首付、生育、读书、换车 & 1–3 年 & 定存、短债、同业存单指数 \\
\textbf{长期} & 十年不打算动的 & 10 年以上 & 权益类为主 \\
\bottomrule
\end{xltabular}
\end{tablewrap}

\begin{pullquote}{}
期限决定了你能承受多少波动，那张风险测评问卷决定不了。
\end{pullquote}

推一步就明白了。一个在问卷上被判定为"积极型"、但两年后要付 80 万首付的人，那 80 万碰都不该碰权益，因为两年不够一次熊市走完。一个被判定为"保守型"、但那笔钱三十年不会动的人，全放定存反而是在温水里慢慢亏掉购买力。同一份问卷，同一个人，两笔钱，答案完全相反。

\section{保险不是理财，保险是防破产}

这一节的全部内容就是一句话：\textbf{保险买的是"倒霉的时候不至于破产"，不是买收益。}

排序的原则也简单，先堵最可能发生并且最贵的那个窟窿。医疗险一年几百块，管大额住院，发生概率最高，放第一位。意外险一年一两百，杠杆最高，放第二位。重疾险赔一笔现金，用来填生病期间没有收入的那段，年缴几千，放第三位。定期寿险排最后，但它对某些人最要命。

\term{定期寿险}{term life insurance}\index{011@定期寿险}只在被保人身故时赔付，赔给还活着的人。它在中国被严重低估，因为买它的人自己一分钱好处都得不到。

林悦要是明年结婚生孩子，房贷还剩两百四十来万。她出事，这 247 万不会消失，会落在伴侣和孩子头上。\textbf{有房贷、有孩子的人，定期寿险保额至少要覆盖房贷余额。}她查了几家，32 岁女性、300 万保额、保到 60 岁、缴 20 年，年缴在 1,900 到 2,600 之间。一年两千块换掉一个 300 万的窟窿，是她整张表上性价比最高的一笔支出。

\begin{warnbox}{用保险做投资之前，先算清楚三件事}
\textbf{一，你交进去的钱有多少真的进了账户。}初始费用、保障成本、账户管理费会先拿走一部分，前几年拿走的比例通常最高。

\textbf{二，写进合同的保证利率是多少。}演示利率和保证利率是两个词，差一个字，差很多年。销售讲的通常是前者，赔付时按的是后者。

\textbf{三，什么时候能取，取的时候扣多少。}很多产品前五到十年退保拿回来的钱少于你交进去的。

这三个数全在合同里，全都不在话术里。如果有人花二十分钟讲收益、一分钟带过费用，你至少知道了他的重点在哪儿。

\end{warnbox}

\section{2026 年的算术题：要不要提前还贷}

\begin{statrow}{3}
  \statitem{3.5}{\%}{2026 年 8 月 20 日，5 年期以上 LPR}
  \statitem{1.4}{\%}{7 天期逆回购政策利率}
  \statitem{1.41}{\%}{商业银行净息差（2026 年二季度末）}
\end{statrow}

净息差 1.41\% 是历史低位。它的含义很直白：银行赚的存贷利差已经薄得像纸。利差压到这个程度，它就没有空间给你高存款利率了。所以 2026 年那种"我想找个稳当的地方拿 4\%"的愿望，基本无解。凡是在这个环境里同时承诺 4\% 以上和"保本"的东西，第 23 章会告诉你它的零件是怎么装的。

在这样的利率环境里，提前还贷成了很多人的常规问题。这里不给结论，给一套算法。

\begin{flowlist}
  \flowitem{01}{先查清楚你的房贷现在到底是多少利率}{不是合同上签的那个数。2023 年 9 月和 2024 年 10 月两轮存量利率调整之后，很多人的实际执行利率已经落到 3.1\%–3.5\% 区间，自己却还在按五年前的印象做判断。打开手机银行查重定价后的执行利率，五分钟的事。}
  \flowitem{02}{再找出你的"确定性收益率"}{不是你希望赚到的，是确定拿得到的：三年期定存、货基七日年化、国债逆回购。2026 年 8 月这一档大致在 1.4\% 到 2\% 之间。别把股票的期望收益放进这一格，期望不是确定。}
  \flowitem{03}{两个数相减}{3.35\% 减 1.8\%，差 1.55 个百分点。这才是提前还贷每年真正给你省下的东西，不是 3.35\%。很多人在这一步把收益算成了两倍。}
  \flowitem{04}{算完先别按确认}{还完之后你的安全垫还剩几个月？房子不能切一块下来救急，还进去的钱想拿回来只能重新申请贷款，而你需要它的那个时刻，往往正好是银行不愿意贷给你的时刻。\textbf{流动性只在你不需要它的时候，看起来像是多余的。}}
  \flowitem{05}{最后一格叫心理成本，没有公式}{月供从 12,170 降到 9,000，晚上睡得踏实一点，这件事值多少钱。有人填 0，有人填得比那 1.55 个百分点还大。两种人都没填错。}
\end{flowlist}

你的利率、你的备用金厚度、你工作的稳定性，我一样都不知道。\hlwarm{这一节不构成投资建议，也不构成还贷建议。}任何人在不知道这三件事的前提下告诉你"该不该提前还贷"，他给你的不是建议，是内容。

\bookbreak

最后说回动手这件事。这张表第一次做，会花掉你一个晚上，因为你得去翻公积金余额、查贷款剩余本金、找小区最近的成交价。第二次做只要二十分钟。

林悦做完之后改了三件事：把货基加到 8 万，买了那份定期寿险，把 AI 主题的仓位砍掉一半换成宽基。她的收益率没有因此提高，她只是不再需要靠猜测来回答"我到底有多少钱"。

\begin{takeaway}{本章一句话}
净值是算出来的，不是感觉出来的。而这张表上最大的一项，是你自己。

\begin{itemize}
  \item 房子按市价记，房贷按余额记，中间那个差才是你的，它可以是负数。
  \item 收益率不归你管，储蓄率归你管。前十年，后者的权重压倒性地大。
  \item 应急储备按刚性支出算 3–6 个月，放在当天能到账的地方。
  \item 32 岁的 383 万人力资本，比 105 万净值更需要被保护，也更容易被重复押注。
\end{itemize}

\end{takeaway}

\begin{bridgebox}
\textbf{下一章}表做完了，下一个问题绕不过去：那些数字要变大，钱到底得从哪儿来？
\end{bridgebox}


\chapter{收益从哪儿来：三种钱和一个残酷等式}
\chaptermeta{20}

\begin{ledetext}
你账户里多出来的那几万块，来路一共只有三条。分不清是哪一条，你就永远不知道自己是靠本事还是靠运气。
\end{ledetext}

\section{钱只有三个来路}

上一章那张表做完了，接下来的问题很朴素：怎么让净值那一栏变大。

抛开所有花哨的说法，你从市场上拿走的每一块钱，只可能来自三个地方。

\begin{flowlist}
  \flowitem{01}{经济增长赚的钱}{企业卖出东西、赚到利润、拿一部分分红、剩下的再投进去扩产。这块钱是新造出来的，蛋糕变大了，所有股东按份额一起分。它是正和的。来源说到底只有两样：生产率和人口。}
  \flowitem{02}{别人口袋里的钱}{你 8 块买 12 块卖，赚了 4 块。这 4 块不是凭空长出来的，它对应着某个人 12 块买了、8 块割了。蛋糕没变大，只是重新切了一遍。它是零和的。}
  \flowitem{03}{被摩擦吃掉的钱}{管理费、托管费、申赎费、佣金、买卖价差、印花税、大单进出时把价格打歪造成的冲击成本。这块钱既不来自增长也不来自对手方，它直接离开了所有参与者，去了基金公司、券商和国库。它是负和的，而且每年都收，不管你赚不赚。}
\end{flowlist}

三条路，性质完全不同。第一种可以持续，第二种不能保证，第三种保证发生。

\begin{pullquote}{}
投资这件事，最有用的一个问题是：我赚的这笔钱，到底是哪一种？
\end{pullquote}

\begin{egbox}{把一年的收益拆开看}
老赵在合肥开了家做汽车线束的小厂，业余炒股。2025 年他证券账户赚了 \textbf{4.27 万}，逢人就说自己波段做得好。

拆开：他常年拿着的 22 万宽基指数基金，那年净值涨加分红，贡献 \textbf{2.01 万}（第一种钱）。他另外用 15 万本金做了 11 次短线，来回赚了 \textbf{2.89 万}（第二种钱，对面是那 11 组卖给他和买他手上货的人）。全年佣金、印花税、基金申赎费合计 \textbf{6,300 元}（第三种钱，走了）。

2.01 + 2.89 − 0.63 = 4.27 万。

问题来了：那 2.89 万里，有多少是本事，有多少是运气？11 次样本，答不出来。同样的操作再来五年，如果还是稳定赚，才有一点点说服力。\textbf{第一种钱可以复制，第二种钱只能事后统计。}老赵把两者混在一起，得出了"我会做波段"这个结论。

\end{egbox}

\section{一个不需要任何假设的等式}

1991 年，威廉·夏普写了一篇一页半的东西，叫《主动管理的算术》。全文没有一个统计检验，没有一个假设，只有加法和除法。它的结论至今没人能反驳。

推理是这样走的。

\begin{chainflow}
\chainnode{所有人的持仓加起来就是整个市场}\chainarrow \chainnode{被动投资者按市值权重持有}\chainarrow \chainnode{剩下的主动投资者合起来也必然是市值权重}\chainarrow \chainnode{所以扣费前，主动的平均收益 = 市场收益}\chainarrow \chainnode{主动的成本更高}\chainarrow \chainnode{扣费后，主动的平均收益必然低于指数}
\end{chainflow}

关键在第三步。市场上的股票总共就那么多股，被动那一堆按市值比例拿走了一块，剩下的所有股票落在主动投资者手里，这一堆的构成按定义也是市值权重的。它们不可能整体上"超配"某只股票，因为没有第三方能提供额外的股份。

所以在扣掉成本之前，全体主动投资者作为一个整体，收益率和市场一模一样。不是接近，是相等。然后他们比被动投资者多付了研究费、管理费、换手成本。

减一个正数，只能变小。

\begin{egbox}{一个只有三个人的市场}
某个村子里只有一家公司，总股本 1000 股，年初每股 10 元，市值 1 万元。全村只有老张、老王、老李三个股东。

年底每股 11 元。整个市场涨了 10\%，这是公司经营的结果，跟谁买谁卖没关系。

这一年他们互相买卖了很多次。老张在 11.5 元卖给老王 200 股，赚了个差价；老王吃了亏。但年底一算：三个人合计还是持有 1000 股，合计价值 1.1 万元，\textbf{三人合起来的收益率就是 10\%，一分不多一分不少}。老张多赚的每一分，都是老王少赚的那一分。

再算摩擦。这一年他们交易了 8 次，卖出方的印花税千分之一，双边佣金每笔最低 5 元，合计 \textbf{217 元}离开了村子，进了券商和国库。

三人的真实合计收益是 1000 − 217 = \textbf{783 元}，年化 7.83\%。而村里那个从头到尾没动过、也没人给他讲行情的老李，如果一直满仓不动，他拿到的就是 10\%。

\end{egbox}

\begin{deepbox}{进阶：可以跳过}
这个结论最狠的地方在于它\textbf{不依赖任何关于市场是否有效的假设}。哪怕市场荒谬到人人都在拿彩票号码选股，这个算术照样成立，因为它只用到了"所有人的持仓加起来等于市场"这一条，而这一条是定义，不是观察。

你可以说"这次不一样"，但你得先反驳加法。

它同时也说明了另一件事：它讲的是平均数，不是每个人。一定有人长期跑赢，只是他赢的每一块钱，都必须来自另一个主动投资者输掉的那一块。零和游戏里赢家的存在，恰恰要求输家的存在。\hlmark{你在坐上牌桌之前，最好先想清楚自己打算当哪一个。}

\end{deepbox}

\section{长期收益的真正来源，是你替别人难受}

如果第二种钱靠不住、第三种钱一直在流出，长期回报到底从哪儿来？

答案是第一种钱，而它有一个更精确的名字：\term{风险溢价}{risk premium}\index{015@风险溢价}。

逻辑很朴素。有人手里有钱，但不想承受本金腰斩、不想把钱锁十年、不想承担对方还不上的可能、不想面对"急用钱时卖不掉"。这些不想，都得有人来替他们承担。替他们承担的人，收一笔报酬。

\begin{tablewrap}
\begin{xltabular}{\linewidth}{>{\raggedright\arraybackslash}p{0.20\linewidth}XX}
\toprule
\thead{溢价} & \thead{你替别人承担了什么} & \thead{2026 年的参照} \\
\midrule
\textbf{股权风险溢价} & 企业可能破产、利润可能下修、股价可能腰斩，而且不知道多久回来 & 历史长期在 3–5 个百分点，具体数字学界吵了三十年，至今没定论 \\
\textbf{期限溢价} & 把钱锁上十年，中间利率一动你就浮亏 & 美国 10 年期国债收益率 2026 年 8 月底升至 2025 年 1 月以来最高 \\
\textbf{信用溢价} & 借钱的那一方可能还不上 & 私募信贷借款利率 8\%–12\%，比银行贷款高 200–400 个基点 \\
\textbf{流动性溢价} & 想卖的时候没有买家，或者只能打折卖 & 私募信贷的二顺位与夹层 12\%–16\%，代价是几年不能赎回 \\
\bottomrule
\end{xltabular}
\end{tablewrap}

\begin{pullquote}{}
你赚的不是"聪明"的钱，是"承担别人不愿意承担的东西"的补偿。
\end{pullquote}

这个视角一旦装上，很多事情当场想通。为什么长期看股票比债券收益高？不是因为股票有魔力，是因为持有股票的过程更难受，2008 年跌 70\% 的时候你得拿得住，而拿不住的人把补偿让给了拿得住的人。为什么私募信贷能给你 10\%？因为它有三年不让你赎回，那多出来的几个点是你为"不能反悔"收的钱。

反过来推也成立：\textbf{一个既不用锁定期限、又没有波动、还能给你远高于无风险利率回报的东西，逻辑上不存在。}它要么在某处藏了风险没告诉你，要么它根本就是第 23 章要拆的那种东西。

\section{阿尔法这个词已经被用坏了}

行业里把收益拆成两块。跟着市场一起动的那部分叫\term{贝塔}{beta}\index{003@贝塔}，市场之外多出来的那部分叫\term{阿尔法}{alpha}\index{001@阿尔法}。前者理论上很便宜就能买到，后者才值得付管理费。

麻烦在于"市场"的定义一直在扩张。

1964 年那版模型里，市场只有一个变量。1993 年 Fama 和 French 加进了规模和价值两个\term{因子}{factor}\index{075@因子}，一大批原本被称作阿尔法的收益当场被重新归类为贝塔：那些基金经理不是选股厉害，他们只是长期偏爱小盘和低估值的股票。后来动量、质量、盈利、投资又陆续被加进来，每加一个，就有一批"能力"被翻译成"暴露"。

我给一个判断：\textbf{2026 年市面上被包装成阿尔法的东西，大部分是还没被识别出来的贝塔。}不是造假，是分不清。分不清的代价，是你按阿尔法的价格买了贝塔。

更麻烦的是因子本身在衰减。

想象一家藏在巷子里的面馆，好吃又便宜，因为没什么人知道。有人发了篇测评，三个月后门口排两小时，老板涨了五块钱。面还是那碗面，性价比没了。

因子就是这么死的：论文发表，做成产品，几千亿资金照着同一套信号买进去，便宜的东西不再便宜。往下推一步，结论有点残酷：\hlmark{任何一个被广泛讨论的策略，等你听说它的时候，它最好的那几年大概率已经过去了。}这不是谁在使坏，这就是价格干活的方式。

拥挤还带来第二个后果。所有人用同一套信号，进场时互相抬价，出场时互相踩踏。因子的超额收益变薄了，波动却变大了。这是 2018 年之后量化行业反复经历的事情。

\section{三十年的过路费}

这一节是全章最有杀伤力的部分，只有一道算术题。

本金 30 万，年化毛收益 7\%，投 30 年不动。唯一的变量是费率。

\begin{tablewrap}
\begin{xltabular}{\linewidth}{>{\raggedright\arraybackslash}p{0.20\linewidth}Xrr}
\toprule
\thead{情形} & \thead{净年化} & \thead{30 年后终值} & \thead{被费用拿走} \\
\midrule
\textbf{不收任何费用（理论值）} & 7.00\% & 228.4 万 & 0 \\
\textbf{宽基指数基金，费率 0.15\%} & 6.85\% & 218.9 万 & 9.5 万 \\
\textbf{主动型基金，费率 1.5\%} & 5.50\% & 149.5 万 & 78.9 万 \\
\bottomrule
\end{xltabular}
\end{tablewrap}

\begin{barfig}
  \barrow{零费用}{100}{228.4 万}
  \barrow{费率 0.15\%}{96}{218.9 万}
  \barrow{费率 1.5\%}{65}{149.5 万}
\end{barfig}

218.9 万和 149.5 万，差 \textbf{69.4 万}，占了高费率那一栏终值的将近一半，占了低费率终值的三成二。

把 69.4 万换个说法：它是本金 30 万的 2.3 倍。你投进去 30 万，管理人从中拿走了 69 万。

合同上写的差距只有 1.35 个百分点。1.35\% 看起来像个零头，很多人在签字的时候连问都不问。

\begin{pullquote}{}
这不是收益率的差别，这是复利被人在门口收了三十年过路费。
\end{pullquote}

要说公道话：如果一只基金能持续跑赢指数 3 个百分点，1.5\% 的费率是便宜的。问题是持续跑赢 3 个百分点的人极少，而且你在买入的那一刻无法辨认他是谁。\hlwarm{这段不构成投资建议}，它只是在说，费率是你事前唯一能确定知道的数字，收益率不是。

\section{为什么亏钱的总是散户}

下面几条，我自己全中过。写出来不是为了嘲笑谁。

\textbf{赚了就跑，亏了就扛。}这个毛病学名叫\term{处置效应}{disposition effect}\index{009@处置效应}。涨 8\% 的时候想落袋为安，跌 30\% 的时候说"反正没卖就不算亏"。后半句在会计上是错的：你的账户余额已经变了，没卖唯一改变的是你的心情。结果是你手里长期留着最差的那些，卖掉的全是最好的那些。

\textbf{过度自信。}各国的调查里，大约八成司机认为自己车技高于平均水平。这在数学上做不到。股票账户里的版本是：亏了怪市场，赚了归自己。

\textbf{买你看到的，不是买你研究的。}热搜、涨停板、朋友圈里的收益截图。这里有一个不对称：没人晒亏损截图。所以你看到的样本从第一眼开始就是歪的，你以为自己在观察世界，其实在观察幸存者。

\textbf{手一直在动。}换手率越高，成本越高，收益越低。这一条在各国的散户账户研究里，几乎是唯一一个稳定复现的结论。它跟你选股水平无关，它是算术。

这四条有个共同点：它们全都发生在你觉得自己很理性的时候。

\section{2026 年，你站在哪儿}

\begin{databox}{2026 数据}
\textbf{全球被动投资规模已首次超越主动管理。}这件事花了五十年。夏普那篇一页半的文章写于 1991 年，市场用了三十五年才开始按它说的做。

\textbf{中国公募基金规模约 38 万亿元}（2026 年初数据）。这个盘子已经大到，公募之间互相博弈的部分远大于它们从散户手里赚走的部分。

\textbf{A 股：2026 年 8 月 31 日收盘，上证指数 3986.30 点，当月上涨 4.02\%，两市成交约 2.13 万亿元。}市场的核心逻辑正从"估值修复"转向"盈利驱动"。

\end{databox}

最后那句话如果翻译成人话，正好落回本章开头。

\begin{talkbox}
  \talkline{研报}{"市场逻辑正从估值修复转向盈利驱动。"}
  \talkline{翻译}{过去这一段，股价涨是因为大家愿意为同样的一块钱利润多付几倍钱（第二种钱，重新定价，有人接手就涨，没人接手就停）。往后要涨，得靠那一块钱利润本身变成一块二（第一种钱，蛋糕真的变大）。}
  \talkline{所以}{估值能修复到哪儿，没人知道，那取决于别人愿意付多少。盈利能不能兑现，财报里会有答案，虽然要等。}
\end{talkbox}

\begin{mythbox}{常见误解}
\mvfrow{误解}{"投资就是找到会涨的股票。"}
\mvfrow{事实}{拉长到十年以上看，宽基指数的累计收益里绝大部分来自企业盈利增长和分红再投入，估值变化的贡献要小得多。"找到会涨的股票"这句话里藏着一个没说出口的前提：会有人在更高的价位从你手里接走。\textbf{研究"哪家公司能赚钱"和研究"下个月谁愿意出更高价"，是两门完全不同的手艺。}很多人以为自己在做前者，账户的换手率出卖了他们。}
\end{mythbox}

\begin{mythbox}{常见误解}
\mvfrow{误解}{"基金经理是专业的，肯定比我强。"}
\mvfrow{事实}{他在信息处理和纪律上多半确实比你强。但"比你强"和"能持续跑赢同类平均、并且把 1.5\% 的费用赚回来"，是两个命题。第一个成立不代表第二个成立。\\ 真正卡住你的是第三件事：\textbf{即使确实存在少数能持续跑赢的人，你在事前挑不出他。}用过去三年排名去挑，前四分之一的基金在下一个三年继续留在前四分之一的比例，长期统计下来接近抛硬币。而销售页面上给你看的，永远是过去三年的排名。}
\end{mythbox}

\begin{warnbox}{小心}
本章讲的全是机制，不是标的。哪个指数会涨、哪只基金明年好、A 股 3986 点是高是低，这些问题本书一个都不回答，也建议你对任何回答得很干脆的人保持警惕。\hlwarm{对机制可以下判断，对价格不能。}

\end{warnbox}

\begin{takeaway}{本章一句话}
你赚的每一块钱只有三个来路：经济增长、别人的口袋、以及正在被摩擦吃掉。第一种可以持续，第二种是零和的，第三种保证发生。

\begin{itemize}
  \item 扣费前主动投资者的平均收益必然等于市场，扣费后必然低于指数。这是恒等式，不是观点。
  \item 长期收益来自承担别人不愿承担的东西：波动、期限、违约、卖不掉。
  \item 30 万投 30 年，费率 0.15\% 和 1.5\% 的终值差 69.4 万，是本金的 2.3 倍。
  \item 费率是你事前能确定的唯一数字。
\end{itemize}

\end{takeaway}

\begin{bridgebox}
\textbf{下一章}既然长期收益主要来自承担风险，那有没有办法在不降低收益的前提下少担一点？有，而且只有这一个。
\end{bridgebox}


\chapter{资产配置：唯一的免费午餐}
\chaptermeta{21}

\begin{ledetext}
金融学里几乎所有东西都要你拿什么去换。只有一件事，收益不减，风险变小。它的条件比多数人以为的苛刻。
\end{ledetext}

\section{那顿饭长什么样}

上一章的结论是：长期回报来自承担风险。听着像是没得选，想要收益就得挨揍。

有一个例外。

假设有两个资产，长期收益都是 8\%，但它们涨跌的节奏对不上。各买一半，每年调回一半一半。看看会发生什么。

\begin{tablewrap}
\begin{xltabular}{\linewidth}{>{\raggedright\arraybackslash}p{0.20\linewidth}rrr}
\toprule
\thead{年份} & \thead{资产 A} & \thead{资产 B} & \thead{各一半，年末调回} \\
\midrule
\textbf{第 1 年} & +32\% & −9\% & +11.5\% \\
\textbf{第 2 年} & −18\% & +24\% & +3.0\% \\
\textbf{第 3 年} & +25\% & −5\% & +10.0\% \\
\textbf{第 4 年} & −11\% & +22\% & +5.5\% \\
\textbf{第 5 年} & +12\% & +8\% & +10.0\% \\
\textbf{\textbf{算术平均}} & 8.0\% & 8.0\% & 8.0\% \\
\textbf{\textbf{实际年化}} & 6.17\% & 7.15\% & 7.95\% \\
\textbf{\textbf{波动（标准差）}} & 19.6\% & 13.5\% & 3.2\% \\
\bottomrule
\end{xltabular}
\end{tablewrap}

10 万块放进去五年：A 变成 13.49 万，B 变成 14.12 万，各一半变成 \textbf{14.66 万}。

组合比它的两个成分都赚得多。这不是错觉，也不是我算错了。

原因在算术平均和实际年化之间那道缝。跌 18\% 之后要涨 22\% 才回本，波动越大，这道缝越宽，实际拿到手的越少于纸面上的平均数。把波动摁下去，缝就窄了，同样的 8\% 平均收益，你能带走更多。

\begin{pullquote}{}
收益没有减少，风险减少了。这是金融学里唯一一件近似免费的事。
\end{pullquote}

\begin{warnbox}{先说清楚这个例子的水分}
上面那两组数字是我故意挑的，它们几乎完全反向，\term{相关性}{correlation}\index{070@相关性}接近 −1。现实中找不到这么听话的一对。

股票和债券的相关性在过去二十年里从 −0.4 到 +0.6 之间来回跑，不是一个常数。所以真实世界的免费午餐没这么大份，效果通常是把波动降低两三成，而不是从 19.6\% 降到 3.2\%。

方向是对的，幅度要打折。任何一本把这个例子讲得像永动机的书，都跳过了这一段。

\end{warnbox}

\section{关键变量是相关性，不是数量}

免费午餐的开关只有一个：两样东西是不是一起动。

如果两个资产永远同涨同跌，各买一半和全买一个，波动一模一样，什么都没省下。买十只、买五十只，同样什么都没省下。

\begin{mythbox}{常见误解}
\mvfrow{误解}{"我买了八只不同的基金，已经很分散了。"}
\mvfrow{事实}{老潘在武汉一家医院的设备科上班，2025 年陆续买了 8 只基金，名字分别叫科技创新、数字经济、高端制造、先进制造、智能汽车、新质生产力、成长精选、产业升级。他觉得自己布局得很周全。\\ 后来他把 8 只基金的前十大重仓拉到一张表上对齐，去重之后只剩 23 家公司；其中 6 家同时出现在 7 只基金里，合计占他总仓位的 \textbf{41\%}。2026 年 3 月这 6 家一起回调，他 8 只全绿，一只都没躲开。\\ \textbf{他不是买了 8 只基金，他是用 8 种方式买了同一个东西。}分散要看穿透之后的持仓，不看产品数量。}
\end{mythbox}

判断分散有没有效，有一个笨办法：回想最近一次市场大跌，你账户里有没有哪一栏是绿的旁边站着一栏红的。如果所有格子都是同一个颜色，你手上就只有一个资产，它只是被拆成了八个名字。

\section{每一类资产，在什么时候救你，在什么时候害你}

大类资产不该按"收益高低"排序，该按"它在什么天气下管用"排序。

\begin{tablewrap}
\begin{xltabular}{\linewidth}{>{\raggedright\arraybackslash}p{0.20\linewidth}XX}
\toprule
\thead{资产} & \thead{它在什么时候救你} & \thead{它在什么时候害你} \\
\midrule
\textbf{股票} & 经济扩张、企业盈利上行、通胀温和。长期分享增长的唯一入口 & 衰退、盈利下修、利率快速上行的时候。而且它跌起来不通知你 \\
\textbf{利率债（国债、政金债）} & 衰退和通缩恐慌里它涨，央行降息时它涨。提供现金流，也提供"别的东西都在跌"时的缓冲 & 通胀上行加加息，久期越长伤得越重。2022 年就是这么发生的 \\
\textbf{现金与类现金} & 所有人都缺钱的时候你有钱。它的价值不在利息，在于它是一张随时能出手的期权 & 通胀高的年份，它被安静地磨掉购买力。放太多是慢性失血 \\
\textbf{黄金与实物资产} & 货币信用被怀疑、地缘冲突、实际利率下行的时候 & 实际利率上行、风险偏好回暖的时候。它不生息，等待是有成本的 \\
\textbf{自住房} & 你得有地方住，这份居住价值天天在兑现。温和通胀下杠杆还能放大收益 & 你需要现金的时候。它卖不快、切不开，而且和你的工作往往在同一个城市里同生共死 \\
\bottomrule
\end{xltabular}
\end{tablewrap}

最后一栏那句得展开一句：房子的坏处不只是流动性差。你在杭州工作，房子也在杭州，如果杭州的主导产业出问题，你的收入和你的房价会一起往下走。第 19 章说的"不要重复押注同一个变量"，在这里同样成立，只是很少有人这么看自己的房子。

\section{60/40 的教训}

六成股票、四成债券，是过去三十年被写进无数教材和产品说明书的默认答案。它背后的逻辑是股债负相关：股票跌的时候债券涨，互相垫着。

2022 年，这套逻辑当场失效。

\begin{egbox}{2022 年发生了什么}
那一年标普 500 全收益 \textbf{−18.1\%}，美国综合债券指数 \textbf{−13.0\%}，经典 60/40 组合跌了 \textbf{16\%} 上下，是 1937 年以来最差的年份之一。

两条腿一起断，因为砍它们的是同一把刀。\textbf{通胀起来 → 央行加息 → 债券价格按久期下跌}（第 10 章那套算术），同时\textbf{折现率抬高 → 股票估值被压缩}（第 6 章那个除数）。同一个变量，两个方向的伤害。

持有 60/40 的人以为自己买了两样东西，2022 年才发现自己买的是同一个宏观赌注的两个版本。

\end{egbox}

这里要下一个判断，而且我认为它比大多数配置建议重要：\textbf{股债负相关不是自然法则，它是低通胀时代的产物。}

1998 到 2020 年那二十来年，冲击主要来自需求侧：需求塌了，企业盈利下滑，股票跌；央行降息救市，债券涨。负相关是这个因果链的副产品，不是资产本身的属性。

当冲击来自供给侧，链条反过来：能源涨价、关税加码、供应链断裂，物价上去了，央行不得不加息，股和债同向下跌。2026 年的世界里，供给侧的麻烦比过去二十年密集得多：中东局势推高油价并恶化贸易条件，关税与去全球化抬高了商品成本，AI 数据中心的资本开支在抢电力和产能。

在一个通胀中枢更高、供给冲击更频繁的世界里，只靠股债两条腿是不够的。这就是通胀对冲类资产在 2026 年被重新拿出来讨论的原因。

\section{央行在买黄金，但他们买的不是收益}

先把黄金的缺点说完，再说它为什么还在桌上。

黄金不产生现金流。一盎司金子放一百年，还是一盎司金子，不分红不生息。它的全部回报来自"下一个人愿意出多少钱"。长期看，它的实际回报低于股票，这一点没有争议。

\begin{mythbox}{常见误解}
\mvfrow{误解}{"黄金能保值，跌不到哪儿去。"}
\mvfrow{事实}{2011 年 9 月金价站上 1900 美元/盎司附近，到 2015 年 12 月跌到 1050 美元附近，\textbf{四年多跌掉四成半}，期间不产生一分钱利息。在 2011 年高点买入的人，要等到 2020 年才回到成本线。\\ 黄金对冲的是特定的风险：货币信用、地缘冲突、极端通胀。它不是"稳"的东西。把一件对冲工具当成保本工具，是 2013 年让很多人亏钱的那个误会。}
\end{mythbox}

说完缺点，说 2026 年最有说服力的那组证据。它不来自任何一位分析师的观点，来自各国央行的实际行为。

\begin{databox}{2026 数据：全球央行储备结构}
欧洲央行 2026 年 6 月《欧元的国际角色》报告：\textbf{截至 2025 年底，黄金已超越美国国债，成为全球央行的第一大储备资产。}

\begin{barfig}
  \barrow{黄金}{27}{27\%}
  \barrow{美国国债}{22}{22\%}
  \barrow{其他美元资产}{20}{20\%}
  \barrow{欧元}{15}{15\%}
\end{barfig}

黄金占比一年之间从 20\% 升到 27\%，美债从 25\% 降到 22\%，其他美元资产从 22\% 降到 20\%，欧元维持 15\%。美元计价资产合计仍占 \textbf{42\%}，仍是最大的单一币种，这一点没变。

\end{databox}

\begin{statrow}{3}
  \statitem{3.6}{万吨}{全球央行黄金储备，接近布雷顿森林体系鼎盛时期的约 3.8 万吨}
  \statitem{850}{吨}{2025 年全球央行净购金量}
  \statitem{5600}{美元}{2026 年 1 月金价一度接近的每盎司水平}
\end{statrow}

还有个有意思的细节：2025 年单一最大的黄金买家不是哪个国家，是稳定币发行商 Tether，一年买了超过 100 吨。一家发行数字美元的公司，把储备的一部分换成了地下室里的金属。

央行为什么买？\textbf{他们买的不是收益，是"不会被没收"。}一国的美债储备存在别人的清算系统里，别人可以按一个键把它冻住。这件事在 2022 年之后从假想题变成了现实题。金库里的实物金砖，放在自己国境内，谁也按不了那个键。

这个动机往下推一步，会推出一个对你有用的结论：\hlmark{央行的理由和你的理由不一样，所以他们的工具你不能照抄。}你在手机上买的黄金 ETF，登记在券商的账户里，依赖的恰恰是同一套清算系统。它能对冲通胀和货币贬值，但它对冲不了央行真正担心的那件事。你买它可以，但别把别人的动机当成自己的理由。

\hlwarm{这不构成投资建议。}本书不预测金价，2026 年 1 月的 5600 美元是一个已经发生的数字，不是一个方向。

\section{再平衡和定投：两条治人的规则}

配置比例定下来只是第一天的事。市场会替你改比例：涨得多的那一块自动变大，跌得多的自动缩小，一年以后你的 60/40 可能已经变成 71/29，而你什么都没做。

\term{再平衡}{rebalancing}\index{079@再平衡}就是把它掰回去。

\begin{flowlist}
  \flowitem{01}{按时间}{每年固定一天，全部调回目标比例。挑一个跟市场无关的日子，比如生日后的第一个工作日。别挑年底，年底所有人都在动，你的情绪会被新闻带着走。}
  \flowitem{02}{按阈值}{目标权益 60\%，涨到 65\% 或者跌到 55\% 就动手，中间不管。这个办法比按时间更贴合市场，代价是你得每个月看一眼。}
  \flowitem{03}{混合}{每季度看一次，只在偏离超过 5 个百分点时才操作。多数人用这个够了。}
  \flowitem{04}{用新钱补，别用卖出调}{每月的储蓄优先买进比例偏低的那一块。同样达到再平衡的效果，省掉申赎费和可能的税。这一条能省下的钱，往往比你纠结选哪只基金省下的多。}
\end{flowlist}

再平衡的价值不在于它聪明，在于它强迫你在下跌的时候买入、在上涨的时候卖出。这两件事人靠意志力做不到。每一轮下跌里基金赎回量都会放大，这是行业里公开的规律，说明大部分人在那个时刻做的恰好是相反的动作。

它是一台自动的高抛低吸机器，而且它不预测，只按比例干活。比喻可以继续推：什么时候这台机器会害你？当某个资产进入不可逆的单边下行时，机械再平衡会不断往一个正在归零的东西里加钱。所以\textbf{再平衡只能在大类资产层面做，不能在个股层面做}。一个国家的宽基指数不会归零，除非这个经济体消失；一家公司会。

说完再平衡说定投，这两件事经常被混为一谈，其实解决的是不同的问题。

\term{定投}{dollar-cost averaging}\index{012@定投}不提高预期收益。如果一个市场长期向上，一次性投入的期望收益高于分批投入，因为你的钱在场的时间更长。这个结论和很多人的直觉相反，但它就是算术。

定投降低的是另外两样东西：择时风险，你不会恰好把全部身家押在最高点；行为风险，市场跌的时候你还在买，因为扣款是自动的，不需要你每个月重新鼓一次勇气。

\begin{pullquote}{}
定投主要治的是人，不是市场。
\end{pullquote}

"微笑曲线"那套话术的问题在于它偷偷塞了个前提：市场一定会回来。

日经 225 在 1989 年 12 月见到 38915 点，之后花了 34 年，到 2024 年 2 月才第一次超越那个位置；2026 年 6 月首次突破 70000 点。从 1990 年开始定投日经的人，最后确实赢了，代价是熬三十四年，中间要眼睁睁看着账户长期浮亏并且每个月照扣不误。

所以定投真正的前提不是"曲线会微笑"，是两条：你买的那个东西背后的经济体还在增长，以及你真的能坚持三十年。第二条比第一条难得多。

\section{你的年龄，和中国投资者的三条绳子}

回到林悦。她 32 岁，人力资本 383 万，不含房的金融资产 25.6 万。她的"总资产"里 94\% 是人力资本。

那 25.6 万如果全部亏光，是她总资产的 6\%。同样一场跌 30\% 的熊市，落在她身上是总资产的 1.8\%，落在一个 62 岁、人力资本已经接近零的人身上，就是 30\%。

年轻人能承受更高的波动，不是因为胆子大，是因为后面还有几十年的现金流可以往坑里填。这就是\term{生命周期投资}{life-cycle investing}\index{047@生命周期投资}的全部道理。

但它有个前提，第 19 章那张职业表已经说过：如果你的人力资本本身就像一只高贝塔股票，这份额外的容忍度要打折。券商营业部里 28 岁的年轻人，未必比 45 岁的公务员能承受更多波动。

接下来是三条只在中国存在的绳子，绕不开。

\begin{egbox}{绳子一：房子已经把仓位占满了}
央行 2019 年城镇居民家庭资产负债调查显示，住房资产占城镇家庭总资产的 \textbf{59.1\%}。口径是 2019 年、城镇家庭，之后没有同口径的公开更新，但结构没有理由发生根本变化。

这意味着一件很尴尬的事：大多数中国家庭在打开基金 App 之前，已经默认满仓了一个加了杠杆、流动性极差、和自己收入高度相关的单一资产。

林悦的例子是现成的：房子 328 万，其余金融资产 25.6 万。她在那 25.6 万里怎么分配，对整体风险的影响不到一成。\textbf{对她来说，最有效的一次资产配置决策不在账户里，在于要不要买第二套。}

\end{egbox}

\textbf{绳子二：跨境配置是有额度的。}每人每年 5 万美元的便利化额度，QDII 基金的额度时紧时松、经常溢价，港股通的标的范围有限。教科书上那种"全球分散"的最优组合，在中国落地时会撞到一堵墙。承认这堵墙的存在，比假装它不在要有用。

\textbf{绳子三：你其实已经持有一大块债券了，只是它不叫债券。}林悦的公积金账户 9.6 万，加上养老保险个人账户里的余额，这些钱利率低、几乎没有波动、期限极长、强制缴纳。它的性质接近一笔久期很长的低息债券。

推论很实用：如果你已经被强制持有了这么一块类债券资产，你在金融账户里再配 60\% 的债，可能是在重复同一件事。至于该配多少，取决于你的公积金余额、你的年龄和你三年内要花的钱，这三个数只有你自己知道。

\begin{deepbox}{进阶：为什么没有通用比例}
本章从头到尾没有出现一个"建议配置比例"，这是故意的。\hlwarm{这不构成投资建议。}

原因不复杂：两个同样 35 岁、同样年入 40 万、在同一份风险测评问卷上得了同一档分数的人，一个两年后要付 90 万首付，一个未来十年没有任何刚性大额支出。他们能承受的波动差着十万八千里，而问卷把他们分进了同一格。

问卷问的是你的性格，配置该问的是你的日期表。任何人给你一个通用比例，都必须先问你两件事：这笔钱什么时候要用，以及你的刚性支出每月是多少。不问这两件事就给比例的，给的不是配置方案。

\end{deepbox}

\begin{takeaway}{本章一句话}
把涨跌不同步的东西放在一起，收益不减，波动变小。这是金融学唯一的免费午餐，而它的开关是相关性，不是数量。

\begin{itemize}
  \item 八只重仓同一批公司的基金，是一只基金的八个名字。分散看穿透持仓。
  \item 股债负相关是低通胀时代的产物，不是自然法则。2022 年两条腿一起断，因为砍它们的是同一把刀。
  \item 央行买黄金买的是"不会被没收"，你买的是账户里的一行数字。动机不同，别照抄。
  \item 再平衡强迫你反着人性做，定投治的是人不是市场。两者都不提高预期收益，都降低你搞砸的概率。
\end{itemize}

\end{takeaway}

\begin{bridgebox}
\textbf{下一章}配置好了，剩下的交给时间。可时间这件事，被误解得比什么都深。
\end{bridgebox}


\chapter{复利的真相，和关于复利的谎言}
\chaptermeta{22}

\begin{ledetext}
复利被讲成了理财界的鸡汤。它是真的，但它有五个前提，缺一个就不成立。而它对亏损的态度，比你想的狠得多。
\end{ledetext}

\section{先做一道小学算术}

你有 10 万块。第一年亏 50\%，第二年赚 50\%。

现在你有多少？

7 万 5。

10 万亏一半剩 5 万，5 万涨一半是 7 万 5。两年下来"平均收益率"是 0\%，你的钱少了四分之一。要回到 10 万，第二年得赚 100\%，不是 50\%。

你的账户不做加法，它只做乘法。所谓复利，算的是\term{几何平均}{geometric mean}\index{026@几何平均}。而几何平均对亏损的惩罚是不对称的：\textbf{亏得越深，回本所需的涨幅涨得越快，而且是加速涨。}

这张表建议你拍下来。

\begin{tablewrap}
\begin{xltabular}{\linewidth}{>{\raggedright\arraybackslash}p{0.20\linewidth}Xrr}
\toprule
\thead{你亏了} & \thead{剩下} & \thead{回本需要涨} & \thead{按年化 8\% 算，要爬多少年} \\
\midrule
\textbf{10\%} & 0.90 & 11.1\% & 1.4 年 \\
\textbf{20\%} & 0.80 & 25.0\% & 2.9 年 \\
\textbf{30\%} & 0.70 & 42.9\% & 4.6 年 \\
\textbf{50\%} & 0.50 & 100\% & 9.0 年 \\
\textbf{70\%} & 0.30 & 233\% & 15.6 年 \\
\textbf{90\%} & 0.10 & 900\% & 29.9 年 \\
\bottomrule
\end{xltabular}
\end{tablewrap}

最后一行读三遍。亏掉九成之后，就算你接下来每年都能稳拿 8\%，回到原点也要将近三十年。三十年是很多人的整个工作生涯。

\begin{egbox}{举个栗子}
汪工在成都一家做汽车线束的厂子里当工艺工程师。2015 年 6 月，他把攒了六年的 37 万 8 全部买进一只创业板主题基金，因为"身边人都在赚"。

后面的事你知道。他没割，一直拿着，最深的时候账户剩 14 万 3，亏掉 62\%。

62\% 意味着什么？意味着他需要涨 163\% 才回到 37 万 8。按每年 8\% 爬，要爬 12 年 7 个月。他那年 34 岁。

这不是"拿得住就行"能解决的问题。这是算术。

\end{egbox}

\section{8\% 和 8\%，不是同一个 8\%}

两个人都告诉你，他的长期年化是 8\%。

第一个是真的每年 8\%，稳得像一张持有到期的国债。第二个是隔年跳：今年 +30\%，明年 −14\%，来回摆。你把 30 和 −14 一加除以二，正好是 8。

各拿 32 万，各走三十年。

\begin{barfig}
  \barrow{每年稳稳 8\%}{100}{322 万}
  \barrow{+30\% 与 −14\% 交替}{53}{170 万}
\end{barfig}

差 151 万。而这两个人的"平均年化"在任何一份宣传材料上都会被写成同一个数字。

波动的那位，真实年化是 5.74\%，不是 8\%。因为 1.30 乘 0.86 等于 1.118，开平方是 1.0574。少掉的那 2.26 个百分点，被波动本身吃掉了。这个东西叫\term{波动拖累}{variance drag}\index{008@波动拖累}，有个粗略但好用的口算公式：\textbf{几何收益率约等于算术收益率减去方差的一半。}上面这个例子波动率是 22 个百分点，0.22 的平方除以 2 是 0.0242，8\% 减 2.42\% 得 5.58\%，跟真实值 5.74\% 只差一点点。

把它想成开车：同样从成都到重庆，一个走高速匀速跑，一个在山路上不停地急刹再地板油。里程一样，油耗差一截。可以继续推下去：既然波动本身就是成本，那么在收益率不变的前提下把波动压下来，等于凭空提高了收益。这正是第 21 章说资产配置是"免费午餐"的真正原因，它不是玄学，是这条公式。

\section{钱几乎全在后半段}

有个心算工具叫 72 法则：用 72 除以年化收益率，得到本金翻一倍需要的年数。6\% 就是 12 年，9\% 就是 8 年。误差在常见区间内很小。

拿 23 万，按 6\% 走。第 12 年 46 万 3，第 24 年 93 万 1，第 30 年 132 万 1。

现在把这条曲线切成两段看。

\begin{statrow}{2}
  \statitem{50.8}{万}{前 20 年一共长出来的钱}
  \statitem{58.3}{万}{第 21 到第 30 年长出来的钱}
\end{statrow}

最后十年的收成，超过了前二十年的总和。

复利的形状不是一条斜坡，是一根被慢慢压弯的鱼竿。前二十年你几乎看不出它在弯，最后十年它突然抬起来。这个形状能推出两个很实际的结论。

第一，早开始比高收益更值钱。

\begin{egbox}{举个栗子}
小林 26 岁，郑州一家社区医院的护士，每月定投 1800 块，一年 21600。她只投十年，36 岁那年停手，之后一分钱不再加，就那么放着。按 6\% 算，她 61 岁时账上大概 122 万。

她表哥 36 岁才开始，同样每月 1800，一直投到 61 岁，投了 25 年，本金掏出去 54 万。他 61 岁时账上大概 118 万 5。

小林掏了 21 万 6，表哥掏了 54 万，多掏了 32 万多。最后表哥还少几万。

差别不在收益率，两个人用的是同一个 6\%。差别在小林那笔钱多在曲线上待了十年，而那十年正好是鱼竿抬头的那一段。

\end{egbox}

第二，中途取钱的代价大得不像话。

还是那 23 万、那 6\%。假设第 12 年你抽走 15 万去付一笔首付。到第 30 年，账户不是少 15 万，是少 42 万 8。因为被抽走的不只是那 15 万，还有它后面 18 年的弯曲。

那 15 万的真实价格，是 42 万 8。

\begin{pullquote}{}
复利不奖励聪明。它奖励活得久、开始得早、以及最不像本事的那一项：不动手。
\end{pullquote}

\section{榜单上看不见的那些基金}

第一个谎言：把"历史年化 X\%"当成承诺。

"这只基金过去五年年化 19\%。"这句话通常是真的，它同时几乎是无用的。

排行榜是从\emph{还活着}的产品里排出来的。清盘的、被合并的、换了基金经理又换了名字重新出发的，都不在榜上。中国公募基金规模 2026 年初约 38 万亿元，几千只产品，每年都有一批悄悄退场。你看到的那张榜单，是一张幸存者的合影，照片里没有没能来的人。这叫\term{幸存者偏差}{survivorship bias}\index{073@幸存者偏差}。

更大的那个幸存者，是国家。

"美股长期年化 10\%"这句话背后，站着一个本土从未被现代战火犁过、打赢了两次世界大战、货币成了全世界储备资产的国家。同一段时间里，1917 年的俄国股市、1949 年的上海证券交易所、两次大战之间中欧的一批市场，投资者最终拿到的数字是零。

\textbf{把一个赢家的历史收益率，当成"资产类别的自然属性"，这个推理是错的。}它不是"股票长期会涨 10\%"，它是"在过去一百年里赢麻了的那个国家，股票涨了 10\%"。这两句话的区别，就是你未来三十年计划里最大的那个假设。

\begin{deepbox}{进阶：可以跳过}
归零的市场有个讨厌的性质：它们不进入任何长期收益率序列。因为序列需要数据，而归零之后就没有数据了。所以所有"长期看股票一定涨"的统计，天然只能用活下来的样本做。

这不是说统计没用。是说你在读那个数字的时候，得知道分母被谁悄悄改小了。

\end{deepbox}

顺便说说第三个谎言，因为它是同一个毛病：\textbf{"巴菲特靠的是复利。"}

他确实靠复利。但有两样东西很少被一起提。

一是杠杆。伯克希尔的核心是保险，保户交了保费、还没出险的那段时间，钱归公司支配，这笔钱叫\term{浮存金}{float}\index{016@浮存金}。它的成本常年低于国债利率，某些年份甚至是负的。学术界对他长期杠杆倍数的估算大约是 1.6 倍。用别人的钱、且几乎不要利息、且不会被追缴，这个条件本身就值一大截超额收益。

二是时间。他 11 岁买第一只股票，绝大部分财富是 50 岁之后堆起来的。他的复利跑道长度，是普通人的两倍还多。

说他"只靠复利"，是把一个幸存者的全部条件砍掉，只留下最容易复制的那一条给你看。

\section{复利的五个前提}

第二个谎言：复利适用于一切。

它需要同时满足五个条件。缺一个，公式就不成立了。

\begin{flowlist}
  \flowitem{01}{时间足够长}{至少二十年，最好三十年。如果你 47 岁才开始规划退休，鱼竿抬头的那一段你多半赶不上。这不是打击，是让你把预期挪到"提高储蓄率"上，那个变量还在你手里。}
  \flowitem{02}{中途不取钱}{结婚、买房、生孩子、父母住院、失业半年。普通人的三十年里，几乎必然有一次要动这笔钱。前面算过，第 12 年抽 15 万，终值少 42 万 8。}
  \flowitem{03}{收益能再投资}{分红要再投，票息要再投。拿去消费的那部分，不参与复利。一只年年高分红但你年年花掉的产品，你拿到的是单利。}
  \flowitem{04}{不破产}{任何一年出现 −100\%，前面所有年份的乘积一次性归零。加杠杆、押单一标的、玩合约的人，把这个条件的概率从"几乎不可能"提到了"迟早会碰上"。}
  \flowitem{05}{税和费吃不掉太多}{管理费、托管费、销售服务费、申赎费、换手带来的摩擦，加起来每年 1.5\% 不算夸张。8\% 变成 6.5\%，三十年后终值只剩原来的三分之二。你没做错任何事，钱少了三分之一。}
\end{flowlist}

\begin{mythbox}{常见误解}
\mvfrow{误解}{"我要求不高，年化 15\% 就行。"}
\mvfrow{事实}{2026 年 8 月，中国 7 天期逆回购政策利率 1.4\%，5 年期以上 LPR 3.5\%；美国联邦基金利率目标区间 3.50\%–3.75\%。要求 15\%，等于要求在无风险利率之上再拿 11 到 13 个百分点的补偿。这个补偿没有人会白给你，它对应的是你必须承担的风险。真正该问的不是"15\% 高不高"，而是"为了这 15\%，我要接受多大概率的一次 −40\%"。回到本章第一张表：−40\% 需要涨 66.7\% 才回本。}
\end{mythbox}

\section{同样的平均收益，倒过来就要命}

这一节讲一个中文世界里几乎没被讲清楚的东西：\term{序列风险}{sequence of returns risk}\index{074@序列风险}。

先说结论。在\textbf{积累期}，收益的先后顺序不重要，甚至前面跌还是好事，你每月定投的钱买到了更便宜的份额。在\textbf{提取期}，顺序决定生死。

郑老师，杭州人，教了三十二年中学化学，2026 年退休。账上 246 万，打算每年取 12 万补贴生活，取二十年。

给她两种未来。收益率的那二十个数字完全一样，只是顺序颠倒。

路径 A：头三年 −27\%、−17\%、−12\%，之后十七年每年 +10\%。\\ 路径 B：先是十七年每年 +10\%，最后三年 −27\%、−17\%、−12\%。

两条路的算术平均都是 5.7\%，几何年化都是 5.08\%，每年都取 12 万，二十年都取走 240 万。数学上，它们是同一组收益率。

\begin{tablewrap}
\begin{xltabular}{\linewidth}{>{\raggedright\arraybackslash}p{0.20\linewidth}rr}
\toprule
\thead{时点} & \thead{路径 A（先跌后涨）} & \thead{路径 B（先涨后跌）} \\
\midrule
\textbf{第 3 年末} & 105 万 & 284 万 \\
\textbf{第 10 年末} & 80 万 & 428 万 \\
\textbf{第 17 年末} & 31 万 & 708 万 \\
\textbf{第 19 年末} & 10 万 & 412 万 \\
\textbf{第 20 年末} & 0 & 352 万 \\
\bottomrule
\end{xltabular}
\end{tablewrap}

路径 B 走完二十年还剩 352 万，够她再活二十九年。路径 A 在第 19 年末只剩 10 万出头，第 20 年那笔 12 万她取不全。同样的平均收益，一个人留下三百多万，另一个人钱没了。

原因不复杂。头三年市场跌掉四成多，她还是得取 12 万，只能在低位卖出份额。卖掉的份额从此不在池子里，后面十七年的 +10\% 是给剩下的份额涨的。

像在旱季从水池里舀水。后来雨季来了，水位涨回去了，但你舀走的那几瓢不在池子里，它们没跟着涨。继续推一步：那么退休前后的三到五年，是一个人整个财务生命里风险最集中的窗口，因为它同时具备"本金最大"和"开始往外抽水"两个条件。所谓"临近退休降低权益仓位"，不是保守，是算术。

\begin{warnbox}{小心}
这一节不是教你择时。没有人知道你退休那年市场会涨还是会跌，任何声称知道的人都可以直接拉黑。

能做的是结构上的安排：让提取的那部分钱不依赖当年的市场。比如留出两到三年生活费的现金和短久期债券，跌的年份就动这一块，不动权益。这只是机制上的说明，不构成投资建议，具体怎么配得看你自己的现金流和承受力。

\end{warnbox}

\section{你只有一条时间线}

最后一件事，也是这本书里我最想让你带走的一个念头。

有个游戏。你押上全部身家抛一次硬币：正面，涨 50\%；反面，跌 40\%。

算期望收益：0.5 乘 1.5 加 0.5 乘 0.6，等于 1.05。每一轮期望赚 5\%。任何一本教科书都会告诉你这个游戏值得玩。

现在分两种玩法。

第一种：一百个人，每人玩一轮。这群人加总起来确实赚了大约 5\%，有人翻倍有人腰斩，平均是正的。

第二种：一个人，连玩一百轮。假设正反各出现五十次，他的 1 万块会变成 51 块 5。

因为 1.5 乘 0.6 等于 0.9。每凑齐一正一反，他就少掉 10\%。

\begin{statrow}{2}
  \statitem{131.5}{倍}{集合平均给出的答案（1.05 的 100 次方）}
  \statitem{0.5}{\%}{时间平均给出的答案（一个人玩 100 轮的剩余）}
\end{statrow}

这两个数字都没算错。它们回答的是两个不同的问题。前者叫\term{集合平均}{ensemble average}\index{025@集合平均}，后者叫\term{时间平均}{time average}\index{050@时间平均}。当两者不相等，这个系统就是\term{非遍历的}{non-ergodic}\index{014@非遍历的}。

\textbf{期望收益是替"很多个你"算的。而你只有一个你。}

\begin{pullquote}{}
一百个人各玩一次俄罗斯轮盘，这个游戏的期望收益是正的。一个人连玩一百次，他必死。
\end{pullquote}

所以在"把收益率从 8\% 提到 10\%"和"把爆仓概率从 3\% 降到 0"之间，永远先做后者。前者是把乘数从 1.08 改成 1.10，后者是保证这串乘数里不会出现一个 0。

任何数乘以 0 等于 0。小学二年级的内容，也是这本书里最贵的一句话。

\begin{mythbox}{常见误解}
\mvfrow{误解}{"拿得住就行，时间能抹平一切亏损。"}
\mvfrow{事实}{时间能抹平波动，抹不平永久性损失。退市的公司、违约清算的债、被强平的杠杆账户、清盘的基金，没有"后来"。就算不归零，时间也可能长到超出你的生命尺度：日经 225 在 1989 年 12 月 29 日收于 38915 点，再次站上这个数字是 2024 年，中间隔了三十四年；2026 年 6 月它第一次突破 70000 点。三十四年，是一个人从入职到退休的全程。}
\end{mythbox}

\begin{databox}{2026 数据：先把尺子摆好}
美国联邦基金利率目标区间 3.50\%–3.75\%，2026 年已连续五次维持不变（截至 7 月）。中国 7 天期逆回购政策利率 1.4\%；2026 年 8 月 20 日 LPR，1 年期 3.0\%，5 年期以上 3.5\%。

拿这几个数当基准线：在人民币资产里，任何明显高于 3.5\% 的"稳定收益"，高出来的部分都是风险的价格，不是谁的善意。这个换算你得会自己做，因为下一章里所有人都指望你不会做。

\end{databox}

\begin{takeaway}{本章一句话}
复利是乘法，不是加法。乘法怕两件事：怕中间有个特别小的数，更怕中间有个 0。

\begin{itemize}
  \item 亏 50\% 要涨 100\% 才回本；亏 90\% 按 8\% 爬要三十年。
  \item 波动本身在吃收益，压低波动等于提高收益。
  \item 复利的钱几乎全在最后十年，所以早开始比高收益重要，中途取钱代价极大。
  \item 提取期最怕开头就跌。同样的平均收益，顺序颠倒可以是"剩 352 万"和"归零"的差别。
  \item 先保证不出局，再谈收益率。你只有一条时间线。
\end{itemize}

\end{takeaway}

\begin{bridgebox}
\textbf{下一章}这一章讲的每一件事，骗子都比你更早学会了。他们知道你算不清几何平均，知道你分不清集合平均和时间平均，也知道你看到"稳定 15\%"时脑子里没有那把 3.5\% 的尺子。好消息是，所有骗局的零件加起来只有五个。
\end{bridgebox}


\input{chapters/ch-23}
\input{chapters/ch-24}
\chapter{金融的本质，是陌生人之间的信任}
\chaptermeta{25}

\begin{ledetext}
走完了二十四章管道。现在回到楼下那家砂锅店，看看那 98 块钱现在长什么样。
\end{ledetext}

\section{回到那 98 块}

序章里你扫了个码，98 块。某个机房的硬盘上，你名下那行从 6432.10 变成 6334.10，老陈名下那行从 21847.60 变成 21945.60。

当时我说，这就是"付钱"的全部物理过程。你大概只当它是个说法。

现在把同一件事再看一遍。

\begin{bsheet}
  \bsside{你的账 · 资产}{%
    \bsrow{银行存款}{6334.10}
    \bsrowtotal{这是你的资产}{6334.10}
  }
  \bsside{银行的账 · 负债}{%
    \bsrow{你的存款}{6334.10}
    \bsrowtotal{这是它欠你的}{6334.10}
  }
\end{bsheet}

你手机里那个数字，从来不是"钱"。它是一家银行写给你的欠条。你把这张欠条转了 98 的份额给老陈，两家银行之间怎么算清，去了央行的账本上。你吃一顿饭，惊动了中央银行，这句话现在你知道它字面上是真的。

老陈拿这笔流水去还了一期经营贷。还掉的那部分本金，从这个世界上消失了：银行资产端的贷款少一块，负债端的存款同步少一块，货币总量真的变少。他剩下的钱买了一只 R2 理财，理财买了信用债，债的钱到了江苏一家做储能电芯的厂子，厂子订了一台德国的极片分切机，人民币换成欧元，欧元进了德国厂商的财资账户，去买了短期债券。

整条链现在应该都亮着。

\begin{chainflow}
\chainnode{钱是记账}\chainarrow \chainnode{记账要有人信}\chainarrow \chainnode{信任靠制度和抵押品}\chainarrow \chainnode{银行放贷时创造存款}\chainarrow \chainnode{钱约等于债}\chainarrow \chainnode{债要付息}\chainarrow \chainnode{付息需要增长}\chainarrow \chainnode{增长需要投资}\chainarrow \chainnode{投资需要给风险定价}\chainarrow \chainnode{定价错误加杠杆等于危机}\chainarrow \chainnode{危机后重置资产负债表}\chainarrow \chainnode{循环重来}
\end{chainflow}

二十四章之前，这行字对你是十二个词组。现在它是一条你能倒着走回来的路。

\section{它只干三件事}

把所有的产品、机构、术语、监管条文全部拿掉，剩下的只有三件事。

\textbf{跨时间搬钱。}把你四十岁的收入搬到七十岁去花，把你二十八岁买不起的房子搬到今天来住。存款、贷款、债券、养老金，都是这一件事的不同穿法。

\begin{egbox}{举个栗子：跨时间搬钱的标价}
老陈那 60 万经营贷，三年期，年利率 5.4\%，等额本息。月供 18091 元，三十六个月一共还回去 65 万 1 千多。

他今天多了一家店，代价是未来三年每个月固定支出 18091 元，以及 5 万 1 千多的利息。

这就是"把明天的钱搬到今天"的标价。所有金融产品的定价，剥到最后都是在给这一件事标价：\textbf{时间是有价格的。}第 3 章那句话，走到这里正好绕回来。

\end{egbox}

\textbf{在人群中分摊风险。}一个人扛不起的事，一万个人一起扛就不算事。保险、股份公司、再保险、衍生品，本质都是把一个人的悬崖，摊成一万个人的台阶。

\textbf{让陌生人之间的承诺变得可执行。}

第三件最难，也是这本书书名的那个意思。前两件事在任何有血缘和村庄的社会里都存在，只是规模小得可怜。真正把规模放大到几十亿人的，是第三件。

\section{最不性感的那些东西}

问一个具体的问题：老陈想再开一家分店，缺 60 万。

\begin{egbox}{举个栗子：银行凭什么借给一个陌生人 60 万}
信贷员根本不认识老陈。他看的是三样东西。

一，三年的收单流水和纳税记录。这依赖会计准则和支付清算系统，让"这家店到底赚不赚钱"变成一个可以被外人核验的数字。

二，有没有能抵押的东西。老陈名下那套小两居能不能抵押，取决于不动产登记系统能不能在一天之内把抵押权登记上，并且这个登记在法庭上算数。

三，还不上怎么办。这依赖破产法和强制执行：清偿顺序是可预期的，抵押物是能被处置的。

三样里缺一样，老陈的利率就从 5\% 跳到 18\%；缺两样，他只能去找亲戚。这就是第 13 章那句"风险的价格"落地之后的样子。

\end{egbox}

会计准则、清算系统、破产法、抵押登记、审计、监管。没有一样能上头条，没有一样有故事性，绝大多数人一辈子不会主动读它们中的任何一行。

\textbf{但信任的物理基础，全在这里。}它们让"我凭什么信一个我永远不会见面的人"这个问题，有了一个不依赖人品的答案。

\begin{pullquote}{}
一个国家的利率下限，很大程度上是由它的法院和登记处决定的，不是由它的央行决定的。
\end{pullquote}

所以我愿意下一个不太讨喜的判断：\textbf{一个社会金融体系的上限，不是由它最聪明的交易员决定的，是由它最不被注意的那些规则决定的。}金融创新的速度可以很快，信任的地基长得很慢。地基长不上去的地方，创新出来的东西最后都会变成第 23 章里的某一行。

\begin{mythbox}{常见误解}
\mvfrow{误解}{"金融不创造任何价值，就是一群人在互相倒手，本质是零和赌场。"}
\mvfrow{事实}{交易环节确实有大量零和的部分，这一点批评得对。但把资金从"手里有钱但没项目的人"送到"有项目但没钱的人"手里，这件事本身创造价值：老陈的第二家店雇了 7 个人，那台德国分切机让电芯厂的良品率上去了。问题不在于金融该不该存在，而在于它有多大比例真的在干这件事，多大比例只是在自己内部循环收费。这个比例，全世界都还没算清楚。}
\end{mythbox}

\bookbreak

\section{三个从来没变过的问题}

雅浦岛沉在海里的石头、罗马的金币、宋朝的交子、汇丰的银行券、你的支付宝余额、链上的稳定币。

技术换了六次，问题一次没换。

\textbf{谁在承诺？}这张欠条是谁写的。是主权国家、是持牌银行、是一家开曼注册的公司，还是一段没人负责的代码。

\textbf{凭什么信他？}他背后有什么：税收权、储备资产、抵押品、法律约束，还是仅仅是一份 PDF。

\textbf{出事了谁兜底？}存款保险、央行的最后贷款人、破产清算顺序，还是"祝你好运"。

拿这三句去问 2026 年任何一个新东西，比如稳定币：发行方是一家私人公司；它凭储备资产（短期国债）和 GENIUS Act 建立的监管框架；出事了看储备够不够、赎回条款怎么写，没有存款保险。三句话问完，你对它的位置就有数了。它不是骗局，也不是新黄金，它是一张利息归别人的、跑得很快的欠条。

\section{九条}

不凑十条。这九条是从前面二十四章里长出来的，不是我编的口号。

\begin{flowlist}
  \flowitem{01}{任何"钱"都是某个人的负债}{你的存款是银行欠你的，现金是央行欠你的，稳定币是发行方欠你的。搞清楚是谁欠你，比搞清楚有多少更重要。}
  \flowitem{02}{一个人的资产是另一个人的负债}{所以"全社会同时去杠杆"在算术上做不到，只能转移。谁接走了，是每一次危机真正的争论点。}
  \flowitem{03}{高收益不是被发现的，是被承担的}{没有藏在角落里等你捡的高收益。多出来的每一个点，都对应着某种你必须扛下来的东西。}
  \flowitem{04}{能不能随时拿回来，比收益率重要}{流动性是唯一一个不能被话术绕过的检验。}
  \flowitem{05}{不要用会被追缴的钱做期限长的事}{短钱长投是所有爆仓的同一个死法，从 1998 年的 LTCM 到你隔壁那个用信用卡套现炒股的人，机制完全一样。}
  \flowitem{06}{先活下来，再谈收益率}{任何数乘以 0 都是 0。你只有一条时间线。}
  \flowitem{07}{储蓄率是你唯一能完全控制的变量}{收益率不听你的，市场不听你的，房价不听你的。每月存多少听你的。}
  \flowitem{08}{波动不等于风险}{永久性损失才是风险。但如果波动会逼你在低点卖出，它就变成了风险，这也是为什么现金缓冲比看上去重要。}
  \flowitem{09}{你看到的收益率，都是幸存者的收益率}{包括所有的历史数据、所有的排行榜、所有的成功故事。包括这本书引用的那些。}
\end{flowlist}

\section{我们不知道的事}

一本在这些问题上装作有答案的书，两年就过期了。所以老实列在这里。

AI 究竟会不会提高全要素生产率，还是只是把成本从一个部门挪到另一个部门。这轮资本开支已经花掉了，产出还没到。

2\% 是不是最优的通胀目标。这个数字最早来自 1980 年代末新西兰的一次政策讨论，从来没有严格的理论证明。

货币政策的传导到底有多强。2021 到 2023 年的那轮通胀里，加息和通胀回落之间有多少是因果，学界还在吵。

加密资产最终会变成什么。支付基础设施、投机资产、监管套利工具，还是三样都有一点。

被动投资占比的临界点在哪里，或者到底有没有临界点。

这五个问题，2026 年 8 月没有人知道答案。任何告诉你他知道的人，你现在应该有能力判断他为什么要这么说。

\begin{takeaway}{本章一句话}
金融是一整套让陌生人之间的承诺变得可执行的技术。技术一直在换，那三个问题一次没换：谁在承诺，凭什么信他，出事了谁兜底。

\begin{itemize}
  \item 它只干三件事：跨时间搬钱、在人群中分摊风险、让陌生人的承诺可执行。
  \item 信任的地基是会计、清算、破产法和登记，不是聪明的交易员。
  \item 不做金融的信徒，也不做它的敌人。做一个看得懂它的人。
\end{itemize}

\end{takeaway}

\begin{databox}{最后一次提醒时间口径}
本书写于 2026 年 8 月。书里所有的数字都带着这个日期：美国联邦基金利率 3.50\%–3.75\%，中国 5 年期以上 LPR 3.5\%，全球央行储备里黄金占 27\%，全球债务约 353 万亿美元。

这些数字明年就不对了。管道明年还在。

\end{databox}

上个月老陈把价格调了。两个菜一碗汤，现在 104。

他说是猪骨和菜价都涨了。他不知道的是，这 6 块钱里有一小块是中东的油价，还有一小块是他去年那笔经营贷的利息。

你知道。

汤还是好喝的。

\begin{bridgebox}
\textbf{下一站}附录：术语速查表，以及 2026 年 8 月的数据快照。前者用来查，后者用来两年后回头对照，看看哪些变了、哪些没变。
\end{bridgebox}



% 后记不编号，但进目录
\part{后记：合上书之后}
\begingroup
\let\ELEGANTchapter\chapter
\renewcommand{\chapter}[1]{%
  \ELEGANTchapter*{#1}%
  \addcontentsline{toc}{chapter}{#1}%
  \markboth{#1}{#1}}
\chapter{把二十六章压成一页}
\chaptermeta{I}

\begin{ledetext}
合上书两年之后，细节会忘。这一页让你在五分钟内，把二十六章重新拼回一条链。
\end{ledetext}

这一页是这本书的替身。

你合上书两年，细节会掉，最后只剩几个模糊的念头；到那天再翻回二十六章，成本太高，也太慢。所以这里把全书压成一张五分钟能扫完的图：一条链，四卷，二十六句，最后五件事。忘得只剩这一页，也够你把整座建筑重新立起来。

用法只有一种：忘了，就翻到这一页扫一遍；扫完还拿不准，顺着那句话翻回原文。

\section{一条链}

全书只讲这一条链，反复讲，换着角度讲。它一共十二环。

\begin{chainflow}
\chainnode{钱是记账}\chainarrow \chainnode{记账要有人信}\chainarrow \chainnode{信任靠制度和抵押品}\chainarrow \chainnode{银行放贷时创造存款}\chainarrow \chainnode{钱约等于债}\chainarrow \chainnode{债要付息}\chainarrow \chainnode{付息需要增长}\chainarrow \chainnode{增长需要投资}\chainarrow \chainnode{投资需要给风险定价}\chainarrow \chainnode{定价错误加杠杆等于危机}\chainarrow \chainnode{危机后重置资产负债表}\chainarrow \chainnode{循环重来}
\end{chainflow}

十二环连起来读，是一口气：钱是记账，记账要有人信，信任靠制度和抵押品垫着；所以银行放贷那一刻，存款是当场造出来的，钱约等于债。债要付息，付息就得有增长，增长靠投资，投资得先给风险定价。定价一旦错了，再叠上杠杆，就是危机；危机逼人重置资产负债表，重置完，下一轮从头再来。

这条链上还有一个位置，常被看漏：你没有站在外面看它，你站在末端，是一张很小的资产负债表。机器怎么转你左右不了，但那张小表怎么摆，是你唯一说了算的事。这也是四卷为什么一路从天上讲到地上。

这条链还能倒着读。一场危机过后，钱没有消失，只是换了名字记在别人账上；从“重置资产负债表”往回摸，能摸到谁在借、谁在还、谁在扛损失。

随便拿一条新闻试。看到“某银行债券利差走阔”，往链条上找位置：这是“投资需要给风险定价”那一环，市场在重估谁会还不上钱。看到“居民提前还贷”，那是“钱约等于债”这一环在收缩。每一条新闻都能挂到其中一环上。

\section{四卷各管哪一段}

四卷不是按主题分的，是按这条链的进度分的。前一卷故意留的坑，后一卷负责填。

地基篇管头三环：钱是记账，记账要有人信，信任靠制度和抵押品。它给的不是知识，是一套读账的语法：资产、负债、净值，利息、通胀、贴现。解决的是“黑话太多”这件事，让没学过经济的人先能听懂饭桌上的对话。它回答“钱凭什么值钱”，不回答“怎么变有钱”；还故意留了一个坑：银行放出去的钱，最初从哪儿冒出来。

管道篇接过那个问题，管第四到第八环：放贷创造存款，钱约等于债，债要付息，付息需要增长，增长需要投资。它把货币怎么被造出来、央行管什么、利率怎么传到你的月供、债股汇怎么标价，一节一节铺开，让新闻标题背后的管道变得能看穿。走到结尾，你已经能把一条新闻翻译成一段因果。它留下一个更大的问号：管道里流的东西，为什么每隔几年就炸一次。

天气篇管第九到第十一环：给风险定价、加杠杆、出危机、重置资产负债表。它讲清楚三件事：风险的价格怎么算，杠杆怎么把小麻烦放大成灾难，泡沫和周期为什么必然反复。解决的是“为什么好日子和坏日子轮流来”。读完你会明白，危机不是意外，是这台机器的固定零件。它不管一件事：这些潮汐落到你个人头上，该怎么办。

落地篇管最后半环，也就是你：怎么做自己的资产负债表，收益到底从哪来，怎么配置才不白担风险，复利哪句是真哪句是假，骗局靠什么零件运转，2026 年哪六个变量在改写算式。从这里往后，书不再教你去看懂别人，只教你摆平自己那张表。它不再往下留问题，因为它绕回了起点：钱是记账，记账要人信，而你此刻已经知道谁在欠你。

四卷可以跳着读。急着防骗，直接翻落地篇；想懂新闻，先补管道篇。但那条链，最好从头到尾各走一遍。

\section{二十六章，压成二十六句}

下面每句都是那一章真正的落点，不是标题的复述。按顺序读，是全书骨架；随手挑一句，那一章就该在眼前立起来。这二十六句也不是并列的，是把上面那条链拆成二十六段；哪一句看着眼生，就说明那一环你已经忘了。

\begin{flowlist}
  \flowitem{00}{一张一百块的环球旅行}{你花掉的 98 块没动过一毫米，只是两行数字一减一加，却已顺手参与了一次清算、一次汇率报价。}
  \flowitem{01}{钱不是财富，钱是一套记账系统}{先有欠账，后有铸币；钱不是藏在金库里的东西，是一群人一起记住的账。}
  \flowitem{02}{你的存款，其实不在银行的金库里}{你把钱存进银行的那一刻，它就不再属于你，剩下的只是银行欠你的一张条。}
  \flowitem{03}{利息：时间是有价格的}{利息拆开看，是时间、通胀、违约、流动性四样东西的价，反过来就是贴现率。}
  \flowitem{04}{通货膨胀：每天从你口袋里拿走一点的那个人}{通胀不改动任何一行数字，它只把量东西的那把尺子悄悄改短，再在债主和欠债人之间搬钱。}
  \flowitem{05}{资产、负债、净值：全世界通用的三个词}{资产等于负债加净值，净值才是财富；而一个人的金融资产，必然是另一个人的负债。}
  \flowitem{06}{所有金融的原点：今天的钱换明天的钱}{金融只干两件事：把今天的钱搬到明天，把风险从一个人摊给一群人，其余都是工程细节。}
  \flowitem{07}{货币是怎么“凭空”出现的}{银行放贷的那一刻钱才出生，你还款的那一刻这笔钱就死了，从没有人问过一句“这能兑现吗”。}
  \flowitem{08}{央行不是印钞厂，是水闸和温度计}{央行决定不了钱有多少，它定的是钱的价格；它能把马牵到河边，逼不了马喝水。}
  \flowitem{09}{利率的传导链：从政策利率到你的房贷}{政策利率到你的月供之间隔着七道闸，每一道都可能把信号拖慢，甚至让它蒸发在半路。}
  \flowitem{10}{债券：世界上最大、最无聊、最要命的市场}{债券的价格和收益率永远反向，期限越长跌得越狠，因为那是一个除法，不是行情。}
  \flowitem{11}{股票：把一家公司切成一亿片}{股票是排在所有人后面的那份剩饭，风险最大，也正因如此，长期回报最高。}
  \flowitem{12}{汇率与美元体系：为什么全世界用别人的钱做生意}{汇率是两种钱互相的报价，今天主要看资本往哪儿跑，货往哪儿走是次要的。}
  \flowitem{13}{风险的价格：为什么有人借钱 3\%，有人 12\%}{利率是一叠拆得开的账单，看不清的那部分按最坏的算，抵押品和征信能把它压下来。}
  \flowitem{14}{杠杆：金融世界的放大镜与绞肉机}{杠杆不改变你的对错，只改变你能撑多久；放大收益和让你归零的是同一个数。}
  \flowitem{15}{衍生品：从农民的保险单到华尔街的赌桌}{衍生品是切下来转手的风险，它是不是赌，取决于你身上原本有没有那份方向相反的风险。}
  \flowitem{16}{影子银行：管道之外的管道}{影子银行做的是同一门生意，拿短钱做长事，只是出事时没有存款保险，也没有最后贷款人。}
  \flowitem{17}{泡沫与崩溃的通用剧本}{泡沫要四个零件：真实的新故事、便宜的信贷、杠杆、新进场的人；故事通常是真的，错的是价格和时间表。}
  \flowitem{18}{周期：为什么好日子和坏日子会轮流来}{周期不是钟表，是四五个不同长度的钟摆挂在一起；你能判断站在哪一段，算不出下一个拐点。}
  \flowitem{19}{给自己做一张资产负债表}{净值是算出来的，不是感觉出来的；你表上最大的一笔资产是你自己，而它没上保险。}
  \flowitem{20}{收益从哪儿来：三种钱和一个残酷等式}{你赚的每一块钱只有三个来路：经济增长、别人的口袋、被摩擦吃掉，后两种早晚找上你。}
  \flowitem{21}{资产配置：唯一的免费午餐}{把不同步的东西放在一起，收益不减、波动变小，它的开关是相关性，不是数量。}
  \flowitem{22}{复利的真相，和关于复利的谎言}{复利是乘法，怕中间有个特别小的数，更怕中间有个 0；所以先别出局，再谈收益。}
  \flowitem{23}{骗局的五个零件}{骗局只有五个零件：看不懂的故事、又高又稳的承诺、紧迫、熟人、退出摩擦，认全了外壳随便换。}
  \flowitem{24}{2026：六个你必须知道的新变量}{六件事同时在变天：AI 烧钱、支付重做、黄金上位、私募膨胀、利率上移、被动资金变巨兽。}
  \flowitem{25}{金融的本质，是陌生人之间的信任}{金融是让陌生人敢签字的机器，技术换了六次，三个问题没换：谁承诺、凭什么信、出事谁兜底。}
\end{flowlist}

书里每个数字都标了 2026 年 8 月。两年后回来看，数字已经不对，但这一页没有一句依赖某个数字。数字当水位读，机制当河道读，河道不搬家。

\begin{pullquote}{}
大多数人研究钱的顺序是反的：先问买什么，最后才问钱从哪儿来。而骗局和危机，吃的都是这个顺序。
\end{pullquote}

\section{如果只记住五件事}

这五条不是排在最前面的五章，是全书里最值钱的五个结论。它们的共同点：不需要任何背景，能直接压在一个具体决定上。

\begin{takeaway}{如果只记住五件事}
从上面二十六句里再蒸馏出五条。忘掉其余全部，也不会被骗，也能看懂新闻。

\begin{itemize}
  \item 任何“钱”都是某个人的负债。先问谁欠你、拿什么撑着，再谈别的。
  \item 钱是被银行借出来的，不是印出来的。大家愿意借，钱才变多；抢着还钱，钱就变少。
  \item 利率是全世界的除数。它一动，股票、房子、债券、你的月供一起重算。
  \item 高收益是风险的价格，不是善意。“又高又稳”四个字凑齐，基本是骗局。
  \item 你只有一条时间线。先保证不出局，再谈收益；储蓄率是唯一全归你管的变量。
\end{itemize}

\end{takeaway}

这五条不用记原文，记意思就够。真拿不准一笔钱，回到第一条：谁欠你，拿什么撑。

回到序章那个扫码的晚上。你扫了 98 块，老陈账上多 98，你账上少 98，什么都没变，又什么都变了：这一回，你看见了那 98 块接下来要去哪儿。

书到这里就完了。它没能让你多赚一分钱，也没打算。它只在往后的每一次扫码、每一次看新闻、每一次签合同里，多给你留了一秒钟的明白。

有这一秒，就够了。

\begin{bridgebox}
\textbf{下一篇}全书讲了二十六章钱怎么跑，却始终没教你发财。这不是漏了，是故意的。后记二《我们为什么从头到尾不谈发财》，说这件事。
\end{bridgebox}


\input{chapters/ep-2}
\chapter{给读者的信}
\chaptermeta{III}

\begin{ledetext}
最后这几页不讲知识。讲点别的。
\end{ledetext}

\section{我不知道你是谁}

可能你 23 岁，刚签完第一份劳动合同，HR 让你填公积金比例，你盯着那个空格不知道填多少，也不好意思问。

可能你 34 岁，房贷还了五年，最近每天晚上算提前还贷划不划算，算到十二点还没算明白，第二天又接着算。

可能你 41 岁，公司上个月做了一轮优化，你没被优化，但你第一次很认真地想了一个问题：如果下一轮是我呢。

可能你 58 岁，退休金刚开始发，银行客户经理这半年对你格外热情，每次去都要拉你坐下聊十分钟。

也可能你只是刷到了这本书，本来打算看两段就关掉，结果不知怎么看到了第 7 章，发现自己活了这么多年，居然一直搞错了钱是从哪儿来的。

这五个人的处境差得很远。但他们面对的是同一台机器。

这本书是写给这台机器里的人的。

\section{三件我希望你带走的东西}

\textbf{第一件，一副眼镜。}

资产 = 负债 + 净值。一个人的资产，是另一个人的负债。

这两句话看起来像会计课的开头，实际上是这本书里最锋利的工具。学会了它，你看新闻的方式会永久地变掉：听到"全球债务 353 万亿美元"，你会下意识地问一句，那这 353 万亿是谁的资产？听到"银行放水"，你会问，钱进了谁的账户，又是谁的负债增加了？

问出这句话的人，和没问出这句话的人，看到的是两个世界。

\textbf{第二件，三个问题。}

\begin{pullquote}{}
谁在承诺？凭什么信他？出事了谁兜底？
\end{pullquote}

以后每出现一个新东西，你不需要先弄懂它的技术细节。先问这三句。三句问完，它在你心里的位置就基本定了。

这三句话大概是这本书里保质期最长的部分。等到 2035 年，书里那些利率、汇率、储备占比全都不对了，这三句还对。

\textbf{第三件，一种和不确定性相处的舒服感。}

这本书里至少有五处写着"目前没人知道答案"。AI 会不会真的提高生产率，2\% 是不是最优的通胀目标，被动投资的临界点在哪儿，加密资产最终变成什么。

刚开始读的时候，这种话可能让你不舒服。你是来找答案的。

但一个能坦然说出"这件事没人知道"的人，比一个什么都有答案的人，在真实世界里活得稳得多。因为骗你的人从不说这句话，他们的答案总是很完整。

\section{一件我希望你别带走的}

一种新的傲慢。

懂了一点之后最容易犯的毛病，是觉得那些还不懂的人很蠢，是觉得自己从此不会栽跟头了。

提醒一句：2008 年之前，全世界最懂金融的那批人，就聚在一起，造出了那场危机。他们不是不懂，他们比谁都懂。知识没能拦住他们，因为知识拦不住激励。每个人都在按自己的奖金规则做最优选择，加起来就是灾难。

所以懂了之后真正该增加的不是自信，是两样别的东西：对自己判断的怀疑，和对"我这样做是因为谁在付我钱"这个问题的敏感。

\section{如果你现在手头很紧}

这本书默认了一件事：你有一点可以支配的钱。

我知道不是所有人都有。

如果你现在的状态是每个月都不太够，那对你最有用的不是第 21 章的资产配置，那一章你可以先跳过。对你最有用的是第 19 章里"储蓄率是你唯一能完全控制的变量"这半页，和整个第 23 章。

因为钱少的时候，被拿走一次的代价，比钱多的时候大得多。三万块对一个有三百万的人是 1\%，对一个有五万块的人是半条命。

骗子最偏爱的从来不是有钱人。有钱人身边有一堆人替他把关。

这一段是这本书里我最希望被读到的一段。

\section{关于"引路人"这三个字}

这本书不能替你走路。

它能做的事情小得多：在你已经走了很多年、以后还要接着走的那条路上，装了几盏灯。

灯照得不远。它照不到十年后的利率，照不到你们行业还剩几年好光景，照不到你父母的身体和你孩子的运气。再往前的路仍然是黑的，所有人的都黑，包括写这本书的人和管着几万亿的人。

灯只做一件事：让你看清脚下这几米有没有坑。

而人生里那些真正把人绊倒的坑，大部分就在脚下这几米。不在远处。

\section{接下来这一周，三件小事}

不是作业。是这本书唯一一次请你动手。

\begin{flowlist}
  \flowitem{01}{花二十分钟，把自己的表填出来}{打开你的银行 App、公积金账户、基金账户，按第 19 章的样子列一张：左边资产，右边负债，中间那个差额就是你的净值。不用给任何人看，写在手机备忘录里就行。大多数人活一辈子没做过这件事，而它只需要二十分钟。}
  \flowitem{02}{翻出你的房贷合同或者租房合同}{找到利率那一页，看清楚三件事：加了多少个基点、重定价日是哪一天、提前还款有没有违约金。如果你还没有贷款，那就看看你手里那笔存款或者理财的实际年化，再拿它跟当月的通胀比一比。}
  \flowitem{03}{挑一句最反直觉的，讲给一个人听}{比如"银行放贷的时候，钱才被创造出来"。讲的过程中你会卡壳，卡壳的地方就是你还没真懂的地方。回去把那一章再看一遍。这是检验理解最便宜的办法。}
\end{flowlist}

\section{最后}

写这本书的过程里，我一直在想一个问题：一个普通人，花十来个小时读完十万字，到底能换到什么。

答案可能没有想象中那么壮观。

你不会因此变富。你的房子不会涨回来。你也拦不住下一次危机，那件事谁都拦不住。

你换到的只是：在往后每一次刷到财经新闻、每一次坐在客户经理对面、每一次在合同上签字的时候，脑子里多出来的那一两秒钟。

那一两秒里，你会想起来问一句"这钱从哪儿来"。

就是这个。

\bookbreak

某个周五晚上，你大概还是会在楼下点两个菜一碗汤。手机叮一声。

那一声，和你读这本书之前的每一声，听起来一模一样。

只有你知道，它已经不一样了。

路还长。慢慢走。

\begin{takeaway}{带走这三样就够}
一副眼镜：资产等于负债加净值，一个人的资产是另一个人的负债。三个问题：谁在承诺，凭什么信他，出事了谁兜底。一种舒服：能坦然说出"这件事没人知道"。

\begin{itemize}
  \item 别带走傲慢。2008 年之前最懂金融的那批人，正是造出那场危机的人。
  \item 钱越少，被拿走一次的代价越大。手头紧的时候，第 23 章比第 21 章重要。
  \item 这本书装的是脚下几米的灯，不是远方的地图。而绊倒人的坑大多在脚下几米。
\end{itemize}

\end{takeaway}

\begin{bridgebox}
\textbf{接下来}附录里是术语速查和 2026 年 8 月的数据快照。前者用来查，后者留着两年后回头对照，看看哪些变了，哪些其实一直没变。
\end{bridgebox}


\endgroup

% ================= 后置部分 =================
\backmatter
\appendix
\part{附录}
\chapter{术语表}
\markboth{术语表}{术语表}

\begin{ledetext}
本表收录正文中出现的核心术语，按汉语拼音排序。每条后面标注它在哪一章展开。看到不认识的词回来查一眼，然后跳回正文，看它是怎么长出来的。
\end{ledetext}

\begin{multicols}{2}

\glossletter{A}
\glossitem{72 法则}{}{72 除以年化收益率，约等于本金翻倍所需年数。按 3.5\% 算，需要约 20.6 年。}{03}

\glossletter{B}
\glossitem{保证金追缴}{margin call}{价格下跌触发追加保证金，拿不出就被强平。判断正确但活不到被证明的那天，就是这么来的。}{14}
\glossitem{被动投资}{}{按指数权重机械买入。2025–2026 年全球规模首次超过主动管理，价格由谁来发现成了新问题。}{11}
\glossitem{遍历性}{ergodicity}{集合平均不等于时间平均。一百个人各玩一次和一个人连玩一百次，是两件完全不同的事。}{22}
\glossitem{波动损耗}{volatility decay}{涨 10\% 再跌 9.09\% 回到原点，三倍杠杆产品却亏了 5.5\%。杠杆 ETF 只适合极短期。}{14}

\glossletter{C}
\glossitem{储蓄率}{}{收入减支出的比例。这是全书里唯一一个你能百分之百控制的变量。}{19}
\glossitem{存款保险}{}{中国 50 万元以内本息全额偿付。它真正的作用不是事后赔钱，是事前掐断挤兑的念头。}{02}

\glossletter{D}
\glossitem{贷款创造存款}{}{银行放贷时同时凭空写下一笔资产（贷款）和一笔负债（存款），新货币由此诞生。全书最反常识、也最重要的一条。}{07}
\glossitem{到期收益率}{YTM}{按当前价格买入并持有到期的年化回报。价格与它永远反向变动。}{10}

\glossletter{F}
\glossitem{法定货币}{fiat money}{国家指定、以税收和法偿性为支撑的货币。它的价值不来自黄金储备，来自你必须用它交税。}{01}

\glossletter{H}
\glossitem{货币乘数}{}{过时的教科书模型：认为央行投放准备金后银行按倍数放大。现代银行体系里因果是反过来的。}{07}

\glossletter{J}
\glossitem{挤兑}{bank run}{一个自我实现的协调博弈：我不去取钱的前提是别人也不去。硅谷银行证明了它在社交媒体时代能有多快。}{02}
\glossitem{净息差}{NIM}{银行放贷收益与负债成本之差。中国商业银行 2026 年二季度末为 1.41\%，处于历史低位。}{02}
\glossitem{净值}{net worth}{资产减负债。房子按市价记，房贷按余额记，中间的差额才是你的。收入不是财富，净值才是。}{19}
\glossitem{久期}{duration}{债券对利率的敏感度。久期 8 年，利率升 1\%，价格大约跌 8\%。硅谷银行死在这两个字上。}{10}

\glossletter{L}
\glossitem{量化宽松 / 紧缩}{QE / QT}{央行买入或卖出债券以影响长端利率。QE 灌的是银行间的准备金池，不会自动流进你的口袋。}{07}

\glossletter{M}
\glossitem{名义利率 / 实际利率}{}{实际利率 ≈ 名义利率 − 通胀。存款 2\% 而通胀 3\%，你的钱在变多，购买力在变少。}{03}
\glossitem{明斯基时刻}{}{稳定本身孕育不稳定。日子越安稳，人越敢借钱，系统越脆弱。}{17}

\glossletter{Q}
\glossitem{期权}{option}{买方付权利金换取权利：损失有限，收益可观。卖方相反，像在压路机前捡硬币。}{15}
\glossitem{期限转换}{maturity transformation}{银行的核心魔法：把随时可取的存款变成三十年才还完的房贷。它创造价值，也制造脆弱。}{02}
\glossitem{前瞻指引}{forward guidance}{央行用语言影响预期。2026 年美联储主席沃什明确表示要淡化它，这是沟通范式的转向。}{08}
\glossitem{权益乘数}{}{总资产 ÷ 净资产。银行通常十几倍，所以 0.8\% 的资产收益率能做出两位数的股东回报。}{14}

\glossletter{R}
\glossitem{人力资本}{human capital}{你未来收入的贴现值。对大多数人来说这是最大的一笔资产，而它不在任何一张对账单上。}{19}

\glossletter{S}
\glossitem{三元悖论}{}{固定汇率、资本自由流动、独立货币政策，三者最多得其二。香港、欧元区、中国各自选了不同的两个。}{12}
\glossitem{剩余索取权}{residual claim}{股东的法律身份：付完供应商、工资、利息、税，剩下的才归你。排在最后，所以风险最高、长期回报也最高。}{11}
\glossitem{市盈率}{PE}{股价除以每股收益。它不是独立真理，只是贴现公式的一个简写。}{11}
\glossitem{收益率曲线}{}{不同期限国债收益率连成的线。倒挂曾是可靠的衰退先行指标，2022–2023 年它失灵了，争议至今未平。}{09}
\glossitem{私募信贷}{private credit}{2025 年全球规模约 2.3 万亿美元，利率 8\%–12\%。中金称之为“2 万亿美元的灰犀牛”。}{16}

\glossletter{T}
\glossitem{套息交易}{carry trade}{借低息货币买高息资产。利差是慢慢赚的，汇率是一夜之间还回去的。}{12}
\glossitem{特里芬难题}{}{世界要用美元做储备，美国就得持续输出美元；而长期逆差又会侵蚀对美元的信心。}{12}

\glossletter{X}
\glossitem{系统性风险}{}{分散不掉的那部分。2008 年的教训是：危机中所有相关性都趋近于 1。}{13}
\glossitem{现值与贴现}{present value}{未来的钱换算成今天值多少。贴现率一变，全世界所有资产的价格都要重算。}{03}
\glossitem{信用利差}{credit spread}{同期限下信用债与国债的收益率之差，用来补偿违约风险。它比股市更早知道事情不对。}{13}
\glossitem{信用违约互换}{CDS}{给别人的债买的保险，而且你可以给不属于你的房子买火险。2008 年 AIG 就栽在这上面。}{15}
\glossitem{序列风险}{sequence risk}{同样的平均收益，顺序不同结果天差地别。在退休提取阶段，头几年的大跌是致命的。}{22}

\glossletter{Y}
\glossitem{一级 / 二级市场}{}{公司发新股拿到钱的地方 / 投资者之间转手的地方。你在二级市场买股票，公司一分钱也拿不到。}{11}
\glossitem{影子银行}{}{干着银行的三件事，却不在银行监管框架内，也没有存款保险。它不违法，它是金融体系的正常器官。}{16}
\glossitem{预期损失}{}{违约概率 × 违约损失率 × 敞口。小微企业融资贵不是银行黑心，先是数学，然后才是摩擦。}{13}

\glossletter{Z}
\glossitem{再平衡}{}{按纪律把偏离目标的仓位调回来。它强迫你在下跌时买入、上涨时卖出，而这两件事靠意志力做不到。}{21}
\glossitem{政策利率}{}{央行直接设定的那个价格。中国以 7 天期逆回购利率（1.4\%）为主，美国是联邦基金利率目标区间（3.50\%–3.75\%）。}{08}
\glossitem{滞胀}{}{增长停滞与通胀并存。2026 年的组合是供给冲击加高利率加高债务，这是它危险的原因。}{18}
\glossitem{准备金}{reserves}{商业银行存在央行的钱，只在银行之间清算时使用。QE 创造的是它，不是你账户里的存款。}{07}
\glossitem{资产负债表衰退}{}{企业和居民的目标从赚钱变成还债，利率再低也不借。日本用三十年演示了一遍。}{09}
\glossitem{最后贷款人}{}{白芝浩原则：危机时对有偿付能力的机构、以惩罚性利率、凭优质抵押品无限放款。}{08}
\glossitem{LPR}{}{贷款市场报价利率。18 家报价行在政策利率基础上加点报出，你的房贷 = LPR ± 合同里固定的加点。}{09}
\glossitem{M0 / M1 / M2}{}{流通中的现金 / 加上企业活期存款 / 再加上定期与储蓄存款。中国 M2 三百多万亿元，M0 只有十几万亿——95\% 以上的钱从没当过纸。}{01}
\end{multicols}
\cleardoublepage

\chapter{2026 年 8 月数据快照}
\markboth{数据快照}{数据快照}

\begin{ledetext}
把这些数字当刻度读，不要当结论读。三个月之后它们就不准了，但它们所处的结构会存在很久。
两年后回头翻这一页，比较哪些变了、哪些其实一直没变，是这份附录真正的用处。
\end{ledetext}

\section{利率：全球央行第一次不再同步}

\begin{tablewrap}
\begin{xltabular}{\linewidth}{>{\raggedright\arraybackslash}p{0.24\linewidth}Xr}
\toprule
\bthead{央行 / 品种} & \bthead{水平与动向} & \bthead{时点} \\
\midrule
\textbf{美国联邦基金利率} & 目标区间 3.50\%--3.75\%，2026 年已连续五次维持不变。主席凯文·沃什称通胀仍然高企，市场一度讨论加息可能 & 2026.07 \\
\textbf{欧洲央行} & 上调存款便利利率，为三年来首次加息 & 2026.06 \\
\textbf{日本银行} & 以 7 比 1 加息 25 个基点至 1.0\%，为 1995 年 9 月以来最高；国债购买规模继续压降 & 2026.06 \\
\textbf{中国 1 年期 LPR} & 3.0\%，连续多月持平 & 2026.08.20 \\
\textbf{中国 5 年期以上 LPR} & 3.5\%，连续多月持平 & 2026.08.20 \\
\textbf{中国 7 天期逆回购} & 1.4\%，主要政策利率 & 2026.08 \\
\textbf{中国商业银行净息差} & 1.41\%，历史低位，构成降息空间的现实约束 & 2026.Q2 \\
\textbf{美国 10 年期国债} & 收益率升至 2025 年 1 月以来最高，长端与短端明显分化 & 2026.08 \\
\bottomrule
\end{xltabular}
\end{tablewrap}

日本央行这一轮正常化的完整路径值得单独记一笔：2024 年 3 月结束负利率并终结收益率曲线控制，
政策利率由 $-0.1\%$ 上调至 $0\sim0.1\%$；2024 年 7 月加至 0.25\%；2025 年 1 月加至 0.5\%；
2025 年 12 月加至 0.75\%；2026 年 4 月因中东能源冲击暂停；2026 年 6 月重启加息至 1.0\%。
横跨二十七个月，负利率时代就此收尾。

\section{债务：353 万亿美元，和它的另一面}

\begin{statrow}{4}
  \statitem{353}{万亿美元}{全球债务总额，占全球 GDP 约 305\%（IIF，2026.03）}
  \statitem{40}{万亿美元}{美国国债规模首次突破，可流通部分逾 30 万亿（2026.08）}
  \statitem{2.3}{万亿美元}{全球私募信贷管理规模，2020 年以来翻倍（2025）}
  \statitem{3000}{亿美元}{全球稳定币市值量级（2026.08）}
\end{statrow}

\begin{databox}{必须同时记住的一句}
353 万亿美元的债务，同时也是 353 万亿美元的资产。讨论「债务太多」之前，先问一句：这笔账是谁欠谁的。
这就是第 5 章那副眼镜的用法。
\end{databox}

\section{储备：黄金第一次超过美国国债}

欧洲央行 2026 年 6 月发布的《欧元的国际角色》报告显示，截至 2025 年底，
全球央行储备资产的构成发生了一次结构性变化。

\begin{barfig}
  \barrow{黄金}{27}{27\%}
  \barrow{美国国债}{22}{22\%}
  \barrow{其他美元资产}{20}{20\%}
  \barrow{欧元}{15}{15\%}
\end{barfig}

黄金占比从 2024 年底的 20\% 升至 27\%，首次超过美国国债的 22\%；
美元计价资产合计仍占 42\%，仍是最大的单一币种。
全球央行黄金储备已超过 3.6 万吨，接近布雷顿森林体系鼎盛时期的约 3.8 万吨水平；
2025 年央行净购金 850 吨；金价在 2026 年 1 月一度接近每盎司 5600 美元。

一个容易被忽略的细节：稳定币发行商 Tether 在 2025 年成为单一最大黄金买家，购入超过 100 吨。

\section{市场与中国经济}

\begin{tablewrap}
\begin{xltabular}{\linewidth}{>{\raggedright\arraybackslash}p{0.26\linewidth}Xr}
\toprule
\bthead{项目} & \bthead{数值} & \bthead{时点} \\
\midrule
\textbf{上证指数} & 3986.30 点，当月上涨 4.02\%，两市成交约 2.13 万亿元 & 2026.08.31 \\
\textbf{深证成指} & 14015 点，当月上涨 3.21\% & 2026.08.31 \\
\textbf{创业板指} & 3438.68 点，当月上涨 2.83\% & 2026.08.31 \\
\textbf{日经 225} & 首次突破 70000 点 & 2026.06 \\
\textbf{中国上半年 GDP} & 同比增长 4.7\% & 2026.H1 \\
\textbf{中国公募基金规模} & 约 38 万亿元 & 2026 年初 \\
\textbf{被动投资} & 全球规模首次超越主动管理 & 2025--2026 \\
\textbf{数据中心资本开支} & 同比增长约 57\%，进入超级周期 & 2025 \\
\bottomrule
\end{xltabular}
\end{tablewrap}

\section{把这一页留到 2028 年}

\begin{takeaway}{使用说明}
两年后再翻开这一页，请按这个顺序看：先看哪些数字变了，再问一句为什么变；
然后看哪些结构没变。大概率的结果是，数字面目全非，而机制一条没换。
如果真的有某条机制变了，那才是值得停下来认真研究的事。
\end{takeaway}

\cleardoublepage
\chapter{延伸阅读}
\markboth{延伸阅读}{延伸阅读}

\begin{ledetext}
这本书只画了地图，画不出地形。下面五本，每一本都比这本书更深，也更难。
不必按顺序读，挑一本你此刻最想弄明白的那件事。
\end{ledetext}

\begin{flowlist}
\flowitem{01}{《现代经济中的货币创造》，英格兰银行季刊，2014}{
本书第 7 章的直接依据。官方出版物，免费，二十页，把「贷款创造存款」这件事讲得比任何教科书都清楚。
如果你只读一份延伸材料，读这个。它会让你确信第 7 章不是奇谈怪论，而是央行自己写的常识。}
\flowitem{02}{《疯狂、恐慌与崩溃》，查尔斯·金德尔伯格}{
本书第 17 章那个五幕剧本的来源。它把四百年的泡沫史压缩成一个可复用的结构，
读完你会发现每一次新的狂热都在按同一个剧本演，只是换了道具。}
\flowitem{03}{《债务危机》，瑞·达利欧}{
第 18 章债务长周期的框架出处。这本书的价值在框架不在结论：
去杠杆的四种手段如何配比，决定了一个国家是走出 1930 年代还是走进日本的三十年。}
\flowitem{04}{《大衰退》，辜朝明}{
理解「资产负债表衰退」最好的入口，也是第 9 章利率传导为什么会失灵的完整版。
读完你会明白，为什么有时候把利率降到零，钱也不动。}
\flowitem{05}{《非对称风险》，纳西姆·塔勒布}{
第 22 章遍历性那一节的思想来源。它讲的其实只有一件事：
集合平均不等于时间平均，而你只有一条时间线。这句话比任何收益率都重要。}
\end{flowlist}

\bookbreak

\section{还想再往下走}

如果读完上面五本还想继续，往三个方向走。

\textbf{往制度走。}去读会计准则、破产法和清算制度。它们无聊到极点，
但第 25 章说过，信任的物理基础全在那里。理解了它们，你会对「一个国家的利率下限由它的法院决定」这句话有实感。

\textbf{往数据走。}把中国人民银行、美联储、欧洲央行和 IIF 的定期报告加进你的阅读列表。
它们是免费的一手材料，而绝大多数财经新闻只是对它们的二手转述，还经常转错。
自己看一遍原文，你会发现新闻标题和报告结论之间的距离有多大。

\textbf{往怀疑走。}找那些和本书结论相反的严肃论述来读。
比如关于货币数量论的辩护、关于被动投资并未损害价格发现的实证、关于私募信贷风险被高估的分析。
一个只读支持自己观点材料的人，读得越多越危险。

\cleardoublepage


\cleardoublepage
\printbibliography[heading=bibintoc, title={参考文献与数据来源}]

\cleardoublepage
\printindex

\end{document}